\documentclass[11pt,a4paper]{article}
\usepackage[margin=26mm]{geometry}
\usepackage[T1]{fontenc}
\usepackage[utf8]{inputenc}
\usepackage{lmodern,microtype}
\usepackage{amsmath,amssymb,amsthm,mathtools}
\usepackage{booktabs,array,longtable}
\usepackage{enumitem}
\usepackage{needspace}
\usepackage{xcolor,xurl}
\usepackage[colorlinks=true,linkcolor=blue!50!black,citecolor=green!35!black,urlcolor=blue!50!black]{hyperref}
\hypersetup{pdftitle={Real asymptotic expansions and inverse functions},
  pdfauthor={Vladimir Reshetnikov, with Claude},
  pdfsubject={Power-logarithmic reversion with real exponents}}
\setlength{\parskip}{0.3em}
\setlength{\emergencystretch}{2.5em}
\allowdisplaybreaks
\numberwithin{equation}{section}

\theoremstyle{plain}
\newtheorem{theorem}{Theorem}[section]
\newtheorem{proposition}[theorem]{Proposition}
\newtheorem{lemma}[theorem]{Lemma}
\newtheorem{corollary}[theorem]{Corollary}
\theoremstyle{definition}
\newtheorem{definition}[theorem]{Definition}
\newtheorem{example}[theorem]{Example}
\newtheorem{algorithm}[theorem]{Algorithm}
\theoremstyle{remark}
\newtheorem{remark}[theorem]{Remark}

\newcommand{\R}{\mathbb R}
\newcommand{\C}{\mathbb C}
\newcommand{\N}{\mathbb N}
\newcommand{\Z}{\mathbb Z}
\newcommand{\Q}{\mathbb Q}
\newcommand{\OO}{\mathcal O}
\newcommand{\oo}{o}
\newcommand{\dd}{\mathrm d}
\newcommand{\D}{\mathcal D}
\newcommand{\val}{\operatorname{val}}
\newcommand{\Res}{\operatorname{Res}}
\newcommand{\supp}{\operatorname{supp}}
\newcommand{\abs}[1]{\left|#1\right|}
\newcommand{\M}{\mathrm M}
\newcommand{\bk}{\boldsymbol k}
\newcommand{\bnu}{\boldsymbol\nu}
\newcommand{\bdelta}{\boldsymbol\delta}
\newcommand{\PL}{\mathcal P}
\newcommand{\W}{\mathrm W}
\DeclareMathOperator{\Cat}{Cat}

\title{Real asymptotic expansions and inverse functions\\[0.4em]
\large Power--logarithmic series, real branches, and error transport}
\author{Vladimir Reshetnikov\thanks{The initial manuscript was prepared with Claude (Anthropic).}}
\date{September 2026}

\begin{document}
\maketitle

\begin{abstract}
Real powers and logarithms lead naturally beyond Taylor and Puiseux series.
For example, the positive inverse of $x+x^2(1+\log x)$ begins
$y-y^2(1+\log y)+y^3(2\log^2y+5\log y+3)$, whereas the inverse of
$x+x^{\sqrt2}$ begins $y-y^{\sqrt2}+\sqrt2\,y^{2\sqrt2-1}$.
We develop a calculus of complete power--logarithmic blocks with real
exponents, valid at finite endpoints and at infinity.

Marker Lagrange--B\"urmann inversion and an Euler-operator formula give the
coefficients of arbitrary real powers of an inverse. Real branch selection,
weighted truncation, and error transport are treated explicitly. Convergence
is proved for the analytic finite-parameter class; finite smoothness yields
separate asymptotic theorems. The calculus extends to logarithmic coordinates,
exact inverse cores, finite flat sectors, and Fourier-valued coefficients.
Special-function tails and exponential normalization illustrate the
distinction between convergent local expansions and Poincar\'e expansions.
Logarithmic normalization of Gamma products supplies an exact growth factor
and a relative correction series. Rational interval bounds convert a residual
estimate into a rigorous real root enclosure. Each extension states its
own hypotheses, ordering, and remainder meaning.
\end{abstract}

\tableofcontents

\section{Introduction}
\label{sec:intro}

\subsection{Two questions}

The first question \cite{q1} asks for an asymptotic expansion, as $y\to0^+$, of the inverse of
\begin{equation}
 f_1(x)=x+x^2(1+\log x),\qquad x>0,
 \label{eq:f1}
\end{equation}
on the branch that maps $(0,\infty)$ to itself. Choosing this branch is essential: near a nonzero complex preimage of $0$, the same expression has an analytic inverse germ of the form
\begin{equation}
 c+\frac{y}{c-1}-\frac{3+2\W_{-1}(-e)}{2(c-1)^3}\,y^2+\OO(y^3),\qquad c=e^{\W_{-1}(-e)-1}\approx-0.1178-0.5332\,i,
 \label{eq:complex-germ}
\end{equation}
Here $c$ is a nonzero complex root of $1+x(1+\log x)=0$, where $f_1$ is analytic with a nonzero derivative (Section~\ref{sec:branch}). The real branch has no Taylor expansion at all. Its expansion is
\begin{equation}
\begin{aligned}
 g_1(y)={}&y-y^2(1+L)+y^3\bigl(2L^2+5L+3\bigr)-y^4\Bigl(5L^3+\tfrac{41}{2}L^2+27L+\tfrac{23}{2}\Bigr)\\
 &+y^5\Bigl(14L^4+\tfrac{241}{3}L^3+\tfrac{335}{2}L^2+151L+\tfrac{299}{6}\Bigr)+\OO\bigl(y^6(1+|L|)^5\bigr),
\end{aligned}
 \label{eq:answer1}
\end{equation}
with $L=\log y$: every power of $y$ carries a polynomial in $\log y$, of degree growing with the power, so the scale of the expansion is neither a Taylor nor a Puiseux scale. The question also asks, as a second example, for
\begin{equation}
 g_2(y)=y-y^{\sqrt2}+\sqrt2\,y^{2\sqrt2-1}-\frac{6-\sqrt2}{2}\,y^{3\sqrt2-2}+\Bigl(\frac{17\sqrt2}{3}-4\Bigr)y^{4\sqrt2-3}+\OO\bigl(y^{5\sqrt2-4}\bigr),
 \label{eq:answer2}
\end{equation}
the inverse of $f_2(x)=x+x^{\sqrt2}$: no logarithms this time, but exponents in the semigroup $1+(\sqrt2-1)\N$, which no rational grid contains.

The second question \cite{q2} concerns the forward asymptotics of $(1+x+x^{\sqrt2})^{\sqrt2}$ as $x\to\infty$. Its leading term is $x^2$, but the subsequent terms require ordering two incommensurable positive gaps. The expansion is
\begin{equation}
\begin{aligned}
 (1+x+x^{\sqrt2})^{\sqrt2}={}&x^2+\sqrt2\,x^{3-\sqrt2}+\Bigl(1-\tfrac1{\sqrt2}\Bigr)x^{4-2\sqrt2}+\Bigl(\tfrac{2\sqrt2}{3}-1\Bigr)x^{5-3\sqrt2}\\
 &+\sqrt2\,x^{2-\sqrt2}+\tfrac{13-9\sqrt2}{12}\,x^{6-4\sqrt2}+(2-\sqrt2)\,x^{3-2\sqrt2}+\OO\bigl(x^{7-5\sqrt2}\bigr),
\end{aligned}
 \label{eq:answer3}
\end{equation}
and its inverse, as $y\to\infty$, begins $\sqrt y-\tfrac1{\sqrt2}y^{(2-\sqrt2)/2}+\tfrac{3-\sqrt2}{4}y^{(3-2\sqrt2)/2}+\OO(y^{(4-3\sqrt2)/2})$.

\subsection{What is common to all of this}

All these expansions are finite or infinite sums of \emph{power--log blocks} $w^\beta P(\log w)$, where $w\to0^+$ is a local variable ($w=y$, $w=1/x$, \dots), $\beta$ is a real exponent, and $P$ is a polynomial. Their supports lie in translates of semigroups generated by finitely many positive gaps. These semigroups are locally finite, even when the groups generated by the same gaps are dense in $\R$, so every truncation is a finite computation. The logarithmic polynomials are handled by the Euler identity $\frac{\dd}{\dd w}\bigl[w^\beta P(\log w)\bigr]=w^{\beta-1}(\beta+\tfrac{\dd}{\dd L})P(L)$; and remainders must keep their logarithmic factor, because $y^3(2\log^2y+5\log y+3)$ is not $\OO(y^3)$.

This article develops the algebra and analysis needed to construct and justify these expansions. Its starting points are the questions \cite{q1,q2} and the phase-coordinate approach of \cite{proveit}. The treatment gives coefficient formulas, branch theorems, convergence proofs and remainder estimates, then extends the construction to finite endpoints from either side, infinity, observables of an inverse, and scales involving iterated logarithms or exponentially small terms.

\subsection{Summary of results}

\begin{itemize}[leftmargin=1.6em]
\item \textbf{Scale} (Section~\ref{sec:scale}). Power--log blocks are totally ordered by size, the finitely generated exponent semigroups are locally finite, and the truncated arithmetic (products, real powers, logarithms, exponentials, composition with polynomials in the logarithm) is exact and terminates. The single primitive is the composition recurrence $Q_{k+1}=\bigl((a-k)Q_k+Q_k'\bigr)/(k+1)$ for $(1+t)^aP(L+\log(1+t))$.
\item \textbf{Forward expansions} (Section~\ref{sec:forward}). Expressions built from constants, the variable, arithmetic, real powers, logarithms, bounded exponentials and analytic functions of expressions with a limit have convergent power--log expansions; structural error propagation supplies a rigorous remainder class, which answers the second question.
\item \textbf{Branch} (Section~\ref{sec:branch}). For $f(x)=ax^p(1+\sum_ix^{\delta_i}B_i(\log x))$ with $a,p\ne0$, $\delta_i>0$, the real branch is the one with $x>0$ and $x/(y/a)^{1/p}\to1$; it is unique, and the complex germ \eqref{eq:complex-germ} is explained.
\item \textbf{Lagrange--B\"urmann with a marker} (Section~\ref{sec:lagrange}). A residue proof of $\Phi(u)=\Phi(q)+\sum_N\frac{(-\eta)^N}{N!}\partial_q^{N-1}[\Phi'(q)h(q)^N]$ at a positive point, valid analytically and formally, with the transport of observables and the general-core form.
\item \textbf{Coefficients} (Section~\ref{sec:coefficients}). For every multi-index $\bk$ of weight $\sigma$, the block of $g(y)^r$ at $z^{r+\sigma}$, $z=(y/a)^{1/p}$, is
 $\frac{(-1)^{|\bk|}r}{p^{|\bk|}\bk!}\prod_{j=1}^{|\bk|-1}(\D+r+\sigma+pj)\prod_iB_i(L)^{k_i}$, $\D=\dd/\dd L$; equal weights are added. Degrees, leading logarithms (Catalan numbers), closed forms for constant coefficients (Fuss--Catalan), $\log g$, and the formal uniqueness through the triangular equation, Picard iteration and Newton iteration.
\item \textbf{Convergence} (Section~\ref{sec:convergence}). The expansion converges absolutely for small positive $y$ after substituting finitely many small parameters $z^{\delta_i}(\log z)^j$, with coefficient bounds $|C_{\bk}(L)|\le MK^{|\bk|}(1+|L|)^{D(\bk)}$; an explicit conditional tail majorant follows from Rouch\'e's theorem.
\item \textbf{Remainders} (Section~\ref{sec:remainders}). Truncation below any exponent cutoff leaves $z^{1+\sigma_*}C_{\sigma_*}(L)+\oo(z^{1+\sigma_*})$ where $\sigma_*$ is the first omitted weight and $C_{\sigma_*}$ the complete coefficient there, computed from the one-step boundary of the retained multi-index set; depth truncation has its own bound.
\item \textbf{Transport} (Section~\ref{sec:transport}). Forward remainders and residuals are transported to the inverse with the correct scaling $z^{1-p}$; a residual bracket gives existence and a two-sided error bound; and a theorem with a complete proof shows that finite smoothness of the perturbation, with derivative bounds in the power--log scale, suffices for the marker formula with an explicit remainder, which covers perturbations such as $x^2\sin(\log x)$.
\item \textbf{Endpoints} (Section~\ref{sec:endpoints}). Finite points from either side and $\pm\infty$ reduce to the same local inversion problem through a local variable and a signed leading power $p\ne0$; the limit $y_0$ may be finite or infinite; powers $(x-x_0)^r$ and $x^r$ of the inverse are obtained directly. Worked examples include the Wright omega function and the inverse of \eqref{eq:answer3}.
\item \textbf{Lambert coordinates and boundaries} (Section~\ref{sec:boundaries}). Leading logarithmic and exponential cores reduce to a convergent Lambert expansion in a reciprocal-logarithm coordinate, with a polynomial in its logarithm at every order. This covers $x\log x$ near $0$ and $xe^x$ at infinity. Zero gaps, flat perturbations and oscillatory coefficients require further scale choices; symbolic exponents can use depth truncation when their branch assumptions are proved.
\item \textbf{Coordinates and extended scales}. Exact changes of source and target coordinate transport inverse expansions and their errors. Weighted-support algorithms organize coefficient calculation; interval residual bounds turn local approximations into enclosures of unique roots. Further sections develop exact-core perturbations, iterated logarithmic scales, Fourier coefficients and exponentially small sectors.
\end{itemize}

\subsection{Conventions}

$\N=\{0,1,2,\dots\}$. For a local variable $w\to0^+$ we write $L=\log w$ and $\M(w)=1+|\log w|$. $\OO$ and $\oo$ refer to $w\to0^+$ with all other data fixed; no uniformity in parameters is asserted unless stated. $\D=\dd/\dd L$ acts on polynomials in $L$. For a multi-index $\bk=(k_1,\dots,k_m)\in\N^m$ we write $|\bk|=\sum k_i$, $\bk!=\prod k_i!$, and $\bk\cdot\bdelta=\sum k_i\delta_i$. Real powers of positive quantities are $w^\beta=e^{\beta\log w}$. Real observables use the selected real branches; auxiliary complex powers may occur in analytic proofs or Fourier coefficient representations. Power and weight cutoffs are \emph{exclusive}: a cutoff $q$ retains the blocks of exponent $<q$. An exponential sector depth $N$ is inclusive and retains degrees through $N$.

\section{The power--log scale and its exact arithmetic}
\label{sec:scale}

\subsection{Blocks and their order}

Throughout, $w$ denotes a positive real variable tending to $0^+$, $L=\log w\to-\infty$ and $\M(w)=1+|\log w|$.

\begin{definition}[Blocks]
A \emph{power--log block} is an expression $w^\beta P(L)$ with $\beta\in\R$ and $P\in\R[L]$, $P\ne0$. Its \emph{exponent} is $\beta$ and its \emph{logarithmic degree} is $\deg P$.
\end{definition}

\begin{lemma}[Power beats logarithm]
\label{lem:dominance}
For every $\varepsilon>0$ and every real $b$, $w^\varepsilon\M(w)^b\to0$ as $w\to0^+$. Consequently, if $\beta<\gamma$ and $P,Q\ne0$, then $w^\gamma Q(L)=\oo\bigl(w^\beta P(L)\bigr)$; and within a fixed exponent, a higher logarithmic degree dominates: $w^\beta L^j=\oo(w^\beta L^k)$ for $j<k$.
\end{lemma}
\begin{proof}
Put $w=e^{-t}$, $t\to+\infty$: $w^\varepsilon\M(w)^b=e^{-\varepsilon t}(1+t)^b\to0$. For the second claim the ratio is $w^{\gamma-\beta}Q(L)/P(L)$, and $|Q(L)/P(L)|\le C\M(w)^{\deg Q}$ for $w$ small because $P(L)$ has no zeros for $L$ sufficiently negative and $|P(L)|\ge c|L|^{\deg P}\ge c$ there; the first claim finishes. The third claim is clear.
\end{proof}

Blocks are therefore totally ordered by size: by increasing exponent, and at equal exponents by decreasing logarithmic degree. A finite sum of blocks with distinct exponents is asymptotic to its block of least exponent (whole polynomial included), and vanishes identically only if every block does. This is why truncations are taken by \emph{complete blocks}: cutting a block inside its polynomial would drop terms larger than the retained terms of the next exponent.

\begin{definition}[Remainder classes]
For $\beta\in\R$ and an integer $k\ge0$ we write $R(w)=\OO_{\beta,k}$, or in the package \wl{PowerLogRemainder[w, $\beta$, k]}, for $R(w)=\OO\bigl(w^\beta\M(w)^k\bigr)$ as $w\to0^+$.
\end{definition}

By Lemma~\ref{lem:dominance}, $\OO_{\beta,k}\subseteq\OO_{\beta',k'}$ whenever $\beta>\beta'$, or $\beta=\beta'$ and $k\le k'$; the class with the smaller exponent, and at equal exponents the larger degree, is the coarser one. A block $w^\beta P(L)$ with $\deg P=k$ belongs to $\OO_{\beta,k}$ and to no $\OO_{\beta,k'}$ with $k'<k$. Sums and products of remainder classes obey
\begin{equation}
 \OO_{\beta,k}+\OO_{\beta',k'}=\OO_{\min(\beta,\beta'),\,k''},\qquad
 \OO_{\beta,k}\cdot\OO_{\beta',k'}=\OO_{\beta+\beta',\,k+k'},
 \label{eq:O-arithmetic}
\end{equation}
where $k''$ is the degree attached to the smaller exponent, and the larger of the two when the exponents coincide.

\subsection{Exponent semigroups}

\begin{lemma}[Local finiteness]
\label{lem:finite}
Let $\delta_1,\dots,\delta_m>0$ and $S=\{\bk\cdot\bdelta:\bk\in\N^m\}$. For every real $H$ the set $\{\bk\in\N^m:\bk\cdot\bdelta\le H\}$ is finite, hence $S\cap(-\infty,H]$ is finite, and every $\sigma\in S$ has only finitely many representations $\sigma=\bk\cdot\bdelta$.
\end{lemma}
\begin{proof}
If $\bk\cdot\bdelta\le H$ then $k_i\le H/\delta_i$ for each $i$.
\end{proof}

When the gaps are incommensurable the \emph{group} generated by them is dense in $\R$ (for instance $\Z+\sqrt2\Z$), but the semigroup of nonnegative combinations has only finitely many elements below any bound. This is the whole reason why expansions with irrational exponents can be computed term by term: the support of every expansion in this article lies in a translate $\beta_0+S$ of such a semigroup. Following \cite{proveit}, a subset of $\R$ that is finite below every bound is called \emph{left-finite}; the semigroups $S$ are left-finite, and so are their translates and finite unions.

\subsection{Jets and their arithmetic}

\begin{definition}[Jets]
Fix gaps as above and $\beta_0\in\R$. A \emph{jet} (of cutoff $H$) is a finite sum $U=\sum_{\sigma\in T}w^\sigma C_\sigma(L)$ with $T\subset\beta_0+S$ finite, $C_\sigma\in\R[L]$, all $\sigma<H$; equal exponents are merged and zero polynomials are discarded. Its \emph{valuation} $\val U$ is its least exponent ($+\infty$ for the empty jet). The jet algebra $\PL_{<H}$ consists of jets truncated below $H$, with the sum and the product computed exactly and then truncated.
\end{definition}

The jet algebra is the quotient of the ring of finite sums of blocks by the ideal of blocks of exponent $\ge H$. In the package a jet is a sorted list of pairs $\{\sigma,C_\sigma\}$ with $\sigma$ an exact real number in canonical form (\wl{RootReduce} for algebraic numbers) and $C_\sigma$ a polynomial in a private symbol $\ell$; two exponents are equal if and only if their canonical forms coincide, and they are ordered exactly (Section~\ref{sec:package}).

\subsubsection*{Differentiation}

\begin{lemma}[Euler identity]
\label{lem:euler}
For $\beta\in\R$, $P\in\R[L]$ and $r\in\N$,
\begin{equation}
 \frac{\dd}{\dd w}\bigl[w^\beta P(L)\bigr]=w^{\beta-1}(\beta+\D)P(L),\qquad
 \frac{\dd^r}{\dd w^r}\bigl[w^\beta P(L)\bigr]=w^{\beta-r}\prod_{j=0}^{r-1}(\beta-j+\D)P(L).
 \label{eq:euler}
\end{equation}
The operators $\beta-j+\D$ commute, and $(\beta+\D)$ preserves $\deg P$ when $\beta\ne0$ and lowers it by one when $\beta=0$.
\end{lemma}
\begin{proof}
$\frac{\dd}{\dd w}[w^\beta P(\log w)]=\beta w^{\beta-1}P+w^\beta P'(\log w)/w$. Iterate.
\end{proof}

The identity says that differentiation acts on a block as the scalar $\beta$ plus $\D$ on the polynomial; the exponent bookkeeping and the polynomial arithmetic separate completely. In the coefficient formulas below the operators $\D+c$ with $c>0$ appear, and their action on the coefficient vector $(z_0,\dots,z_d)$ of $P=\sum z_sL^s$ is bidiagonal: $(c+\D)P=\sum_s\bigl(cz_s+(s+1)z_{s+1}\bigr)L^s$.

\subsubsection*{Unit series}

Let $U$ be a jet with $\val U>0$. Then $U^k$ has valuation $\ge k\val U$, so in $\PL_{<H}$ only finitely many powers of $U$ are nonzero, and every power series $\Phi(t)=\sum_k\phi_kt^k$ can be substituted:
\begin{equation}
 \Phi(U)=\sum_{k<H/\val U}\phi_kU^k\in\PL_{<H}.
 \label{eq:unit-series}
\end{equation}
In particular $(1+U)^a=\sum_k\binom akU^k$ for every real $a$ (with $\binom ak=a(a-1)\cdots(a-k+1)/k!$), $\log(1+U)=\sum_{k\ge1}(-1)^{k+1}U^k/k$, $e^U=\sum_kU^k/k!$, and $\phi(c+U)=\sum_k\phi^{(k)}(c)U^k/k!$ for any $\phi$ analytic at $c$. These are identities between finite sums of blocks modulo exponents $\ge H$; their analytic content is discussed in Section~\ref{sec:convergence}.

\subsubsection*{Composition with the logarithm}

The one operation that mixes the exponent bookkeeping with the polynomial coefficients is the substitution of a jet into the logarithm: if $w$ is replaced by $w(1+U)$ then $L$ becomes $L+\log(1+U)$, and a block $w^a P(L)$ becomes $w^a(1+U)^aP(L+\log(1+U))$. The following recurrence computes this without expanding $\log(1+U)$ separately.

\begin{lemma}[Composition recurrence]
\label{lem:composition}
Let $P\in\R[L]$ and $a\in\R$. For $|t|<1$,
\begin{equation}
 (1+t)^aP\bigl(L+\log(1+t)\bigr)=\sum_{k\ge0}Q_k(L)\,t^k,\qquad
 Q_0=P,\qquad Q_{k+1}=\frac{(a-k)Q_k+Q_k'}{k+1},
 \label{eq:composition}
\end{equation}
with $Q_k\in\R[L]$, $\deg Q_k\le\deg P$, and, when $a\in\N$ and $P$ is constant, $Q_k=0$ for $k>a$. Consequently, for a jet $U$ of positive valuation, $(1+U)^aP(L+\log(1+U))=\sum_kQ_k(L)U^k$ in $\PL_{<H}$.
\end{lemma}
\begin{proof}
Put $F(t)=(1+t)^aP(L+\log(1+t))$; it is analytic in $|t|<1$ with values in $\R[L]$, so $F=\sum Q_kt^k$ with $Q_k\in\R[L]$ and $Q_0=P$. Differentiating, $(1+t)\partial_tF=aF+\partial_LF$, because $\partial_t$ acting on the two factors gives $a(1+t)^{a-1}P+(1+t)^{a-1}P'$. Comparing coefficients of $t^k$ in $(1+t)\sum(k+1)Q_{k+1}t^k=\sum(aQ_k+Q_k')t^k$ gives $(k+1)Q_{k+1}+kQ_k=aQ_k+Q_k'$. The degree does not increase, and for $a\in\N$ and constant $P$ the factor $(a-k)$ annihilates $Q_{a+1}$ and hence all later $Q_k$ (for nonconstant $P$ the derivative keeps the sequence alive: $\log(1+t)$ is not a polynomial). The jet statement follows because substitution of $U$ for $t$ is a ring homomorphism on finite sums.
\end{proof}

\begin{remark}
The recurrence is the whole composition calculus needed here: the forward model $ax^p\bigl(1+\sum x^{\delta_i}B_i(\log x)\bigr)$ evaluated at $x=z(1+U)$ is
$az^p(1+U)^p\bigl(1+\sum_iz^{\delta_i}(1+U)^{\delta_i}B_i(L+\log(1+U))\bigr)$, each factor of which is an instance of \eqref{eq:composition} (with $P=1$ for the first). Its derivative is obtained from the same lemma with $a$ replaced by $a-1$ and $P$ by $aP+P'$, since $\frac{\dd}{\dd x}[x^aP(\log x)]=x^{a-1}(aP+P')(\log x)$.
\end{remark}

\subsection{Blocks with nonconstant leading coefficient}

Everything above allows the polynomial in the leading block to be nonconstant. Two operations do not: taking a real non-integer power or a logarithm of a sum whose leading block is $w^\beta P(L)$ with $\deg P\ge1$ produces $P(L)^a$ or $\log P(L)$, which are not polynomials in $L$. Such expressions live in larger scales (rational functions of $L$, or iterated logarithms, Section~\ref{sec:boundaries}) and are outside the class treated here; the package reports them as unsupported. Nonnegative integer powers of any jet are computed by repeated multiplication and stay in the class.

\section{Forward expansions}
\label{sec:forward}

Before inverting anything we need the forward function in the normal form $ax^p(1+\sum_ix^{\delta_i}B_i(\log x))$, and the second question \cite{q2} is a forward problem in its own right. This section shows that a large class of expressions has convergent power--log expansions, and describes the precision-tracking evaluator that computes them with a rigorous remainder class.

\subsection{Convergent power--log expansions}

\begin{definition}
\label{def:pl-expansion}
A function $F$ on $(0,w_0)$ has a \emph{convergent power--log expansion} with support in $\beta_0+S$ ($S$ a finitely generated positive semigroup with gaps $\delta_1,\dots,\delta_m$) if there are polynomials $C_\sigma\in\R[L]$, constants $M,K\ge1$, and integers $d_0,d\ge0$ such that
\begin{equation}
 F(w)=\sum_{\sigma\in S}w^{\beta_0+\sigma}C_\sigma(L)\quad\text{for }0<w<w_1,
 \qquad \|C_\sigma\|_1\le M K^{n(\sigma)},\quad
 \deg C_\sigma\le d_0+d\,n(\sigma),
 \label{eq:pl-expansion}
\end{equation}
where $n(\sigma)$ is the least $|\bk|$ with $\bk\cdot\bdelta=\sigma$, $\|P\|_1$ is the sum of the absolute values of the coefficients of $P$, and the series converges absolutely. Thus $|C_\sigma(L)|\le M K^{n(\sigma)}\M(w)^{d_0+dn(\sigma)}$. Zero coefficients satisfy the degree bound by convention. We write $F\in\mathcal E$. The fixed allowance $d_0$ is essential: it admits a polynomial in $\log w$ at the leading exponent, including $\log w$ itself.
\end{definition}

The coefficient bound makes the tail after any cutoff controllable (Lemma~\ref{lem:tail} below), and it is stable under all the operations we need.

\begin{lemma}[Tail of a convergent expansion]
\label{lem:tail}
Let $F\in\mathcal E$ as in \eqref{eq:pl-expansion}, $\delta_*=\min_i\delta_i$, and $H>0$. Then
\begin{equation}
 F(w)-\sum_{\sigma<H}w^{\beta_0+\sigma}C_\sigma(L)=\OO\bigl(w^{\beta_0+H}\M(w)^{d_0+Nd}\bigr),\qquad N=\lceil H/\delta_*\rceil.
 \label{eq:tail}
\end{equation}
More precisely, the omitted part equals $w^{\beta_0+\sigma_*}C_{\sigma_*}(L)+\oo(w^{\beta_0+\sigma_*})$ where $\sigma_*=\min\{\sigma\in S:\sigma\ge H\}$.
\end{lemma}
\begin{proof}
Split the omitted terms into those with $n(\sigma)\le N-1$ and the rest. The first group is finite (Lemma~\ref{lem:finite}) and each term is $\OO(w^{\beta_0+H}\M^{d_0+(N-1)d})$. For the rest, the number of weights whose shortest representation has length $n$ is at most $m^n$, and $\sigma\ge n\delta_*$, so their sum is bounded by
\[
 Mw^{\beta_0}\M^{d_0}\sum_{n\ge N}\bigl(mKw^{\delta_*}\M^d\bigr)^n
 =\OO(w^{\beta_0+N\delta_*}\M^{d_0+Nd}).
\]
Since $N\delta_*\ge H$, this proves the bound. For the refined statement let $\sigma_{**}$ be the next element of $S$ after $\sigma_*$, which exists by local finiteness when $S$ has a generator, and choose $\sigma_*<H'<\sigma_{**}$. Applying the bound at $H'$ shows that everything beyond $\sigma_*$ is $\OO(w^{\beta_0+H'}\M^{d_0+N'd})=\oo(w^{\beta_0+\sigma_*})$. The difference $\sigma_{**}-\sigma_*$ need not be at least $\delta_*$. Finite support with no further terms has zero tail.
\end{proof}

\begin{theorem}[Closure]
\label{thm:closure}
The class $\mathcal E$ contains the constants and $w$ and is closed under sums, products, and the following operations, with the supports and gaps indicated.
\begin{enumerate}[leftmargin=1.8em]
\item If $F\in\mathcal E$ has leading block $cw^\beta$ with a \emph{constant} $c>0$, then $F^a\in\mathcal E$ for every real $a$ and $\log F\in\mathcal E$; their supports are contained in $a\beta+S'$ and $\{0\}\cup S'$ respectively, with block $\log c+\beta L$ at exponent $0$ in the logarithm. Here $S'$ is a finitely generated positive semigroup containing the exponents of $F/(cw^\beta)-1$; the proof establishes that such a semigroup exists.
\item If $F\in\mathcal E$ has no negative-exponent blocks and its exponent-$0$ block is $c+kL$ with $k\in\R$, then $e^F\in\mathcal E$ with leading block $e^cw^k$. The argument $F$ may grow logarithmically.
\item If $F\in\mathcal E$ has $\beta_0\ge0$ and constant exponent-$0$ block $c$, and $\phi$ is analytic at $c$, then $\phi\circ F\in\mathcal E$.
\end{enumerate}
\end{theorem}
\begin{proof}
First note that a positive-valuation expansion $U$ can be represented with support in a positive semigroup based at $0$: if necessary, add its positive leading exponent as a generator and enlarge the constants in \eqref{eq:pl-expansion}. Choose one shortest multi-index $\bk$ for each resulting weight $\sigma$, put $D=d_0+d$, and introduce
\[
 T_{ij}=w^{\delta_i}L^j,\qquad 1\le i\le m,\quad 0\le j\le D.
\]
For $n=|\bk|\ge1$, every $w^\sigma L^s$ with $s\le d_0+dn\le Dn$ is a product of $n$ such parameters: distribute the $s$ logarithmic factors among the $n$ generator factors, each receiving at most $D$. Fix one such distribution for every monomial. The coefficient-norm bound makes the resulting series analytic in a sufficiently small polydisc, since its terms of degree $n$ have total absolute coefficient at most $M(mK)^n$. Analytic composition with $(1+U)^a$, $\log(1+U)$, $e^U$, or $\phi(c+U)$ is therefore legitimate. Cauchy estimates give exponential coefficient bounds in the new parameters, and their logarithmic degree grows at most linearly in parameter degree.

Regrouping by weight preserves bounds of the form \eqref{eq:pl-expansion}: a representation of weight $\sigma$ has length at most $\sigma/\delta_*$, whereas its shortest length is at least $\sigma/\delta_{\max}$; their ratio is bounded, and the number of parameter monomials of a given total degree is at most exponential. Enlarge $M,K,d$ to absorb both factors.

For sums and products, use the union of the generator sets and, for a sum of differently shifted supports, the positive difference between their shifts as an additional generator. Coefficient norms are subadditive and submultiplicative, so the same bounds apply. If leading blocks cancel, rebasing at the next nonzero exponent still admits finitely many positive generators: the multi-indices above that exponent form an upward-closed subset of $\N^m$, with finitely many coordinatewise minimal elements. Subtracting the new leading exponent from the weights of those elements supplies finitely many nonnegative shifts, together with the original generators. This also justifies normalizing $F$ by its actual leading monomial in (1). Separate the finite leading polynomial before applying the small-parameter construction; in (2) separate $c+kL$ to obtain the factor $e^cw^k$.
\end{proof}

The theorem supplies sufficient closure conditions, and every truncation is a genuine asymptotic expansion by Lemma~\ref{lem:tail}. They are not necessary conditions for every individual expression: for example $\sqrt{L^2}=-L$ for $L<0$, and polynomials may be applied to an unbounded logarithm. In general, non-integer powers of a nonconstant leading polynomial, its logarithm, exponentials of negative-power arguments, and oscillatory functions of an unbounded logarithm require larger scales (Section~\ref{sec:boundaries}).

\begin{corollary}[Laurent and Puiseux composition]
If $U\in\mathcal E$ has a positive constant leading coefficient and positive valuation, and a selected real branch of $\phi(c+v)$ has a convergent Puiseux expansion $v^{n_0/d}A(v^{1/d})$ for $v\to0^+$, with $d\ge1$ an integer and $A$ analytic at $0$, then $\phi(c+U)\in\mathcal E$. The same holds from below using $\phi(c-v)$.
\end{corollary}
\begin{proof}
Apply the closure theorem to $U^{1/d}$, the analytic function $A$, and $U^{n_0/d}$. Negative $n_0$ includes poles. The side of approach fixes the real branch of a ramified head.
\end{proof}

Variable powers also reduce to the closure theorem through $F^G=\exp(G\log F)$ when $F>0$ and the resulting exponential argument satisfies its hypotheses; this includes $w^w=\exp(w\log w)$. For a nonzero leading block $w^\beta P(L)$, the eventual sign is the sign of $(-1)^{\deg P}[L^{\deg P}]P$, so absolute values reduce to multiplication by that sign whenever it is provable.

\subsection{The precision-tracking evaluator}
\label{sec:evaluator}

The package computes forward expansions by structural recursion over the expression tree, carrying for every subexpression a triple $(T,P,D)$: a jet $T$ and the assertion
\begin{equation}
 \text{subexpression}=T+\OO\bigl(w^P\M(w)^D\bigr),\qquad P\in\R\cup\{+\infty\},
 \label{eq:triple}
\end{equation}
where $P=+\infty$ means that $T$ is exact. A global working order $K_w$ bounds the number of terms of every unit series that is expanded. The rules, all justified by \eqref{eq:O-arithmetic} and Lemma~\ref{lem:tail}, are:
\begin{itemize}[leftmargin=1.6em]
\item constants and $w$ are exact; \wl{Plus}: compute $T_1+T_2$ and truncate below $P=\min(P_1,P_2)$. Combine the input remainder degrees at $P$ with the degree of every explicit block discarded at exponent $P$;
\item \wl{Times}: if $R_i$ denotes the input remainder, the omitted terms include $R_1T_2$, $R_2T_1$ and $R_1R_2$. Their exponent/degree pairs follow from \eqref{eq:O-arithmetic}; use the coarsest pair and truncate $T_1T_2$ below that exponent. Also include the degrees of explicit product blocks discarded exactly at this boundary. A nonempty jet has valuation strictly below its remainder exponent; an empty jet contributes its remainder's exponent and degree when estimating its size;
\item \wl{Power} with an exact real exponent $a$: nonnegative integers by binary exponentiation with precision tracked at each product (exact for exact input). Otherwise the leading block $cw^\beta$ of $T$ must have constant $c>0$ (or $c<0$ with integer $a$); with $U=T/(cw^\beta)-1$, relative precision $P-\beta$ and relative cutoff $c_{\mathrm{rel}}=\min(P-\beta,K_w-a\beta)$, the result is $c^aw^{a\beta}\sum_k\binom akU^k$ with precision $a\beta+c_{\mathrm{rel}}$. For a nonempty $U$, its series-truncation degree is bounded by $\lceil c_{\mathrm{rel}}/\val U\rceil\deg_{\max}U$, combined with the input error degree $D$; an empty $U$ has no computed series tail;
\item \wl{Log}: the leading block must be $cw^\beta$ with constant $c>0$; the result is $\log c+\beta L+\log(1+U)$;
\item \wl{Exp}: $T$ must have no negative exponents and its exponent-$0$ block must be $c+kL$ with exact numeric $k$; the result is $e^cw^k\exp(U)$;
\item an analytic head $\phi$ (\wl{Sin}, \wl{ArcTan}, \wl{BesselJ[0,\,$\cdot$]}, \dots) applied to $T=c+U$ with constant $c$: $\sum_k\phi^{(k)}(c)U^k/k!$ with the Taylor coefficients obtained from \wl{Series[$\phi$[c + t], \{t, 0, N\}]};
\item any other subexpression is expanded by \wl{Series[expr, \{w, 0, n\}, Assumptions -> w > 0]}; the coefficients (which may contain $\log w$) and the irrational-power monomials that \wl{Series} leaves outside its \wl{SeriesData} are parsed as jets, the precision is the order of the \wl{O} term, and the logarithmic degree $D$ is read from the first omitted block obtained by expanding one order further.
\end{itemize}
The whole expression is re-expanded with a larger $K_w$ until the final precision reaches the requested cutoff. The result is reported as the retained blocks plus $\OO_{\sigma_*,k_*}$, where $\sigma_*$ is the exponent of the first omitted block if one lies below the final precision (its complete coefficient having been computed), and otherwise the precision pair $(P,D)$ itself.

\begin{proposition}[Soundness]
\label{prop:soundness}
If every leaf and every Taylor, Laurent, Puiseux or fallback expansion supplied by \wl{Series} satisfies its asserted branch and remainder contract, then the root satisfies \eqref{eq:triple}, and the reported remainder class is valid.
\end{proposition}
\begin{proof}
Induct over the expression tree. For sums and products, use \eqref{eq:O-arithmetic} for propagated errors and add the explicit discarded blocks to the remainder; in particular, their logarithmic degrees at the precision boundary cannot be ignored. For unit series, apply Lemma~\ref{lem:tail} to the finite jet's positive generators and the mean value estimate $\Phi(1+U)-\Phi(1+U_{\mathrm{known}})=\OO(U-U_{\mathrm{known}})$ for its unknown part. Laurent and Puiseux heads are reduced to these same operations on a known signed increment. The final remainder statement follows from the refined tail bound whenever a complete first omitted block has been computed below the precision limit.
\end{proof}

When \wl{Series} supplies a Taylor, Laurent, Puiseux or fallback expansion, its branch and remainder assertions remain inputs to the soundness proposition. Observing the degree of one omitted block alone does not prove that an arbitrary unknown tail satisfies the resulting bound. The structural operations preserve valid assertions, but the package does not prove every analytic contract of a special function or a user-supplied remainder. If cancellation has hidden a leading increment needed by a pole or branch point, the working precision must increase before composition.

\subsection{The second question}

For $F(x)=(1+x+x^{\sqrt2})^{\sqrt2}$ at $x\to+\infty$ the local variable is $w=1/x$ and
\[
 F=w^{-2}\bigl(1+w^{\sqrt2-1}+w^{\sqrt2}\bigr)^{\sqrt2}=w^{-2}\sum_{k\ge0}\binom{\sqrt2}{k}\bigl(w^{\sqrt2-1}+w^{\sqrt2}\bigr)^k ,
\]
an instance of item (1) of Theorem~\ref{thm:closure} with gaps $\sqrt2-1$ and $\sqrt2$. The exponents of $x$ are $2-(\sqrt2-1)k_1-\sqrt2k_2$; ordering the semigroup elements gives \eqref{eq:answer3}, and the first omitted exponent is $7-5\sqrt2$ (multi-index $(5,0)$). The workaround in the question, which peels off one leading term per limit computation, produces the same seven terms; the evaluator supplies them with the remainder $\OO(x^{7-5\sqrt2})$ and can continue to larger finite cutoffs within its resource budget.

\section{The model and the real branch}
\label{sec:branch}

\subsection{The normal form}

Let $u\to0^+$ be the local variable (later $u=x-x_0$, $x_0-x$, $1/x$ or $-1/x$). The central object is a finite \emph{power--log germ}
\begin{equation}
 f(u)=y_0+a\,u^p\Bigl(1+\sum_{i=1}^m u^{\delta_i}B_i(\log u)\Bigr),\qquad
 a\ne0,\ p\ne0,\ \delta_i>0,\ B_i\in\R[L]\setminus\{0\},
 \label{eq:model}
\end{equation}
with $y_0\in\R$ when $p>0$ and $y_0=0$ when $p<0$ (a constant term is then a correction with gap $-p$). Equal gaps are merged. The \emph{leading block} $au^p$ has a constant coefficient; a logarithmic factor in the leading block ($u\log u$, $u(\log u)^{-1}$) changes the problem (Section~\ref{sec:boundaries}). Everything below is stated for this exact finite germ; functions that are only known through a truncated expansion are handled in Section~\ref{sec:transport}.

The \emph{uniformizer} is
\begin{equation}
 z=\Bigl(\frac{v}{a}\Bigr)^{1/p},\qquad v=y-y_0\ (p>0),\quad v=y\ (p<0),\qquad L=\log z=\frac{\log(v/a)}{p},
 \label{eq:uniformizer}
\end{equation}
so that $z\to0^+$ in every case: for $p>0$ because $v\to0$ with $v/a>0$, for $p<0$ because $v\to\pm\infty$ with the sign of $a$. The role of $z$ is that the inverse satisfies $u\sim z$.

\subsection{Existence and uniqueness of the local inverse}

\begin{theorem}[Positive local branch]
\label{thm:branch}
For the germ \eqref{eq:model} there is $\varepsilon>0$ such that $f$ is strictly monotone on $(0,\varepsilon)$, with
\begin{equation}
 f(u)-y_0\sim a\,u^p,\qquad f'(u)\sim a\,p\,u^{p-1}\qquad(u\to0^+).
 \label{eq:leading}
\end{equation}
Its inverse $g$, defined on the image of $(0,\varepsilon)$, is the unique function with $g>0$, $g\to0$ and $f(g(y))=y$ there, and it satisfies
\begin{equation}
 g(y)\sim z,\qquad g'(y)\sim\frac{z}{p\,v}\qquad(z\to0^+).
 \label{eq:inverse-leading}
\end{equation}
\end{theorem}
\begin{proof}
Write $R(u)=\sum_iu^{\delta_i}B_i(\log u)$. By Lemma~\ref{lem:dominance}, $R(u)\to0$ and, since $uR'(u)=\sum_iu^{\delta_i}(\delta_iB_i+B_i')(\log u)$ by Lemma~\ref{lem:euler}, also $uR'(u)\to0$. Hence
\[
 f'(u)=apu^{p-1}\Bigl(1+R(u)+\frac{uR'(u)}{p}\Bigr)=apu^{p-1}(1+\oo(1)),
\]
which has the sign of $ap$ on some $(0,\varepsilon)$; this gives monotonicity and \eqref{eq:leading}. The inverse exists on the image, tends to $0$ because $f$ is a homeomorphism of $(0,\varepsilon)$ onto its image, and is the only solution of $f(u)=y$ in $(0,\varepsilon)$. Substituting $u=g(y)$ into $f(u)-y_0\sim au^p$ gives $v\sim ag(y)^p$, i.e.\ $g\sim z$; and $g'=1/f'(g)\sim1/(apg^{p-1})=g/(pag^p)\sim z/(pv)$.
\end{proof}

The theorem is local. Whether the local branch extends to a global inverse is a separate question, answered for the two examples of \cite{q1} as follows.

\begin{proposition}[Global invertibility of the examples]
\label{prop:global}
$f_1(x)=x+x^2(1+\log x)$ and $f_2(x)=x+x^{\sqrt2}$ are strictly increasing bijections of $(0,\infty)$ onto itself, with
\begin{equation}
 f_1'(x)=1+x(3+2\log x)\ \ge\ m_1:=1-2e^{-5/2}>0.835,\qquad f_2'(x)=1+\sqrt2\,x^{\sqrt2-1}>1.
 \label{eq:global-slopes}
\end{equation}
Their inverses $g_1,g_2$ are $C^1$ up to $y=0$ with $g_i(0)=0$, $g_i'(0)=1$, but $g_1$ is not $C^2$ at $0$: $g_1''(y)=-f_1''(g_1)/f_1'(g_1)^3\sim-(5+2\log y)\to+\infty$.
\end{proposition}
\begin{proof}
$x(3+2\log x)$ has derivative $5+2\log x$, minimum at $x=e^{-5/2}$ with value $-2e^{-5/2}$; the limits of $f_1$ at $0$ and $\infty$ are $0$ and $\infty$. The rest is immediate.
\end{proof}

So for both examples the branch of Theorem~\ref{thm:branch} is \emph{the} real inverse, and the expansions \eqref{eq:answer1}--\eqref{eq:answer2} describe it near $0$; the absence of a second-order Taylor expansion of $g_1$ at $0$ is exactly why $\log y$ must appear at $y^2$.

\subsection{The complex germ returned by \wl{InverseFunction}}

The equation $f_1(x)=0$, $x\ne0$, reads $1+x(1+\log x)=0$. With $t=1+\log x$ it becomes $te^t=-e$, so $t=\W_k(-e)$ and
\begin{equation}
 c_k=e^{\W_k(-e)-1}=-\frac1{\W_k(-e)},\qquad f_1'(c_k)=c_k-1,\qquad f_1''(c_k)=3+2\W_k(-e),
 \label{eq:complex-root}
\end{equation}
where the derivative identities follow from $c_k\W_k(-e)=-1$, provided the logarithm near $c_k$ is chosen to have value $\W_k(-e)-1$ there. This branch qualification matters: with the fixed principal logarithm used by \wl{Log}, the relevant roots are the conjugate pair $k=0,-1$; other $k$ require other local logarithm branches. For each compatible choice the derivative is nonzero, and the analytic inverse function theorem gives an inverse germ through $c_k$. The reported \wl{Series[InverseFunction[...][y], \{y, 0, 2\}]} result chooses $k=-1$:
\[
 c+\frac{y}{c-1}-\frac{f_1''(c)}{2f_1'(c)^3}y^2+\cdots=c+\frac{y}{c-1}-\frac{3+2\W_{-1}(-e)}{2(c-1)^3}y^2+\OO(y^3),
\]
which is \eqref{eq:complex-germ}. Nothing is wrong with this answer: it is the inverse of a different branch of $f_1^{-1}$, the one through a complex root, and adding \wl{Assumptions -> y > 0} cannot repair it, because the reality of $y$ says nothing about which preimage is intended. The real branch is characterised by $g(y)\to0$, and that condition, not an assumption on $y$, is what an inversion routine must be given; the package builds the branch of Theorem~\ref{thm:branch} directly and never asks \wl{InverseFunction} to choose.

\section{Lagrange--B\"urmann inversion with a marker}
\label{sec:lagrange}

The classical Lagrange inversion theorem expands the inverse of an analytic map at a point where its derivative does not vanish. Our maps are not analytic at $0$, and the inverse has no Taylor series there; but at every \emph{positive} point the map is analytic, and the trick is to invert a perturbation $u+\eta h(u)=q$ around the positive point $q$ in the perturbation parameter $\eta$. The logarithms and irrational powers in $h$ are then analytic functions of $u$ near $q>0$, and the derivatives that appear are computed by the Euler identity without ever evaluating anything at $0$.

\subsection{The marker theorem}

\begin{theorem}[Lagrange--B\"urmann with a marker]
\label{thm:marker}
Let $q>0$ and let $h$ and $\Phi$ be analytic in a complex neighbourhood of $q$. For $|\eta|$ sufficiently small the equation
\begin{equation}
 u+\eta\,h(u)=q
 \label{eq:marker-equation}
\end{equation}
has a unique solution $u(\eta)$ near $q$ with $u(0)=q$, it is analytic in $\eta$, and
\begin{equation}
 \Phi\bigl(u(\eta)\bigr)=\Phi(q)+\sum_{N\ge1}\frac{(-\eta)^N}{N!}\,\frac{\dd^{N-1}}{\dd q^{N-1}}\Bigl[\Phi'(q)\,h(q)^N\Bigr].
 \label{eq:marker-formula}
\end{equation}
The same identity holds coefficientwise in $\eta$ for arbitrary formal power series $h,\Phi$ with coefficients in any commutative $\Q$-algebra, i.e.\ as an identity of formal power series in $\eta$ whose coefficients are polynomials in finitely many derivatives of $h$ and $\Phi$ at $q$.
\end{theorem}
\begin{proof}
Choose $\rho>0$ so small that $h$ and $\Phi$ are analytic on the closed disc $|s-q|\le\rho$, and let $M=\max_{|s-q|=\rho}|h(s)|$. For $|\eta|<\rho/M$ we have $|\eta h(s)|<|s-q|$ on the circle $C=\{|s-q|=\rho\}$, so by Rouch\'e's theorem $F_\eta(s)=s-q+\eta h(s)$ has exactly one zero $u(\eta)$ in the disc, and by the argument principle
\[
 \Phi(u(\eta))=\frac1{2\pi i}\oint_C\Phi(s)\,\frac{F_\eta'(s)}{F_\eta(s)}\,\dd s
 =\frac1{2\pi i}\oint_C\Phi(s)\,\frac{1+\eta h'(s)}{s-q+\eta h(s)}\,\dd s .
\]
On $C$ the geometric series $\frac1{s-q+\eta h(s)}=\sum_{N\ge0}\frac{(-\eta)^Nh(s)^N}{(s-q)^{N+1}}$ converges uniformly, so
\[
 \Phi(u(\eta))=\sum_{N\ge0}(-\eta)^N\Bigl[\frac1{2\pi i}\oint_C\frac{\Phi(s)h(s)^N}{(s-q)^{N+1}}\,\dd s
 -\frac1{2\pi i}\oint_C\frac{\Phi(s)h'(s)h(s)^N\,(-\eta)}{(s-q)^{N+1}}\,\dd s\Bigr].
\]
By Cauchy's formula the first integral is $\frac1{N!}\frac{\dd^N}{\dd q^N}[\Phi h^N](q)$ and the second, after shifting the index $N\to N-1$, contributes $-\frac1{(N-1)!}\frac{\dd^{N-1}}{\dd q^{N-1}}[\Phi h'h^{N-1}](q)$ to the coefficient of $(-\eta)^N$ for $N\ge1$. Since $\frac{\dd}{\dd q}[\Phi h^N]=\Phi'h^N+N\Phi h'h^{N-1}$, the coefficient of $(-\eta)^N$ is
\[
 \frac1{N!}\frac{\dd^{N-1}}{\dd q^{N-1}}\bigl[\Phi'h^N+N\Phi h'h^{N-1}\bigr]-\frac1{(N-1)!}\frac{\dd^{N-1}}{\dd q^{N-1}}\bigl[\Phi h'h^{N-1}\bigr]
 =\frac1{N!}\frac{\dd^{N-1}}{\dd q^{N-1}}\bigl[\Phi'h^N\bigr],
\]
which is \eqref{eq:marker-formula}; the constant term is $\Phi(q)$. Analyticity of $u(\eta)$ follows from the implicit function theorem at $(q,0)$, where $\partial_uF_0=1$. For the formal statement, both sides of \eqref{eq:marker-formula} have coefficients that are universal polynomials in the Taylor coefficients of $h$ and $\Phi$ at $q$; two polynomials that agree for all analytic data agree identically (they agree on a Zariski-dense set), so the identity holds for arbitrary formal coefficients.
\end{proof}

\begin{remark}
Other proofs exist \cite{gessel,surya,dlmf}; the one above is chosen because it exhibits the analytic content: the solution is the unique root in a disc around $q$, a fact reused in Section~\ref{sec:convergence} to obtain explicit constants. The formal statement is what makes the coefficient formula of Section~\ref{sec:coefficients} an identity, independent of any convergence question.
\end{remark}

\subsection{Consequences}

\paragraph{Near-identity form.} With $\Phi(q)=q$ and $\eta=1$,
\begin{equation}
 f(u)=u+h(u)\quad\Longrightarrow\quad g(y)=y+\sum_{N\ge1}\frac{(-1)^N}{N!}\frac{\dd^{N-1}}{\dd y^{N-1}}\bigl[h(y)^N\bigr],
 \label{eq:near-identity}
\end{equation}
formally; Sections~\ref{sec:convergence} and~\ref{sec:transport} say when setting $\eta=1$ is legitimate as an asymptotic statement. The formula \eqref{eq:near-identity} is a complete answer for the log example: with $h(y)=y^2(1+\log y)$, the $N$-th term is $\frac{(-1)^N}{N!}\partial_y^{N-1}[y^{2N}(1+\log y)^N]$, computed by Lemma~\ref{lem:euler}; for $N=2$, $\frac12\partial_y[y^4(1+L)^2]=y^3\bigl(2(1+L)^2+(1+L)\bigr)=y^3(2L^2+5L+3)$.

\paragraph{Several markers.} If $h=\sum_i\eta_ih_i$ with independent markers, expanding $h^N$ multinomially in \eqref{eq:marker-formula} gives the coefficient of $\prod\eta_i^{k_i}$, $|\bk|=N$, as
\begin{equation}
 \frac{(-1)^N}{\bk!}\,\frac{\dd^{N-1}}{\dd q^{N-1}}\Bigl[\Phi'(q)\prod_ih_i(q)^{k_i}\Bigr],
 \label{eq:multi-marker}
\end{equation}
because the multinomial coefficient $N!/\bk!$ cancels the $1/N!$.

\paragraph{Transport of observables.} Formula \eqref{eq:marker-formula} gives at once the expansion of any analytic function of the inverse. With $\Phi(u)=u^r$ one obtains powers of the inverse; with $\Phi=\log$ its logarithm; differentiating the identity for $u^r$ with respect to $r$ at $r=0$ gives the logarithm again. These are used in Section~\ref{sec:coefficients}.

\paragraph{General core.} Suppose $F=F_0+R$ where $F_0$ has a known inverse $\phi$ on the branch of interest and $R$ is a perturbation. Put $t=F_0(x)$, so that $F(x)=y$ becomes $t+\eta\,R(\phi(t))=y$ with $\eta=1$; Theorem~\ref{thm:marker} with $\Phi=\phi$ yields
\begin{equation}
 x=\phi(y)+\sum_{N\ge1}\frac{(-1)^N}{N!}\frac{\dd^{N-1}}{\dd y^{N-1}}\Bigl[\phi'(y)\,R\bigl(\phi(y)\bigr)^N\Bigr].
 \label{eq:general-core}
\end{equation}
This is the ``phase-coordinate'' inversion of \cite{proveit}; it is implemented as \wl{PerturbativeInverse[$\phi$, R, \{x, y\}, n]} and is the natural tool when the dominant part of $F$ is not a monomial (Section~\ref{sec:boundaries}). For the monomial core $F_0(x)=ax^p$, $\phi(y)=(y/a)^{1/p}$, it is equivalent to the substitution $t=x^p$ used in the next section.

\section{The coefficient theorem}
\label{sec:coefficients}

\subsection{Statement}

Let $f$ be the germ \eqref{eq:model}, $z$ the uniformizer \eqref{eq:uniformizer}, $L=\log z$, $\D=\dd/\dd L$. For a multi-index $\bk\in\N^m$ write $n=|\bk|$, $\sigma=\bk\cdot\bdelta$ and $B^{\bk}=\prod_iB_i^{k_i}$.

\begin{theorem}[Euler-operator coefficient formula]
\label{thm:coefficients}
For every real $r$, the branch $g$ of Theorem~\ref{thm:branch} satisfies, formally and (Theorem~\ref{thm:convergence}) as a convergent expansion for small $z$,
\begin{equation}
 \boxed{\;g(y)^r=z^r+\sum_{\bk\ne0}z^{r+\bk\cdot\bdelta}\,C^{(r)}_{\bk}(L),\qquad
 C^{(r)}_{\bk}(L)=\frac{(-1)^{n}\,r}{p^{n}\,\bk!}\,\prod_{j=1}^{n-1}\bigl(\D+r+\sigma+pj\bigr)\,B^{\bk}(L).\;}
 \label{eq:master}
\end{equation}
The empty product ($n=1$) is the identity. Multi-indices of equal weight are summed: the block of $g^r$ at $z^{r+\sigma}$ is $C^{(r)}_\sigma=\sum_{\bk\cdot\bdelta=\sigma}C^{(r)}_{\bk}$. In particular
\begin{equation}
\begin{aligned}
 g(y)&=z\Bigl(1+\sum_{\bk\ne0}z^{\bk\cdot\bdelta}C_{\bk}(L)\Bigr),\qquad C_{\bk}=C^{(1)}_{\bk},\\
 \log g(y)&=L+\sum_{\bk\ne0}z^{\bk\cdot\bdelta}\frac{(-1)^n}{p^n\bk!}\prod_{j=1}^{n-1}(\D+\sigma+pj)B^{\bk}(L).
\end{aligned}
 \label{eq:master-1-log}
\end{equation}
\end{theorem}
\begin{proof}
Put $t=u^p$ and $q=v/a$ (so $q=z^p$). The equation $f(u)=y$ becomes
\[
 q=t+\sum_it^{1+\delta_i/p}B_i\Bigl(\frac{\log t}{p}\Bigr),
\]
which is \eqref{eq:marker-equation} with markers: $t+\sum_i\eta_ih_i(t)=q$, $h_i(t)=t^{1+\delta_i/p}B_i(\log t/p)$, at $\eta_i=1$. The observable is $\Phi(t)=t^{r/p}=u^r$. All $h_i$ and $\Phi$ are analytic near any $q>0$. By \eqref{eq:multi-marker} the coefficient of $\prod\eta_i^{k_i}$ is
\[
 \frac{(-1)^n}{\bk!}\frac{\dd^{n-1}}{\dd q^{n-1}}\Bigl[\frac rp\,q^{r/p-1}\prod_i\bigl(q^{1+\delta_i/p}B_i(\tfrac{\log q}p)\bigr)^{k_i}\Bigr]
 =\frac{(-1)^n r}{p\,\bk!}\frac{\dd^{n-1}}{\dd q^{n-1}}\Bigl[q^{\,n-1+(r+\sigma)/p}\,B^{\bk}\bigl(\tfrac{\log q}p\bigr)\Bigr].
\]
Apply the repeated Euler identity \eqref{eq:euler} in the variable $q$, with $\Lambda=\log q$ and the polynomial $\tilde P(\Lambda)=B^{\bk}(\Lambda/p)$: the derivative of order $n-1$ of $q^b\tilde P(\Lambda)$, $b=n-1+(r+\sigma)/p$, is $q^{b-n+1}\prod_{j=0}^{n-2}(b-j+\dd/\dd\Lambda)\tilde P$. Since $L=\Lambda/p$, $\dd/\dd\Lambda=\D/p$, and $b-j=\bigl(r+\sigma+p(n-1-j)\bigr)/p$; hence each factor is $\frac1p\bigl(\D+r+\sigma+p(n-1-j)\bigr)$, and with $j'=n-1-j$ running from $1$ to $n-1$ the product is $p^{-(n-1)}\prod_{j'=1}^{n-1}(\D+r+\sigma+pj')B^{\bk}(L)$. Finally $q^{b-n+1}=q^{(r+\sigma)/p}=z^{r+\sigma}$, and the prefactor $r/p$ combines with $p^{-(n-1)}$ into $r/p^n$. The formal validity is Theorem~\ref{thm:marker}'s formal clause, applied to the finitely many markers; setting all $\eta_i=1$ is legitimate coefficientwise because by Lemma~\ref{lem:finite} only finitely many multi-indices contribute to any exponent. The formula for $\log g$ is the derivative of \eqref{eq:master} with respect to $r$ at $r=0$, where the factor $r$ vanishes and only the term with the factor differentiated survives; equivalently it is Theorem~\ref{thm:marker} with $\Phi=\log$.
\end{proof}

\begin{remark}
The normalization in \eqref{eq:master} contains four essential features: sign $(-1)^n$, the multi-index factorial $\bk!$ (not $n!$: the multinomial coefficient from expanding $h^n$ has been absorbed), the factor $p^{-n}$ (one $1/p$ from $\Phi'$ and one from each of the $n-1$ derivatives, through $\dd/\dd\Lambda=\D/p$), and the operator constants $r+\sigma+pj$. Forgetting the substitution $L=\log(v/a)/p$ in the coefficients, or the factor $p^{-n}$, are the two classical mistakes for nonunit $p$.
\end{remark}

\subsection{First orders and cross terms}

For a single block ($m=1$, gap $\delta$, polynomial $B$):
\begin{equation}
 g=z-\frac{z^{1+\delta}}{p}B+\frac{z^{1+2\delta}}{2p^2}\bigl((1+2\delta+p)B^2+2BB'\bigr)
 -\frac{z^{1+3\delta}}{6p^3}(\D+1+3\delta+p)(\D+1+3\delta+2p)B^3+\cdots .
 \label{eq:single-block}
\end{equation}
For two blocks the mixed multi-index $e_i+e_j$ contributes
\begin{equation}
 \frac{z^{1+\delta_i+\delta_j}}{p^2}\bigl(\D+1+\delta_i+\delta_j+p\bigr)(B_iB_j)
 =\frac{z^{1+\delta_i+\delta_j}}{p^2}\bigl((1+\delta_i+\delta_j+p)B_iB_j+B_i'B_j+B_iB_j'\bigr),
 \label{eq:cross}
\end{equation}
which shows why the inverse of a sum of perturbations is not the sum of the separately computed inverse corrections.

\subsection{Degrees and leading logarithms}

\begin{proposition}[Degrees]
\label{prop:degrees}
Let $d_i=\deg B_i$ and $D(\bk)=\sum_ik_id_i$. Then $\deg C^{(r)}_{\bk}\le D(\bk)$, with coefficient at the largest possible degree
\begin{equation}
 [L^{D(\bk)}]\,C^{(r)}_{\bk}=\frac{(-1)^nr}{p^n\bk!}\prod_{j=1}^{n-1}(r+\sigma+pj)\prod_i\bigl([L^{d_i}]B_i\bigr)^{k_i},
 \label{eq:leading-coefficient}
\end{equation}
which is nonzero exactly when $r\ne0$ and $r+\sigma+pj\ne0$ for $1\le j\le n-1$ (always true for $r>0$, $p>0$). Under those conditions the degree is exactly $D(\bk)$. Zero operator constants differentiate and may lower the degree or annihilate the polynomial, while $r=0$ annihilates every correction. An aggregated block $C^{(r)}_\sigma$ has degree at most $\max\{D(\bk):\bk\cdot\bdelta=\sigma\}$; equality requires a nonzero sum of its top-degree coefficients.
\end{proposition}
\begin{proof}
Each factor $\D+c$ with $c\ne0$ preserves the degree and multiplies the top coefficient by $c$; a factor with $c=0$ differentiates. The product $B^{\bk}$ has degree $D(\bk)$, and the prefactor contains $r$.
\end{proof}

For the log example ($a=p=\delta=1$, $B=1+L$, single index $n$) the formula reads
\begin{equation}
\begin{aligned}
 P_n(L)&=\frac{(-1)^n}{n!}\prod_{j=n+2}^{2n}(\D+j)\,(1+L)^n,\\
 [L^n]P_n&=\frac{(-1)^n}{n!}\prod_{j=n+2}^{2n}j=(-1)^n\frac{(2n)!}{n!\,(n+1)!}=(-1)^n\Cat_n,
\end{aligned}
 \label{eq:log-example}
\end{equation}
the signed Catalan numbers. The next coefficient is also explicit: selecting the derivative in exactly one factor of the product gives
\begin{equation}
 [L^{n-1}]P_n=(-1)^n\Cat_n\,n\,\bigl(1+H_{2n}-H_{n+1}\bigr),\qquad H_k=\textstyle\sum_{j\le k}1/j,
 \label{eq:harmonic}
\end{equation}
because $\prod_{j=n+2}^{2n}(\D+j)(1+L)^n$ has $L^{n-1}$-coefficient $\bigl(\prod j\bigr)\bigl(n+\sum_{j=n+2}^{2n}n/j\bigr)$ (the term $n$ from expanding $(1+L)^n$ and the terms $n/j$ from applying $\D$ in the factor $j$). Every $P_n$ is divisible by $1+L$: in the variable $t=1+L$ the operator product $\prod_{j=n+2}^{2n}(\D+j)$ consists of $n-1$ factors, each of which either differentiates (lowering the power of $t$ by one) or multiplies by a constant, so every monomial of $\prod(\D+j)\,t^n$ has degree at least $n-(n-1)=1$ in $t$. Thus
\begin{gather*}
 P_1=-(1+L),\qquad P_2=(1+L)(2L+3),\qquad P_3=-\tfrac12(1+L)(10L^2+31L+23),\\
 P_4=\tfrac16(1+L)(84L^3+398L^2+607L+299),
\end{gather*}
which are the blocks of \eqref{eq:answer1}.

\subsection{Constant coefficients: Fuss--Catalan numbers}

When every $B_i$ is a constant $b_i$, the operators act as scalars and
\begin{equation}
 C^{(r)}_{\bk}=\frac{(-1)^nr}{p^n\bk!}\prod_{j=1}^{n-1}(r+\sigma+pj)\prod_ib_i^{k_i}
 =\frac{r}{r+\sigma}\binom{-(r+\sigma)/p}{n}\frac{n!}{\bk!}\prod_ib_i^{k_i}.
 \label{eq:constant-coefficients}
\end{equation}
The binomial expression in \eqref{eq:constant-coefficients} is read by its removable limit when $r+\sigma=0$; the operator-product expression has no division by $r+\sigma$ and remains directly valid. For $f(x)=x+bx^\alpha$ ($p=1$, $\delta=\alpha-1$, $r=1$) this is the generalised Fuss--Catalan law
\begin{equation}
\begin{aligned}
 g(y)&=\sum_{n\ge0}c_n\,b^n\,y^{1+n(\alpha-1)},\\
 c_n&=\frac{(-1)^n}{n!}\prod_{j=0}^{n-2}(n\alpha-j)=\frac{(-1)^n}{1+n(\alpha-1)}\binom{n\alpha}{n}=(-1)^n\frac{\Gamma(n\alpha+1)}{\Gamma(n+1)\Gamma(n(\alpha-1)+2)},
\end{aligned}
 \label{eq:fuss-catalan}
\end{equation}
valid for every real $\alpha>1$ (no rationality is used), with $c_0=1$ understood separately in the product formula. For $\alpha=2$ the $c_n$ are the signed Catalan numbers and, when $b\ne0$, $g=(\sqrt{1+4by}-1)/(2b)$; the limit at $b=0$ is $g=y$. For $b=1$ and $\alpha=\sqrt2$,
\begin{gather*}
 c_1=-1,\qquad c_2=\sqrt2,\qquad c_3=-\frac{6-\sqrt2}{2},\\
 c_4=\frac{17\sqrt2}{3}-4,\qquad c_5=\frac{51\sqrt2}{4}-\frac{305}{12},\qquad c_6=\frac{322\sqrt2}{5}-77,
\end{gather*}
which gives \eqref{eq:answer2}. In this one-gap case $g(y)=yV(by^{\alpha-1})$ with $V$ the solution of $V+tV^\alpha=1$, $V(0)=1$, analytic at $t=0$ with radius of convergence
\begin{equation}
 R_\alpha=\frac{(\alpha-1)^{\alpha-1}}{\alpha^\alpha},\qquad |c_n|\sim\frac{\sqrt\alpha}{\sqrt{2\pi}\,(\alpha-1)^{3/2}}\,n^{-3/2}R_\alpha^{-n},
 \label{eq:radius}
\end{equation}
by Stirling's formula, so the expansion \eqref{eq:answer2} converges (and equals the positive inverse, by analytic continuation of the implicit equation) for $0<y<R_{\sqrt2}^{1/(\sqrt2-1)}\approx0.12686$. The singularity at $t=-R_\alpha$ is the fold $V=\alpha/(\alpha-1)$ of the implicit equation on the negative axis.

\subsection{Formal existence, uniqueness, and iterative construction}
\label{sec:formal}

Theorem~\ref{thm:coefficients} produces the coefficients; the following independent argument shows that they are the only possible ones and provides two alternative algorithms. Write $g=z(1+U)$ with $U$ an unknown jet of positive valuation. Then $f(g)=y$ is equivalent to
\begin{equation}
 E(U):=(1+U)^p\Bigl(1+\sum_iz^{\delta_i}(1+U)^{\delta_i}B_i\bigl(L+\log(1+U)\bigr)\Bigr)-1=0,
 \label{eq:normalized}
\end{equation}
an equation in the jet algebra (Lemma~\ref{lem:composition} evaluates every term).

\begin{proposition}[Triangular structure]
\label{prop:triangular}
Equation \eqref{eq:normalized} has exactly one formal solution $U$ of positive valuation with support in the semigroup $S$ generated by the gaps. Its coefficient at $z^\sigma$ is determined by the coefficients at smaller weights through one division by $p$.
\end{proposition}
\begin{proof}
$E(U)=pU+\sum_iz^{\delta_i}B_i(L)+(\text{terms of weight}\ge\min(2\val U,\val U+\delta_*))$, because $(1+U)^p=1+pU+\OO(U^2)$ and every correction carries a factor $z^{\delta_i}$. Order the elements of $S$ increasingly (Lemma~\ref{lem:finite}); at a new weight $\sigma$ the coefficient of $z^\sigma$ in $E(U)$ is $pU_\sigma$ plus an expression in the $U_{\sigma'}$ with $\sigma'<\sigma$, which fixes $U_\sigma$. If $U,V$ are two solutions with $U-V$ of least weight $\sigma$, then $E(U)-E(V)=(U-V)(p+\text{positive weight})$ has a nonzero block at $z^\sigma$, a contradiction.
\end{proof}

\paragraph{Picard iteration.} Rewriting \eqref{eq:normalized} as $U=\Psi(U):=\bigl(1+\sum_iz^{\delta_i}(1+U)^{\delta_i}B_i(L+\log(1+U))\bigr)^{-1/p}-1$, the map $\Psi$ gains one gap per step: $\val(\Psi(U)-\Psi(V))\ge\val(U-V)+\delta_*$. Hence $U_{k+1}=\Psi(U_k)$, $U_0=0$, agrees with the solution below weight $(k+1)\delta_*$, and $\lceil H/\delta_*\rceil$ iterations give all blocks below $H$.

\paragraph{Newton iteration.} With $E'(U)=p(1+U)^{p-1}+\sum_iz^{\delta_i}(1+U)^{p+\delta_i-1}\bigl((p+\delta_i)B_i+B_i'\bigr)(L+\log(1+U))$, a unit of the jet algebra with constant leading coefficient $p$, the iteration $U_{k+1}=U_k-E(U_k)/E'(U_k)$ doubles the valuation of the error: if $\val(U_k-U)\ge s$ then $\val(U_{k+1}-U)\ge2s$, because $E(U_k)=E'(U)(U_k-U)+\OO((U_k-U)^2)$ and $E'(U_k)^{-1}=E'(U)^{-1}(1+\OO(U_k-U))$. Starting from $U_0=0$ with $s=\delta_*$, about $\log_2(H/\delta_*)$ iterations suffice. Formal uniqueness implies that Newton iteration and the multi-index formula give the same coefficients at every weight.

\section{Convergence}
\label{sec:convergence}

A formal identity such as \eqref{eq:master} says nothing, by itself, about the function $g$: the series might diverge, or converge to something else. This section proves that the expansion converges absolutely for small $z$ and equals the branch of Theorem~\ref{thm:branch}, first by an implicit-function argument that also yields the coefficient bounds needed for the remainder theory, then by a Rouch\'e argument that yields explicit constants.

\subsection{Convergence in finitely many small parameters}

\begin{theorem}[Convergence]
\label{thm:convergence}
Let $f$ be the germ \eqref{eq:model} and $g$ its branch. For every real $r$ the series \eqref{eq:master} converges absolutely for all sufficiently small $z>0$, its sum is $g(y)^r$, and there are constants $M,K$ such that
\begin{equation}
 \bigl|C^{(r)}_{\bk}(L)\bigr|\le M\,K^{|\bk|}\,\M(z)^{D(\bk)}\qquad(L\le\log z_1),\quad D(\bk)=\textstyle\sum_ik_i\deg B_i .
 \label{eq:coefficient-bound}
\end{equation}
In particular $g^r\in\mathcal E$ (Definition~\ref{def:pl-expansion}) with support $r+S$.
\end{theorem}
\begin{proof}
Write $g=z\,w$ with $w\to1$; by \eqref{eq:normalized} the unknown $w$ solves
\[
 w^p\Bigl(1+\sum_iz^{\delta_i}w^{\delta_i}B_i(L+\log w)\Bigr)=1 .
\]
Expand each $B_i(L+\log w)=\sum_{j=0}^{d_i}L^jA_{ij}(\log w)$ with $A_{ij}(v)=\sum_{s\ge j}b_{is}\binom sjv^{s-j}$, where $B_i=\sum_sb_{is}L^s$. Introduce independent complex parameters $T_{ij}$ in place of $z^{\delta_i}L^j$ and consider
\begin{equation}
 \mathcal E(w,T)=w^p\Bigl(1+\sum_{i,j}T_{ij}\,w^{\delta_i}A_{ij}(\log w)\Bigr)-1 .
 \label{eq:parametric}
\end{equation}
On the disc $|w-1|<1$ with the principal logarithm, $\mathcal E$ is analytic in $(w,T)$; $\mathcal E(1,0)=0$ and $\partial_w\mathcal E(1,0)=p\ne0$. The analytic implicit function theorem gives an analytic $W(T)$ on a polydisc $|T_{ij}|<\rho$ with $W(0)=1$ and $\mathcal E(W(T),T)=0$, and $W^r$ is analytic there too. Its Taylor series $W(T)^r=\sum_{\bnu}c_{\bnu}T^{\bnu}$ converges absolutely on the smaller polydisc $|T_{ij}|\le\rho/2$, with the Cauchy estimates $|c_{\bnu}|\le M(2/\rho)^{|\bnu|}$.

Now specialise $T_{ij}=z^{\delta_i}L^j$. Since $z^{\delta_i}\M(z)^{d_i}\to0$, the point lies in the polydisc for small $z$, and $W(T(z))$ is real (all data are real and the solution is unique) and close to $1$, so $zW(T(z))$ is a solution of $f(u)=y$ with $u\sim z$; by Theorem~\ref{thm:branch} it is $g(y)$. Group the absolutely convergent series $\sum c_{\bnu}\prod_{ij}(z^{\delta_i}L^j)^{\nu_{ij}}$ by the multi-index $k_i=\sum_j\nu_{ij}$: the terms with a given $\bk$ are finitely many (at most $\prod_i(d_i+1)^{k_i}$), and their sum is $z^{\bk\cdot\bdelta}$ times a polynomial in $L$ of degree at most $\sum_j j\nu_{ij}\le D(\bk)$. Absolute convergence permits this regrouping, so $g^r=z^r\sum_{\bk}z^{\bk\cdot\bdelta}\tilde C_{\bk}(L)$ with the bound \eqref{eq:coefficient-bound} for $K=(2/\rho)\max_i(d_i+1)$. Finally the $\tilde C_{\bk}$ are the coefficients $C^{(r)}_{\bk}$ of \eqref{eq:master}: both are formal solutions of the same triangular problem (Proposition~\ref{prop:triangular} for $r=1$; for general $r$ apply the observable $w\mapsto w^r$), which has a unique solution.
\end{proof}

\begin{remark}
The theorem does not give a radius of convergence in $y$: logarithms are present, and the expansion is not a power series in any single variable (except in the one-gap constant-coefficient case, \eqref{eq:radius}). It gives convergence of a multivariate power series after the substitution $T_{ij}=z^{\delta_i}(\log z)^j$, uniformly on $0<z<z_1$. Nor is anything uniform in the parameters $a,p,\delta_i,B_i$: as a gap $\delta\to0$ the hierarchy collapses, and the expansion of $x+x^\alpha$ at $\alpha\downarrow1$ has radius \eqref{eq:radius} tending to $0$.
\end{remark}

\subsection{An explicit conditional majorant}

The implicit-function argument has no explicit constants. Rouch\'e's theorem in the marker equation gives them, at the price of a hypothesis that must be checked at each $y$.

\begin{theorem}[Explicit tail bound]
\label{thm:majorant}
Consider the near-identity germ $f(t)=t+h(t)$, $h(t)=\sum_it^{1+\alpha_i}P_i(\log t)$ with $\alpha_i>0$ and $P_i(L)=\sum_sp_{is}L^s$ (the general germ reduces to this form by $t=u^p$, $q=v/a$). Fix $0<r<1$, put $c_r=-\log(1-r)$ and
\begin{equation}
 \mathcal Q_r(t)=\frac1r\sum_it^{\alpha_i}(1+r)^{1+\alpha_i}\sum_s|p_{is}|\bigl(|\log t|+c_r\bigr)^s .
 \label{eq:Qr}
\end{equation}
Let $t>0$ satisfy $\mathcal Q_r(t)<1$. Then:
\begin{enumerate}[leftmargin=1.8em]
\item for every complex $\lambda$ with $|\lambda|<1/\mathcal Q_r(t)$ the equation $u+\lambda h(u)=t$ has exactly one solution $u(\lambda)$ in the disc $|u-t|<rt$; it is analytic in $\lambda$, $u(0)=t$, and $u(\lambda)$ is real for real $\lambda$;
\item the Lagrange terms $u_N=\frac{(-1)^N}{N!}\frac{\dd^{N-1}}{\dd t^{N-1}}[h(t)^N]$ satisfy $|u_N|\le\frac{rt}{N}\mathcal Q_r(t)^N$, and $u(\lambda)=t+\sum_{N\ge1}u_N\lambda^N$ converges absolutely for $|\lambda|<1/\mathcal Q_r(t)$, in particular at $\lambda=1$;
\item $u(1)$ is the positive solution of $f(u)=t$ in $(t(1-r),t(1+r))$, it coincides with the branch $g(t)$ for all sufficiently small $t$, and
\begin{equation}
 \Bigl|u(1)-t-\sum_{N=1}^{N_0}u_N\Bigr|\le\frac{r\,t}{N_0+1}\,\frac{\mathcal Q_r(t)^{N_0+1}}{1-\mathcal Q_r(t)} .
 \label{eq:tail-bound}
\end{equation}
\end{enumerate}
\end{theorem}
\begin{proof}
On the circle $|s-t|=rt$ write $s=t(1+\zeta)$ with $|\zeta|=r<1$. Then $|s|\le(1+r)t$ and, for the principal logarithm, $|\log s|\le|\log t|+|\log(1+\zeta)|\le|\log t|+\sum_{k\ge1}r^k/k=|\log t|+c_r$. Hence $|h(s)|\le\sum_i((1+r)t)^{1+\alpha_i}\sum_s|p_{is}|(|\log t|+c_r)^s=rt\,\mathcal Q_r(t)$. For $|\lambda|\mathcal Q_r(t)<1$ this gives $|\lambda h(s)|<rt=|s-t|$ on the circle, and Rouch\'e's theorem applied to $s-t+\lambda h(s)$ and $s-t$ shows that the former has exactly one zero in the disc, which is (1); analyticity in $\lambda$ follows from the argument-principle integral as in Theorem~\ref{thm:marker}, and reality from uniqueness together with the invariance of the equation under conjugation. The coefficients $u_N$ of the analytic function $\lambda\mapsto u(\lambda)-t$ are given by Theorem~\ref{thm:marker}; Cauchy's estimate for the $(N-1)$-st derivative of $h^N$ on the circle of radius $rt$ gives $|u_N|\le\frac1{N!}\,\frac{(N-1)!\,(rt\mathcal Q_r)^N}{(rt)^{N-1}}=\frac{rt}{N}\mathcal Q_r^N$, which is (2). For (3), $u(1)$ is real and in the disc, hence in $(t(1-r),t(1+r))$; $g(t)\sim t$ lies in that interval for small $t$ and is also a root, so uniqueness gives $u(1)=g(t)$; and the tail bound is $\sum_{N>N_0}\frac{rt}{N}\mathcal Q_r^N\le\frac{rt}{N_0+1}\sum_{N>N_0}\mathcal Q_r^N$.
\end{proof}

\begin{corollary}[Exponent cutoff]
\label{cor:cutoff-bound}
Let $G_B$ retain all multi-indices of weight $\bk\cdot\boldsymbol\alpha\le B$. Every multi-index with $|\bk|\le N_0:=\lfloor B/\alpha_{\max}\rfloor$ is retained, so under $\mathcal Q_r(t)<1$
\[
 |g(t)-G_B(t)|\le\frac{rt}{N_0+1}\,\frac{\mathcal Q_r(t)^{N_0+1}}{1-\mathcal Q_r(t)} .
\]
\end{corollary}
\begin{proof}
The omitted part of the multi-marker expansion is a subseries of $\sum_{N>N_0}u_N$ (the expansion of $u_N$ into multi-indices is absolutely convergent by the same Cauchy estimates, since $\sum_{|\bk|=N}\frac{\prod M_i^{k_i}}{\bk!}=\frac{(\sum M_i)^N}{N!}$).
\end{proof}

\begin{example}
For $f_1(x)=x+x^2(1+\log x)$ and $r=\frac12$: $\mathcal Q_{1/2}(y)=\frac92\,y\,(1+|\log y|+\log2)$, so for $y\le\frac1{100}$ (where $\mathcal Q_{1/2}\le0.284$) the expansion converges and the error after $N_0$ blocks is at most $\frac{y}{2(N_0+1)}\mathcal Q_{1/2}(y)^{N_0+1}/(1-\mathcal Q_{1/2}(y))$; at $y=10^{-2}$ and $N_0=3$ this is $1.1\cdot10^{-5}$ while the actual error is $1.5\cdot10^{-7}$. The bound is conditional (it says nothing when $\mathcal Q_r\ge1$) and not tight, but it is rigorous, cheap, and turns the asymptotic statements into inequalities at specific points.
\end{example}

\section{Remainders after a cutoff}
\label{sec:remainders}

\subsection{Exponent cutoff: the first omitted block}

Let the target cutoff be an exponent $q$ of the local variable of the result, exclusive. In terms of weights, with $r$ the observable power, a block $z^{r+\sigma}C(L)$ has exponent $(r+\sigma)/|p|$ in the local variable of $y$ (Section~\ref{sec:endpoints}), so the retained set is
\begin{equation}
 I_H=\{\bk\in\N^m:\bk\cdot\bdelta<H\},\qquad H=|p|\,q-r ,
 \label{eq:IH}
\end{equation}
a finite downward-closed set (Lemma~\ref{lem:finite}). Set $C^{(r)}_0=1$. When $H>0$ its \emph{one-step boundary} is
\begin{equation}
 \partial I_H=\{\bk+e_i:\bk\in I_H,\ 1\le i\le m\}\setminus I_H ,
 \label{eq:boundary}
\end{equation}
also finite. When $H\le0$, $I_H$ is empty and we define its boundary to be $\{0\}$, since the leading monomial is already omitted. With at least one correction gap, put
\begin{equation}
 \sigma_*=\min\{\bnu\cdot\bdelta:\bnu\in\partial I_H\},\qquad
 C^{(r)}_{\sigma_*}=\sum_{\bk\cdot\bdelta=\sigma_*}C^{(r)}_{\bk},\qquad k_*=\deg C^{(r)}_{\sigma_*}\ (k_*=0\text{ if }C^{(r)}_{\sigma_*}=0).
 \label{eq:frontier-data}
\end{equation}

\begin{lemma}
\label{lem:frontier}
$\sigma_*=\min(S\setminus[0,H))$, and every multi-index of weight $\sigma_*$ lies in $\partial I_H$. Thus the complete coefficient at the first omitted weight is obtained from the boundary alone.
\end{lemma}
\begin{proof}
For $H\le0$ the assertion is immediate, with $\sigma_*=0$. For $H>0$, let $\bnu\notin I_H$ have weight $\sigma$. Along a path from $0$ to $\bnu$ that increases one coordinate at a time, the weights increase strictly, and the first point outside $I_H$ is a boundary point of weight $\le\sigma$; hence $\sigma_*\le\sigma$, so $\sigma_*$ is the least weight outside $[0,H)$. If $\sigma=\sigma_*$ the first point outside $I_H$ has weight exactly $\sigma_*$ by minimality, and no later point on the path can have the same weight, so $\bnu$ is that boundary point.
\end{proof}

\begin{theorem}[Remainder after an exponent cutoff]
\label{thm:remainder}
Let $G^{(r)}_H=z^r\sum_{\bk\in I_H}z^{\bk\cdot\bdelta}C^{(r)}_{\bk}(L)$. Then, as $z\to0^+$,
\begin{equation}
 g(y)^r-G^{(r)}_H(y)=z^{r+\sigma_*}C^{(r)}_{\sigma_*}(L)+\oo\bigl(z^{r+\sigma_*}\bigr)
 =\OO\bigl(z^{r+\sigma_*}\M(z)^{k_*}\bigr).
 \label{eq:remainder}
\end{equation}
If $C^{(r)}_{\sigma_*}\ne0$ the bound is attained by its top logarithm; if $C^{(r)}_{\sigma_*}=0$ (a resonant cancellation) the remainder is $\oo(z^{r+\sigma_*})$, and the true order is that of the next nonzero block, which the theorem does not identify.
\end{theorem}
\begin{proof}
By Theorem~\ref{thm:convergence}, $g^r\in\mathcal E$ with the bound \eqref{eq:coefficient-bound}, and the omitted multi-indices are those of weight $\ge\sigma_*$ by Lemma~\ref{lem:frontier}; Lemma~\ref{lem:tail} (refined form) gives the claim.
\end{proof}

\begin{corollary}[In the variable of the result]
\label{cor:remainder-y}
With $w$ the local variable of $y$ ($w=|y-y_0|$ for finite $y_0$, $w=1/|y|$ for infinite $y_0$),
\[
 g(y)^r-G^{(r)}_H(y)=\OO\bigl(w^{(r+\sigma_*)/|p|}\,\M(w)^{k_*}\bigr),\qquad\frac{r+\sigma_*}{|p|}\ge q ,
\]
since $z=(w/|a|)^{1/p}$ for $p>0$, $z=(|a|w)^{1/|p|}$ for $p<0$, and $\M(z)\asymp\M(w)$.
\end{corollary}

For a pure monomial ($m=0$), the inverse observable is exactly $z^r$: the tail is zero if this term is retained and equals $z^r$ otherwise. The pair $\bigl((r+\sigma_*)/|p|,k_*\bigr)$ specifies the remainder class in the target coordinate, while $z^{r+\sigma_*}C^{(r)}_{\sigma_*}$ is its complete first omitted block. If that coefficient vanishes, a sharper order requires computing further weights or proving that the remaining series terminates.

\begin{example}[Resonances]
For $f=x+x^2+x^3$ the gaps are $1$ and $2$; at weight $2$ the indices $(2,0)$ and $(0,1)$ contribute $2$ and $-1$, at weight $3$ the indices $(3,0)$ and $(1,1)$ contribute $-5$ and $+5$, so
\[
 g(y)=y-y^2+y^3+0\cdot y^4-4y^5+14y^6-30y^7+33y^8+\cdots ,
\]
The coarse boundary estimate at cutoff $4$ is $\OO(y^4)$; summing the weight-$3$ coefficients improves it to $\OO(y^5)$. For $f=x+x^2+2x^3$ the block at $y^3$ cancels instead. The cancellation concerns the sum of all contributions at an exactly equal weight, not each multi-index separately.
\end{example}

\begin{example}[The log example]
Here $m=1$, $\delta=1$, $B=1+L$, so $I_H=\{0,\dots,N\}$ for $H=N+1$, the boundary is $\{N+1\}$, and by \eqref{eq:log-example} the frontier block $y^{N+2}P_{N+1}(\log y)$ has degree exactly $N+1$ with top coefficient $(-1)^{N+1}\Cat_{N+1}$. Hence
\begin{equation}
\begin{aligned}
 g_1(y)&=y+\sum_{n=1}^Ny^{n+1}P_n(\log y)+\OO\bigl(y^{N+2}\M(y)^{N+1}\bigr),\\
 \frac{g_1(y)-y-\sum_{n\le N}y^{n+1}P_n(\log y)}{y^{N+2}(\log y)^{N+1}}&\to(-1)^{N+1}\Cat_{N+1},
\end{aligned}
 \label{eq:log-remainder}
\end{equation}
and in particular after $y-y^2(1+\log y)$ the error is $y^3(2\log^2y+5\log y+3)+\cdots$, not $\OO(y^3)$. The ratio in \eqref{eq:log-remainder} converges slowly (like $1+c/\log y$): for $N=1$ it is $1.89$ at $y=10^{-20}$ and $1.98$ at $y=10^{-100}$.
\end{example}

\subsection{Depth truncation}

When the exponents are symbolic (Section~\ref{sec:boundaries}) their order may be undecidable, but the total degree $|\bk|$ of a multi-index is always known.

\begin{theorem}[Remainder after a depth cutoff]
\label{thm:depth}
Let $G^{(r)}_{\le N}=z^r\sum_{|\bk|\le N}z^{\bk\cdot\bdelta}C^{(r)}_{\bk}(L)$, $\delta_*=\min_i\delta_i$ and $d_*=\max_i\deg B_i$. Then
\begin{equation}
 g(y)^r-G^{(r)}_{\le N}(y)=\OO\bigl(z^{r+(N+1)\delta_*}\M(z)^{(N+1)d_*}\bigr).
 \label{eq:depth-remainder}
\end{equation}
\end{theorem}
\begin{proof}
The omitted terms have $|\bk|\ge N+1$; use \eqref{eq:coefficient-bound} and the count $m^n$ of multi-indices of degree $n$ exactly as in the proof of Lemma~\ref{lem:tail}.
\end{proof}

Depth truncation retains every block of exponent below $r+(N+1)\delta_*$, but it may also retain blocks above that bound while omitting smaller ones when several gaps are present: for $f=x+x^2+x^4$ depth $1$ gives $y-y^2-y^4$ although the omitted depth-two term $2y^3$ dominates $y^4$. The remainder \eqref{eq:depth-remainder} is nevertheless correct as stated.

\section{Transport of remainders and finite smoothness}
\label{sec:transport}

The theory so far concerns an exact finite germ. In practice the forward function is often known only through a truncated expansion, and one wants to know how the omitted part affects the inverse; conversely, one wants to turn a small residual $f(G(y))-y$ of an approximation $G$ into a bound for $G-g$. Both are mean-value arguments, but the scaling factor $z^{1-p}$ of the derivative must not be forgotten.

\subsection{From the forward function to the inverse}

\begin{theorem}[Transport of a forward remainder]
\label{thm:transport}
Let $f$ be the germ \eqref{eq:model} with $p>0$ and let $F=f+R$ on $(0,u_0)$, where
\begin{equation}
 R(u)=\OO\bigl(u^\rho\M(u)^k\bigr),\qquad R'(u)=\OO\bigl(u^{\rho-1}\M(u)^k\bigr),\qquad \rho>p,\ k\in\N .
 \label{eq:R-hypothesis}
\end{equation}
Then $F$ has a local inverse $G$ on the same branch as $g$, and
\begin{equation}
 G(y)-g(y)=\OO\bigl(z^{\rho-p+1}\M(z)^k\bigr)=\OO\bigl(w^{(\rho-p+1)/p}\M(w)^k\bigr),\qquad w=|y-y_0| .
 \label{eq:transport}
\end{equation}
The same conclusion holds if the derivative hypothesis is replaced by the assumption that $F$ is monotone near $0$.
\end{theorem}
\begin{proof}
From \eqref{eq:R-hypothesis} and Lemma~\ref{lem:dominance}, $F(u)-y_0\sim au^p$ and $F'(u)\sim apu^{p-1}$, so $F$ is strictly monotone near $0$ and its inverse $G$ satisfies $G\sim z$ (Theorem~\ref{thm:branch} applies verbatim). Put $u=G(y)$, $v=g(y)$; both are $\sim z$. Then $F(v)-y=F(v)-f(v)=R(v)=\OO(z^\rho\M(z)^k)$, while $F(u)-y=0$. On the interval between $u$ and $v$, which lies in $[z/2,2z]$ for small $z$, $F'$ is $\ge\frac12|ap|z^{p-1}2^{-|p-1|}$ in absolute value, so the mean value theorem gives $|u-v|\le|R(v)|/\min|F'|=\OO(z^{\rho-p+1}\M(z)^k)$. Under monotonicity instead of the derivative bound, use $f$'s derivative: $f(u)-f(v)=y-R(u)-y=-R(u)$, and $f'\asymp z^{p-1}$ on $[z/2,2z]$ gives the same estimate.
\end{proof}

\begin{remark}
The derivative hypothesis is not a consequence of the value bound: $R(u)=u^2\sin(u^{-2})$ is $\OO(u^2)$, but $R'(u)=2u\sin(u^{-2})-2u^{-1}\cos(u^{-2})$ destroys the monotonicity of $u+R(u)$, and $u+R$ has no inverse near $0$. Truncated expansions produced by \wl{Series} of log-exp functions do satisfy \eqref{eq:R-hypothesis}, because termwise differentiation of their asymptotic expansions is valid; the package records the hypothesis as \wl{"InputRemainder" -> \{$\rho$, k\}} and never proves it.
\end{remark}

\begin{corollary}[Cutoff cap]
If the inverse of $F$ is computed from the truncated germ $f$ with an exponent cutoff $q$ in $w$, the result is a valid expansion of $G$ only for $q\le(\rho-p+1)/p$; for such $q$ the remainder class is the coarser of the truncation remainder (Corollary~\ref{cor:remainder-y}) and \eqref{eq:transport}, compared by exponent first and by logarithmic degree at equal exponents.
\end{corollary}

For example $\sin x+x^2\log x=x+x^2\log x-x^3/6+\OO(x^5)$ is a germ with $p=1$, gaps $1,2$ and $\rho=5$, $k=0$; the inverse can be requested up to the cutoff $q=5$ and is $y-y^2\log y+y^3(\tfrac16+\log y+2\log^2y)+\OO(y^4\M^3)$ for $q=4$.

\subsection{From a residual to the error of the inverse}

\begin{theorem}[Residual bracket]
\label{thm:bracket}
Let $F$ be continuous on $[A,B]$ and differentiable inside with $F'\ge\mu>0$. Let $x_0\in(A,B)$, $y\in\R$, $\varrho=F(x_0)-y$. If $[x_0-|\varrho|/\mu,\ x_0+|\varrho|/\mu]\subseteq[A,B]$, then $F(x)=y$ has exactly one solution $x_*$ in $[A,B]$, it lies in that interval, and
\begin{equation}
 \frac{|\varrho|}{\max F'}\ \le\ |x_*-x_0|\ \le\ \frac{|\varrho|}{\mu} .
 \label{eq:bracket}
\end{equation}
\end{theorem}
\begin{proof}
$F(x_0+|\varrho|/\mu)\ge F(x_0)+|\varrho|\ge y$ and $F(x_0-|\varrho|/\mu)\le y$; the intermediate value theorem gives a root in the interval, monotonicity its uniqueness in $[A,B]$, and the mean value theorem both inequalities.
\end{proof}

This gives an a posteriori certificate that needs no expansion theory, only a lower bound on the derivative. For the two examples the global bounds \eqref{eq:global-slopes} give, for every $y>0$ and every $G>0$,
\begin{equation}
 |G-g_1(y)|\le\frac{|G+G^2(1+\log G)-y|}{1-2e^{-5/2}},\qquad |G-g_2(y)|\le|G+G^{\sqrt2}-y| .
 \label{eq:certificates}
\end{equation}
For a general germ the derivative is $\asymp z^{p-1}$, and the bracket must be applied on $[z/2,2z]$ with $\mu=\frac12|ap|z^{p-1}2^{-|p-1|}$ (valid for small $z$), giving $|G-g|\le\frac{2^{1+|p-1|}}{|ap|}z^{1-p}|f(G)-y|$.

\paragraph{The residual of the truncated expansion.} Combining Theorem~\ref{thm:remainder} with the mean value theorem, the truncated inverse $G_H$ satisfies
\begin{equation}
 \frac{f(G_H(y))}{a\,z^p}-1=-p\,z^{\sigma_*}C_{\sigma_*}(L)+\oo(z^{\sigma_*}),
 \label{eq:residual-identity}
\end{equation}
since $f(G_H)-y\approx f'(g)(G_H-g)=-apz^{p-1}\cdot z^{1+\sigma_*}C_{\sigma_*}$. The normalized residual has the same first block as the error, up to the factor $-p$; the package's \wl{InverseResidual} computes it exactly in the jet algebra (with a higher cutoff it displays the first omitted block; below the cutoff it must vanish identically, which is an independent check of every coefficient), and its ``forward residual cutoff'' in the variable of $y$ is $q+(p-1)/p$.

\subsection{Finite smoothness suffices}

The finite germ is analytic away from $0$, and the convergence proof used that. For perturbations outside the finite power--log class (a coefficient $\sin(\log x)$, a term $e^{-1/x}$, a function known only through derivative bounds) the marker formula \eqref{eq:near-identity} can still be applied, and the following theorem, which appears in report~05 with a sketched proof, shows that it is a genuine asymptotic expansion whenever the perturbation and finitely many of its derivatives are small in the power--log scale. We give the complete proof.

\begin{theorem}[Tame near-identity inversion]
\label{thm:tame}
Let $N\ge0$, $\delta>0$ and $M\ge0$. Let $\varphi$ be real-valued and $C^{N+1}$ on $(0,x_0)$ with
\begin{equation}
 \varphi^{(j)}(x)=\OO\bigl(x^{1+\delta-j}\M(x)^M\bigr)\qquad(x\to0^+),\quad j=0,1,\dots,N+1 .
 \label{eq:tame-hypothesis}
\end{equation}
Then $f=x+\varphi$ has a unique local inverse $g$ on $(0,y_0)$ with $g(y)\sim y$, and
\begin{equation}
 g(y)=y+\sum_{n=1}^N\frac{(-1)^n}{n!}\frac{\dd^{n-1}}{\dd y^{n-1}}\bigl[\varphi(y)^n\bigr]+\OO\bigl(y^{1+(N+1)\delta}\M(y)^{(N+1)M}\bigr).
 \label{eq:tame}
\end{equation}
\end{theorem}
\begin{proof}
\emph{Step 1: the homotopy.} For $\lambda\in[0,1]$ put $F_\lambda(x)=x+\lambda\varphi(x)$. By \eqref{eq:tame-hypothesis} with $j=1$, $\varphi'(x)\to0$, so there is $x_1$ with $F_\lambda'\ge\frac12$ on $(0,x_1)$ for all $\lambda\in[0,1]$; each $F_\lambda$ is an increasing bijection of $(0,x_1)$ onto an interval containing $(0,y_1)$ for some $y_1$ independent of $\lambda$. Let $G(\lambda,y)$ be its inverse, so that
\begin{equation}
 G+\lambda\varphi(G)=y,\qquad G(0,y)=y .
 \label{eq:homotopy}
\end{equation}
Since $|\varphi(x)|\le x/2$ for $x<x_1$ (after shrinking $x_1$), $F_\lambda(y/2)<y<F_\lambda(2y)$ and therefore $G(\lambda,y)\in[y/2,2y]$ for all $\lambda$ and all small $y$. The map $(\lambda,x)\mapsto x+\lambda\varphi(x)-y$ is $C^{N+1}$ with $x$-derivative $\ge\frac12$, so by the implicit function theorem $G$ is $C^{N+1}$ in $(\lambda,y)$, with
\begin{equation}
 \partial_yG=\frac1{1+\lambda\varphi'(G)}\in\bigl[\tfrac23,2\bigr],\qquad \partial_\lambda G=-\frac{\varphi(G)}{1+\lambda\varphi'(G)}=\psi(G)\,\partial_yG,\quad\psi:=-\varphi .
 \label{eq:transport-pde}
\end{equation}

\emph{Step 2: the transport identity.} For every $C^{n}$ function $\Psi$ and $1\le n\le N+1$,
\begin{equation}
 \partial_\lambda^n\bigl[\Psi(G)\bigr]=\partial_y^{n-1}\bigl[\psi(G)^n\,\partial_y\Psi(G)\bigr].
 \label{eq:transport-identity}
\end{equation}
For $n=1$: $\partial_\lambda\Psi(G)=\Psi'(G)\partial_\lambda G=\Psi'(G)\psi(G)\partial_yG=\psi(G)\partial_y[\Psi(G)]$. Assume \eqref{eq:transport-identity} for $n$ and all $\Psi$. Then, using it once more for the functions $\psi^n$ and $\Psi$ in the second step,
\begin{align*}
 \partial_\lambda^{n+1}[\Psi(G)]&=\partial_y^{n-1}\,\partial_\lambda\bigl[\psi(G)^n\partial_y\Psi(G)\bigr]
 =\partial_y^{n-1}\bigl[\psi(G)\,\partial_y(\psi(G)^n)\,\partial_y\Psi(G)+\psi(G)^n\,\partial_y\bigl(\psi(G)\partial_y\Psi(G)\bigr)\bigr]\\
 &=\partial_y^{n-1}\,\partial_y\bigl[\psi(G)^{n+1}\partial_y\Psi(G)\bigr],
\end{align*}
where the last equality is the product rule read backwards: $\partial_y[\psi^{n+1}\Psi_y]=(n+1)\psi^n\psi_y\Psi_y+\psi^{n+1}\Psi_{yy}$ and the bracket equals $n\psi^n\psi_y\Psi_y+\psi^n(\psi_y\Psi_y+\psi\Psi_{yy})$, the same thing. (Here $\psi_y$ stands for $\partial_y[\psi(G)]$.)

\emph{Step 3: the coefficients.} At $\lambda=0$, $G=y$, so \eqref{eq:transport-identity} with $\Psi=\mathrm{id}$ gives $\partial_\lambda^nG(0,y)=\partial_y^{n-1}[\psi(y)^n]=(-1)^n\partial_y^{n-1}[\varphi(y)^n]$, the Lagrange terms. Taylor's formula in $\lambda$ with integral remainder gives
\[
 g(y)=G(1,y)=y+\sum_{n=1}^N\frac{(-1)^n}{n!}\partial_y^{n-1}[\varphi^n]+\frac1{N!}\int_0^1(1-\lambda)^N\,\partial_\lambda^{N+1}G(\lambda,y)\,\dd\lambda .
\]

\emph{Step 4: the remainder.} Put $\kappa(y)=y^\delta\M(y)^M$. Since $G\in[y/2,2y]$ and $\M(G)\le2\M(y)$, hypothesis \eqref{eq:tame-hypothesis} gives $\varphi^{(j)}(G)=\OO(\kappa\,y^{1-j})$ uniformly in $\lambda$, for $j\le N+1$. We claim that for $1\le j\le N+1$,
\begin{equation}
 \partial_y^jG=\OO(y^{1-j})\quad(j\ge1),\qquad \partial_y^j[\varphi(G)]=\OO(\kappa\,y^{1-j})\quad(j\ge0),
 \label{eq:derivative-bounds}
\end{equation}
uniformly in $\lambda\in[0,1]$. Differentiating \eqref{eq:homotopy} $j$ times in $y$ and solving for $\partial_y^jG$: $(1+\lambda\varphi'(G))\partial_y^jG$ equals $\delta_{j1}$ minus $\lambda$ times a sum of terms $\varphi^{(m)}(G)\prod\partial_y^{i_s}G$ with $2\le m$, $\sum i_s=j$, $i_s\ge1$ (Fa\`a di Bruno). By induction on $j$ each such term is $\OO(\kappa y^{1-m})\prod\OO(y^{1-i_s})=\OO(\kappa y^{1-j})$, so $\partial_y^jG=\OO(\kappa y^{1-j})\subseteq\OO(y^{1-j})$ for $j\ge2$ (and $\partial_yG=\OO(1)$). The same Fa\`a di Bruno expansion with $m\ge1$ gives the second claim. By Leibniz, $\partial_y^j[\psi(G)^{N+1}]=\OO(\kappa^{N+1}y^{N+1-j})$, and $\partial_y^{N-j}[\partial_yG]=\partial_y^{N-j+1}G=\OO(y^{-(N-j)})$. Applying \eqref{eq:transport-identity} with $\Psi=\mathrm{id}$ and $n=N+1$,
\[
 \partial_\lambda^{N+1}G=\partial_y^N\bigl[\psi(G)^{N+1}\partial_yG\bigr]=\sum_{j=0}^N\binom Nj\,\partial_y^j[\psi(G)^{N+1}]\,\partial_y^{N-j+1}G=\OO\bigl(\kappa^{N+1}y\bigr)
\]
uniformly in $\lambda$, and the integral remainder is $\OO(y\kappa(y)^{N+1})=\OO(y^{1+(N+1)\delta}\M(y)^{(N+1)M})$.
\end{proof}

\begin{remark}
Applied to $\varphi(x)=x^2(1+\log x)$ ($\delta=1$, $M=1$) the theorem re-proves the remainder $\OO(y^{N+2}\M^{N+1})$ of \eqref{eq:log-remainder} with a proof that uses no analyticity at all; applied to $\varphi(x)=x^2\sin(\log x)$ ($\delta=1$, $M=0$) it justifies $g(y)=y-y^2\sin L+y^3\sin L(2\sin L+\cos L)+\OO(y^4)$, $L=\log y$, a perturbation whose coefficients oscillate and which no power--log expansion describes. The package's \wl{PerturbativeInverse} produces the terms; the theorem is the user's tool for asserting the remainder.
\end{remark}

\section{Finite endpoints, infinity, and powers of the inverse}
\label{sec:endpoints}

The branch theorem, coefficient formula and parametric convergence proof only need $p\ne0$ and $u\to0^+$: the substitution $t=u^p$ and the marker theorem are applied at positive points regardless of the endpoint approached by $t$. Statements that use positivity explicitly, such as the exponent-cutoff corollary to the Rouch\'e majorant, retain their stated restrictions. This section makes the endpoint bookkeeping explicit.

\subsection{Local variables}

Let $f$ be given near $x_0$, and let $u\to0^+$ be the local variable of $x$:
\begin{equation}
 u=\begin{cases}x-x_0 & x\to x_0^+,\\ x_0-x & x\to x_0^-,\\ 1/x & x\to+\infty,\\ -1/x & x\to-\infty .\end{cases}
 \label{eq:local-u}
\end{equation}
Substituting into $f$ and expanding (Section~\ref{sec:forward}) gives a germ $y_0+au^p(1+\sum u^{\delta_i}B_i(\log u))$ in the sense of \eqref{eq:model}, where $y_0=\lim f$ is finite when $p>0$. For $p<0$ the target endpoint is $\pm\infty$ with the sign of $a$, while the additive offset $y_0$ in \eqref{eq:model} is set to $0$; any constant term is a correction with gap $-p$. The local variable of $y$ is
\begin{equation}
 w=\begin{cases}\ \ \,y-y_0 & a>0,\ p>0,\\ -(y-y_0) & a<0,\ p>0,\\ \ \ \,1/|y| & p<0 .\end{cases}
 \label{eq:local-w}
\end{equation}
The uniformizer $z=(v/a)^{1/p}$ tends to $0^+$ in all cases. Its relation to $w$ is $w=|a|z^p$ for $p>0$ and $w=z^{|p|}/|a|$ for $p<0$.

\subsection{The theorem at an arbitrary endpoint}

\begin{theorem}[Signed leading power]
\label{thm:endpoints}
Theorems~\ref{thm:branch}, \ref{thm:coefficients}, \ref{thm:convergence} and \ref{thm:remainder} hold verbatim for the germ \eqref{eq:model} with $p<0$: the unique branch with $u\to0^+$ is $u=g(y)$ with $g\sim z$, $g^r$ has the expansion \eqref{eq:master} convergent for small $z$ (i.e.\ for $|y|$ large), and the remainder after the exponent cutoff $q$ in $w=1/|y|$ is $\OO(w^{(r+\sigma_*)/|p|}\M(w)^{k_*})$.
\end{theorem}
\begin{proof}
The derivative $f'(u)=apu^{p-1}(1+\oo(1))$ has a constant sign, so $f$ is monotone on $(0,\varepsilon)$ with $f\to\pm\infty$; the inverse exists on a neighbourhood of $\pm\infty$ and $u\sim z$. The proof of Theorem~\ref{thm:coefficients} is a computation at a fixed positive $q=v/a$ and never uses the sign of $p$. The parametric equation \eqref{eq:parametric} has $\partial_w\mathcal E(1,0)=p\ne0$ regardless of sign, and the specialisation $T_{ij}=z^{\delta_i}L^j\to0$ as $z\to0^+$ is the same. The remainder theorem is Lemma~\ref{lem:tail}.
\end{proof}

In the variable $y$ the blocks of $g^r$ read $(v/a)^{(r+\sigma)/p}C^{(r)}_\sigma\bigl(\tfrac1p\log(v/a)\bigr)$, with negative exponents of $y$ when $p<0$; the exponent of $w$ is $(r+\sigma)/|p|$, which is the quantity the cutoff refers to.

\subsection{Powers of the inverse and the original variable}

The observable power $r$ in \eqref{eq:master} gives directly the expansion of $u^r$. In terms of $x$:
\begin{equation}
\begin{aligned}
 (x-x_0)^r&=g(y)^r\ (x\to x_0^+),& (x_0-x)^r&=g(y)^r\ (x\to x_0^-),\\
 x^r&=g(y)^{-r}\ (x\to+\infty),& x^r&=(-1)^rg(y)^{-r}\ (x\to-\infty),
\end{aligned}
 \label{eq:power-back}
\end{equation}
The last equality in \eqref{eq:power-back}, involving $(-1)^r$, is used here only for integer $r$. Thus $x^r$ at infinity is obtained from \eqref{eq:master} with local-coordinate power $-r$. Both signed distances $x-x_0$ from above and $x_0-x$ from below are positive and admit arbitrary real powers. For $r=1$ and a finite endpoint, the original inverse is $x=x_0\pm g(y)$. A non-integer power of a negative original variable needs an additional branch convention and is outside the positive-real-power convention of this article.

\subsection{Examples}

\begin{example}[The Wright omega function]
$f(x)=x+\log x$ as $x\to+\infty$: $u=1/x$, $f=u^{-1}-\log u=u^{-1}(1-u\log u)$, so $a=1$, $p=-1$, one gap $\delta=1$ with $B=-L$, and target endpoint $+\infty$. The uniformizer is $z=1/y$, $L=-\log y$, and \eqref{eq:master} with $r=-1$ gives the original variable $x=1/g(y)$:
\begin{equation}
 x=y-\log y+\frac{\log y}{y}+\frac{\tfrac12\log^2y-\log y}{y^2}+\frac{\tfrac13\log^3y-\tfrac32\log^2y+\log y}{y^3}+\OO\Bigl(\frac{\log^4y}{y^4}\Bigr),
 \label{eq:omega}
\end{equation}
the classical expansion of the Wright omega function $\omega(y)=\W(e^y)$ \cite{cghjk,proveit}; the blocks are the Lambert polynomials $P_n(\log y)$ of \cite{proveit}.
\end{example}

\begin{example}[Reciprocal corrections]
$f(x)=x+1/x$ at $+\infty$: $u=1/x$, $f=u^{-1}(1+u^2)$, $p=-1$, gap $2$, constant coefficient. The inverse is the larger root of $x^2-yx+1=0$, $x=\bigl(y+\sqrt{y^2-4}\bigr)/2=y-y^{-1}-y^{-3}-2y^{-5}-\cdots$ (the coefficients are Catalan numbers), which \eqref{eq:constant-coefficients} reproduces.
\end{example}

\begin{example}[Half-integer gaps]
$f(x)=x+\sqrt x$ at $+\infty$: $f=u^{-1}(1+u^{1/2})$, gap $\frac12$. The inverse is $x=\bigl(\sqrt{4y+1}-1\bigr)^2/4=y-\sqrt y+\frac12-\frac1{8\sqrt y}+\frac1{128y^{3/2}}+\OO(y^{-5/2})$; the constant term $\frac12$ is the block of weight $2\cdot\frac12$, i.e.\ exponent $0$ in $w=1/y$, despite the infinite target endpoint.
\end{example}

\begin{example}[A finite endpoint from below]
$f(x)=x-x^2$ near $x_0=0$ with $x\to0^-$: $u=-x$, $f=-u-u^2$, $a=-1$, $p=1$, $y_0=0$, and $w=-y$ (the image is $y<0$). The inverse is $u=-y-y^2-2y^3-\cdots$, i.e.\ $x=y+y^2+2y^3+\OO(y^4)$, the branch of $x=(1-\sqrt{1-4y})/2$ continued to negative $y$.
\end{example}

\begin{example}[A pole]
$f(x)=1+1/x$ near $x_0=0^+$: $f=u^{-1}(1+u)$, $p=-1$, target endpoint $+\infty$, and $x=1/(y-1)=y^{-1}+y^{-2}+\cdots$ with remainder $\OO(y^{-3})$ after two terms.
\end{example}

\begin{example}[The second question, inverted]
For $F(x)=(1+x+x^{\sqrt2})^{\sqrt2}$ at $+\infty$ the forward germ is $u^{-2}(1+\sqrt2u^{\sqrt2-1}+\sqrt2u^{\sqrt2}+\cdots)$ from \eqref{eq:answer3}, with gaps generated by $\sqrt2-1$ and $\sqrt2$, $p=-2$, target endpoint $+\infty$, and $z=y^{-1/2}$. The inverse is
\begin{equation}
\begin{aligned}
 x={}&\sqrt y-\frac{1}{\sqrt2}\,y^{(2-\sqrt2)/2}+\Bigl(\frac34-\frac1{2\sqrt2}\Bigr)y^{(3-2\sqrt2)/2}\\
 &+\Bigl(1-\frac{5}{3\sqrt2}\Bigr)y^{(4-3\sqrt2)/2}-\frac{1}{\sqrt2}\,y^{(1-\sqrt2)/2}+\OO\bigl(y^{(5-4\sqrt2)/2}\bigr),
\end{aligned}
 \label{eq:second-inverse}
\end{equation}
the terms being ordered by decreasing exponent $-(r+\sigma)/2$ of $y$, with local-coordinate power $r=-1$. The forward remainder is transported by \eqref{eq:transport-power} with that value of $r$; the available forward order must be sufficient for the desired inverse cutoff. An independent derivation removes the outer power exactly: $1+x+x^{\sqrt2}=y^{1/\sqrt2}$, then applies the finite-germ coefficient formula.
\end{example}

\section{Worked examples}
\label{sec:examples}

All expansions in this section were produced by the package and checked by the residual composition of Section~\ref{sec:transport} and numerically; the displayed logarithms are $\log$ of the indicated argument.

\subsection{The two examples of the first question, once more}

For $f_1=x+x^2(1+\log x)$ the package call \wl{AsymptoticInverse[x + x\^{}2 (1 + Log[x]), \{x, 0\}, \{y, 6\}]} returns \eqref{eq:answer1} with the remainder \wl{PowerLogRemainder[y, 6, 5]}; the sixth block is
\begin{align*}
 P_5(L)&=-42L^5-\tfrac{3727}{12}L^4-\tfrac{5365}{6}L^3-1256L^2-\tfrac{5177}{6}L-\tfrac{2789}{12}\\
 &=-\tfrac1{12}(1+L)\bigl(504L^4+3223L^3+7507L^2+7565L+2789\bigr),
\end{align*}
and the seventh $P_6=\tfrac1{60}(1+L)(7920L^5+63852L^4+200833L^3+308537L^2+231883L+68307)$. With rational exponents the same object is available as \wl{SeriesData[y, 0, \{1, -1 - Log[y], 3 + 5 Log[y] + 2 Log[y]\^{}2, ...\}, 1, 6, 1]}, on which the built-in series arithmetic works (with the caveat that its \wl{O[y]\^{}6} hides the factor $\M(y)^5$).

For $f_2=x+x^{\sqrt2}$, \wl{SeriesTermGoal -> 4} returns exactly the four terms requested in the question with remainder \wl{PowerLogRemainder[y, 4 Sqrt[2] - 3, 0]}; the exponent cutoff that produces the same result is $4\sqrt2-3$.

\subsection{Nonunit leading power}

$f(x)=3x^2\bigl(1+x(1+\log x)\bigr)$: $a=3$, $p=2$, $\delta=1$, $B=1+L$. With $z=\sqrt{y/3}$ and $L=\log z=\tfrac12\log(y/3)$, \eqref{eq:single-block} gives
\begin{equation}
 g(y)=z-\frac{z^2}{2}(1+L)+\frac{z^3}{8}\bigl(5L^2+12L+7\bigr)-\frac{z^4}{48}(\D+6)(\D+8)(1+L)^3+\OO\bigl(z^5\M^4\bigr),
 \label{eq:nonunit}
\end{equation}
where $\tfrac1{48}(\D+6)(\D+8)(1+L)^3=(1+L)^3+\tfrac78(1+L)^2+\tfrac18(1+L)$. Two ingredients are easy to forget: the factor $p^{-n}=2^{-n}$ and the logarithm $L=\tfrac12\log(y/3)$, not $\log y$. In the variable $y$ the expansion reads
\begin{align*}
 g(y)={}&\frac{\sqrt y}{\sqrt3}-\frac y3\Bigl(\frac12+\frac{\log(y/3)}{4}\Bigr)+\frac{y^{3/2}}{3\sqrt3}\Bigl(\frac78+\frac34\log\frac y3+\frac5{32}\log^2\frac y3\Bigr)\\
 &-\frac{y^2}{9}\Bigl(2+\frac{39}{16}\log\frac y3+\frac{31}{32}\log^2\frac y3+\frac18\log^3\frac y3\Bigr)+\OO\bigl(y^{5/2}\M^4\bigr).
\end{align*}

\subsection{Several gaps}

\paragraph{Irrational and logarithmic gaps together.} $f(x)=x+x^{\sqrt2}+x^2(1+\log x)$ has gaps $\delta_1=\sqrt2-1$ and $\delta_2=1$. Below the exponent $3$ the support is $\{1,\sqrt2,2\sqrt2-1,2,3\sqrt2-2,1+\sqrt2,4\sqrt2-3,2\sqrt2\}$ and
\begin{align*}
 g(y)={}&y-y^{\sqrt2}+\sqrt2\,y^{2\sqrt2-1}-y^2(1+L)-\frac{6-\sqrt2}{2}\,y^{3\sqrt2-2}+y^{1+\sqrt2}\bigl((\sqrt2+2)L+\sqrt2+3\bigr)\\
 &+\Bigl(\frac{17\sqrt2}{3}-4\Bigr)y^{4\sqrt2-3}-y^{2\sqrt2}\Bigl((3\sqrt2+5)L+5\sqrt2+\frac{13}{2}\Bigr)+\OO\bigl(y^3\M^2\bigr),\qquad L=\log y ,
\end{align*}
the mixed blocks coming from \eqref{eq:cross}: the index $(1,1)$ gives $(\D+2+\sqrt2)(1+L)=(\sqrt2+2)(1+L)+1$. The remainder is the block $y^3(2L^2+5L+3)$ of the index $(0,2)$, of degree $2$.

\paragraph{Two irrational gaps.} $f(x)=x+x^{\sqrt2}(1+\log x)+2x^{\sqrt3}$ has gaps $\sqrt2-1$ and $\sqrt3-1$, whose sums are the numbers $k_1(\sqrt2-1)+k_2(\sqrt3-1)$; the package represents them by \wl{RootReduce} (for instance $\sqrt2+\sqrt3=\wl{Root[1 - 10 \#\^{}2 + \#\^{}4 \&, 4]}$) and displays them with \wl{ToRadicals}. Below the exponent $3$ there are eleven blocks; the residual check passes only when the coefficients, which involve $\sqrt{5+2\sqrt6}=\sqrt2+\sqrt3$, are canonicalised with \wl{RootReduce}: this is precisely the zero-recognition defect that made one test of report~01 fail (Section~\ref{sec:reports}).

\subsection{Automatically expanded forward functions}

$\sin x+x^2\log x$ is not a finite power--log sum. The evaluator of Section~\ref{sec:evaluator} expands it to $x+x^2\log x-\tfrac16x^3+\tfrac1{120}x^5+\OO(x^7)$, records $(\rho,k)=(7,0)$, and the inverse below the exponent $4$ is
\[
 g(y)=y-y^2\log y+y^3\Bigl(\frac16+\log y+2\log^2y\Bigr)+\OO\bigl(y^4\M^3\bigr);
\]
here the $y^3$ block combines $(\D+4)L^2/2=2L^2+L$ from the index $(2,0)$ (gap $1$, $B_1=L$) with $+\tfrac16$ from the index $(0,1)$ (gap $2$, $B_2=-\tfrac16$). Similarly $\tan x$ gives $\arctan$: \wl{SeriesTermGoal -> 4} yields $y-\tfrac13y^3+\tfrac15y^5-\tfrac17y^7$, the forward order being raised automatically until the transported remainder allows the fourth term.

\subsection{Symbolic data}

With symbolic coefficients the exponents remain numeric and everything goes through under assumptions that make the leading coefficient's sign and the reality of the coefficients provable: $f=ax+bx^2$ with $a>0$, $b\in\R$ gives $g=y/a-by^2/a^3+2b^2y^3/a^5+\OO(y^4)$. With a symbolic exponent the ordering of the support is undecidable in general and depth truncation (Theorem~\ref{thm:depth}) is used: $f=x+x^\alpha$ with $\alpha>1$ gives, at depth $3$,
\[
 g=y-y^\alpha+\alpha y^{2\alpha-1}-\frac{\alpha(3\alpha-1)}{2}y^{3\alpha-2}+\OO\bigl(y^{4\alpha-3}\bigr),
\]
which is \eqref{eq:fuss-catalan} with a symbolic $\alpha$.

\subsection{Powers and logarithms of the inverse}

For $f=x+x^2$, \eqref{eq:master} with $r=2$ gives $g^2=y^2-2y^3+5y^4-14y^5+\cdots$ (Catalan numbers again), with $r=-1$ it gives $1/g=y^{-1}+1-y+2y^2-5y^3+\cdots$, and \eqref{eq:master-1-log} gives $\log g=\log y-y+\tfrac32y^2-\tfrac{10}3y^3+\cdots$. These expansions are exact consequences of the same coefficient formula and cost nothing extra; expanding \wl{Normal[g]\^{}2} instead would silently lose the remainder information.

\section{Lambert coordinates and boundaries of the theory}
\label{sec:boundaries}

The convergent power--log class is closed under the operations of Theorem~\ref{thm:closure} and inversion of germs with a monomial leading block. A leading logarithm or exponential changes the natural coordinate. For exact Lambert cores, that coordinate change reduces the problem to the same polynomial--logarithmic arithmetic. More general perturbations require additional bookkeeping, described below.

\subsection{A leading logarithm: Lambert \texorpdfstring{$W$}{W} and the \texorpdfstring{$\log\log$}{log-log} scale}

Let $a\ne0$, $p\ne0$, $b\ne0$, and $x\to0^+$ with $x<1$. The exact core $ax^p(-\log x)^b=y$ has $Y=y/a>0$. With $\tau=-\log x\to+\infty$, it reads $\tau^be^{-p\tau}=Y$, hence
\begin{equation}
 x=\exp\Bigl[\frac bp\,\W_k\Bigl(-\frac pb\,Y^{1/b}\Bigr)\Bigr],\qquad
 k=\begin{cases}-1,&p/b>0,\\0,&p/b<0.\end{cases}
 \label{eq:lambert-core}
\end{equation}
The branch is forced by $\tau\to\infty$: when $p/b>0$ the argument approaches $0$ from below and $\W_{-1}\to-\infty$; when $p/b<0$ it tends to $+\infty$ and $\W_0\to+\infty$. For $p,b>0$, the other real root on $\W_0$ approaches $x=1$ as $Y\to0$, while the two branches meet at the maximum $x=e^{-b/p}$. For instance the inverse of $x\log(1/x)$ near $0^+$ is $-y/\W_{-1}(-y)$, and that of $x\log x$ on the same source branch is $y/\W_{-1}(y)$ for $y\to0^-$. For $p>0$, with $\Lambda=\log(a/y)\to\infty$ and $\Lambda_2=\log\Lambda$, the expansion begins
\begin{equation}
 x=\Bigl(\frac ya\Bigr)^{1/p}\Bigl(\frac\Lambda p\Bigr)^{-b/p}\Bigl[1-\frac{b^2}{p\,\Lambda}\log\frac\Lambda p+\OO\Bigl(\frac{\Lambda_2^2}{\Lambda^2}\Bigr)\Bigr],
 \label{eq:loglog}
\end{equation}
a series in $1/\Lambda$ with polynomial coefficients in $\Lambda_2$, multiplying the leading factor. The following construction computes all its logarithmic orders, with convergence and a remainder bound.

\subsection{A convergent universal Lambert expansion}
\label{sec:lambert-expansion}

Let $Z$ approach $+\infty$ on the principal real branch or $0^-$ on the lower real branch. Set
\[
 s=\begin{cases}1,&\W_0(Z),\ Z\to+\infty,\\-1,&\W_{-1}(Z),\ Z\to0^-,\end{cases}
 \qquad A=|\log|Z||,\quad t=A^{-1},\quad\ell=\log t.
\]
Then $\log|\W|+\W=\log|Z|=sA$. Writing $\W=s(1+V)/t$ gives the normalized equation
\begin{equation}
 V+s t\log(1+V)-s t\ell=0.
 \label{eq:lambert-normalized}
\end{equation}

\begin{theorem}[Lambert expansion in a logarithmic coordinate]
\label{thm:lambert}
Equation \eqref{eq:lambert-normalized} has a unique solution $V\to0$ of the form $\sum_{n\ge1}t^nP_n(\ell)$ with $P_1=s\ell$ and $\deg P_n\le n-1$ for $n\ge2$. This series converges for sufficiently small $t>0$. For $N\ge1$, truncating $V$ after $t^N$ leaves $\OO(t^{N+1}\M(t)^N)$, and multiplying by $s/t$ gives the corresponding remainder for $\W$. In particular,
\begin{equation}
 \W=\frac{s}{t}+\ell-s\ell t+
 \Bigl(\ell+\frac{\ell^2}{2}\Bigr)t^2
 -s\Bigl(\ell+\frac{3\ell^2}{2}+\frac{\ell^3}{3}\Bigr)t^3
 +\OO(t^4\M(t)^4).
 \label{eq:lambert-universal}
\end{equation}
\end{theorem}
\begin{proof}
Replace $t\ell$ by an independent parameter $T$. The function $V+s t\log(1+V)-sT$ is analytic near $(V,t,T)=(0,0,0)$ and its $V$-derivative there is $1$. The analytic implicit function theorem gives a convergent Taylor series $V(t,T)$. Moreover $V=sT+t H(t,T)$ for an analytic $H$ with $H(0,0)=0$. After $T=t\log t$, every term of total degree $n\ge2$ therefore has logarithmic degree at most $n-1$, and $P_1=s\ell$. Cauchy estimates bound the coefficient norms exponentially. The tail starting at degree $N+1$ is bounded by a geometric series with first term $t^{N+1}\M(t)^N$, proving the remainder. Coefficients can be computed by the iteration $V\gets s t\ell-s t\log(1+V)$, which gains one power of $t$ each time, or by Newton iteration. The real solution has $\W\sim sA$ and hence belongs to the specified real Lambert branch. Expanding the equation gives \eqref{eq:lambert-universal}.
\end{proof}

An explicit sufficient convergence condition is
\begin{equation}
 q_L=4t(1+|\log t|)<1.
 \label{eq:lambert-convergence-condition}
\end{equation}
Indeed, for complex $|t|,|T|\le1/4$, the map $V\mapsto sT-s t\log(1+V)$ sends $|V|\le1/2$ into itself, since $(1+\log2)/4<1/2$, and has derivative bounded by $1/2$. Its fixed point is jointly analytic and bounded by $1/2$, so its Taylor coefficients are bounded by $4^{a+b}/2$. Beyond the leading term $sT$, every coefficient has $a\ge1$ because $V(0,T)=sT$. If $V_N$ retains total degrees through $N\ge1$, the geometric majorant therefore gives
\begin{equation}
 |V-V_N|\le\frac{2t\,q_L^N}{1-q_L},\qquad
 \left|\W-\frac{s}{t}(1+V_N)\right|
 \le\frac{2q_L^N}{1-q_L}.
 \label{eq:lambert-explicit-tail}
\end{equation}
The construction follows the two-small-parameter method in the supplied transseries notes \cite{proveit}; it establishes a sufficient domain and does not claim convergence from a fixed-logarithm implicit-function argument alone.

For the leading-log core, put $Z=-(p/b)Y^{1/b}$ and choose $s$ from its branch. Since $-bs/p>0$, an arbitrary real observable power is best assembled as
\begin{equation}
 x^r=Y^{r/p}\Bigl(-\frac{bs}{p}\Bigr)^{-br/p}
 t^{br/p}(1+V)^{-br/p}.
 \label{eq:lambert-log-observable}
\end{equation}
This uses a positive real power of the unit $1+V$. Directly exponentiating a truncated $\W$ would shift the required precision by one order; \eqref{eq:lambert-log-observable} makes that precision shift explicit.

\begin{example}[The inverse of $x\log x$]
For $y\to0^-$ let $A=\log(1/(-y))$ and $B=\log A$. The branch with $x\to0^+$ is
\[
 x=\frac{-y}{A}\left[1-\frac{B}{A}
 +\frac{B^2-B}{A^2}
 +\frac{-B^3+\tfrac52B^2-B}{A^3}
 +\OO\left(\frac{(1+B)^4}{A^4}\right)\right].
\]
The endpoint $x\to1$ selects a different inverse germ and is not represented by this expansion.
\end{example}

\subsection{Exponential cores at infinity}

For a positive source coordinate $x$ with $X=x^p\to+\infty$, consider $y=a x^b e^{c x^p}$, where $a\ne0$, $p\ne0$, $c\ne0$, and $Y=y/a>0$. If $b\ne0$, taking the power $p/b$ gives
\begin{equation}
 X=\frac{b}{cp}\,\W_k\left(\frac{cp}{b}Y^{p/b}\right),\qquad
 k=\begin{cases}0,&cp/b>0,\\-1,&cp/b<0.\end{cases}
 \label{eq:exponential-core}
\end{equation}
These branch choices ensure $X>0$ and $X\to\infty$: the argument tends to $+\infty$ in the first case and $0^-$ in the second. If $b=0$, inversion is elementary, $x=(\log Y/c)^{1/p}$. In the Lambert cases, with $s,t,V$ as above,
\[
 x^r=\left(\frac{bs}{cp}\right)^{r/p}t^{-r/p}(1+V)^{r/p},
\]
whose positive prefactor and unit series again have explicit precision tracking.

\begin{example}[The inverse of $xe^x$ at infinity]
Here $x=\W_0(y)$. With $A=\log y$ and $B=\log A$,
\[
 x=A-B+\frac{B}{A}+\frac{B(B-2)}{2A^2}
 +\frac{B(2B^2-9B+6)}{6A^3}
 +\OO\left(\frac{(1+B)^4}{A^4}\right).
\]
Both this expansion and the $x\log x$ expansion use $t=1/A$ as the local variable. An exponent cutoff in $t$ specifies logarithmic accuracy; it is not a cutoff in $y$ or $1/y$. A nonconstant external factor such as $-y$ must also multiply the remainder descriptor.
\end{example}

\subsection{Perturbations of a Lambert core}

Relative perturbations $F(x)=ax^p(-\log x)^b(1+\sum_i x^{\delta_i}B_i(\log x))$ can be approached with the general-core formula \eqref{eq:general-core}. If $\phi$ is the exact core inverse and $\tau=-\log\phi(y)$, then
\[
 \phi'(y)=\left[a x^{p-1}\tau^{b-1}(p\tau-b)\right]^{-1}_{x=\phi(y)}.
\]
The marker coefficients have powers of $\phi(y)$ and rational functions of $\tau$, with denominators arising from $p\tau-b$. A convergent marker expansion requires analytic control on the chosen branch, and re-expanding its coefficients introduces a second order of truncation: perturbation power and logarithmic depth. The exact-core Lambert construction proves the logarithmic expansion. Sections~\ref{sec:core-perturbations} and~\ref{sec:exponential-core-perturbation} retain the exact core and compute separate perturbation corrections for the admitted finite families.

The leading logarithmic polynomial itself can also be generalized before introducing a second perturbation scale.

\begin{corollary}[An arbitrary polynomial at the leading exponent]
\label{cor:polynomial-log-core}
Let $F_0(u)=a u^p Q(\log u)$, where $p\ne0$ and $Q$ is a nonconstant real polynomial of degree $q\ge1$, on the real branch $u\to0^+$. Its inverse has a convergent relative power--log expansion in a reciprocal-logarithm coordinate, even when $Q$ is not a power of an affine polynomial.
\end{corollary}
\begin{proof}
Absorb the leading coefficient of $Q$ into $a$, so $Q$ is monic, and set $b=[L^{q-1}]Q/q$ and $\tau=-\log u-b$. Then
\[
 Q(-\tau-b)=(-\tau)^q R(1/\tau),\qquad
 R(v)=1+\sum_{j=2}^q r_jv^j.
\]
With $Y=y/(a(-1)^q)>0$ (after subtracting any constant target offset), $k=-p/q$, $s=\operatorname{sign}k$ and $Z=kY^{1/q}e^{pb/q}$, take $A=s\log|Z|$, $t=1/A$, and $\tau=(1+V)/(|k|t)$. The inverse equation becomes
\begin{equation}
 V+s t\left[-\ell+\log(1+V)
 +\frac1q\log R\left(\frac{|k|t}{1+V}\right)\right]=0.
 \label{eq:polynomial-log-normalized}
\end{equation}
Replacing $t\ell$ by an independent parameter again gives an analytic equation with $V$-derivative $1$ at the origin. Its solution is jointly analytic, and $R(0)=1$ makes its real powers analytic too. For any real local-coordinate observable power $r$,
\[
 u^r=Y^{r/p}\left(\frac{A}{|k|}\right)^{-qr/p}
 (1+V)^{-qr/p}
 R\left(\frac{|k|t}{1+V}\right)^{-r/p}.
\]
This yields convergence and the same type of finite-order logarithmic remainder as before. The sufficient numerical bound \eqref{eq:lambert-explicit-tail} was proved for $R=1$ and is not asserted unchanged for a general $R$.
\end{proof}

The polynomial correction modifies the logarithmic coefficients and must be included in the equation; it is not beyond all logarithmic orders. In general this inverse has no claimed exact \wl{ProductLog} expression. The construction also applies to logarithmic cores at either infinite source endpoint after $u=1/x$ or $u=-1/x$, using the observable and sign rules of Section~\ref{sec:endpoints}.

\begin{proposition}[Perturbations beyond every logarithmic order]
\label{prop:lambert-flat-perturbation}
Suppose a finite leading-log core $F_0(u)=a u^p Q(\log u)$ from Corollary~\ref{cor:polynomial-log-core}, or an affine-log power core, is supplemented by finitely many polynomial--logarithmic terms of strictly larger $u$-exponent. Alternatively, suppose a growing exponential core $F_0(x)=a x^b e^{c x^p}$ with $c,p>0$ at $x\to+\infty$ is supplemented by a finite power--log sum. Let $g_0$ and $g$ be the corresponding inverse branches, with any common target offset subtracted, and let $t$ be the Lambert coordinate of $g_0$. For every fixed real $r$ and every $N>0$,
\[
 g(y)^r-g_0(y)^r=\oo\bigl(g_0(y)^r t^N\bigr).
\]
Consequently every finite relative logarithmic expansion of the core inverse is also an expansion of the full inverse, with the same finite-order remainder class.
\end{proposition}
\begin{proof}
In the leading-log case $F_0'(u)\sim pF_0(u)/u$, and the perturbation and its derivative are smaller by a factor $\OO(u^\delta\M(u)^D)$ for some $\delta>0$ and finite $D$. A mean-value bracket near $g_0$ yields $(g-g_0)/g_0=\OO(g_0^\delta\M(g_0)^D)$. Here $1/t$ is a positive constant times $|\log g_0|$ plus lower logarithmic terms, so this bound is $\oo(t^N)$ for every $N$.

In the exponential case $F_0'/F_0\sim cp x^{p-1}$, whereas the additive perturbation $R$ and its derivative have only power--log growth. A bracket of width $\OO(|R(g_0)|/|F_0'(g_0)|)$ has slope comparable to $F_0'(g_0)$: the logarithm of that derivative changes by $\OO(R/F_0)=\oo(1)$ across the bracket. Thus $(g-g_0)/g_0$ is exponentially small in $g_0^p$, while $t$ is comparable to $g_0^{-p}$. Finally apply the mean value theorem to the positive real power $u^r$. Continuous local inverse branches exist because the perturbation derivative is smaller than the nonzero core derivative.
\end{proof}

The package marks this situation with \wl{"LeadingCoreOnly" -> True}, stores the discarded expression in \wl{"BeyondLogarithmicOrders"}, and leaves \wl{"ExactInverseExpression"} unavailable. Those terms affect an exponentially smaller scale even though they do not change any finite coefficient in the reported logarithmic expansion.

\subsection{Zero gaps: corrections of relative size \texorpdfstring{$1/\log$}{1/log}}

For $f(x)=x+x/\log x$ the correction has power gap $0$ and relative size $1/\log x$. All corrections of the inverse then sit at the same power $y$ with increasing negative powers of $\log y$, the index set $I_H$ is infinite, and the exponent cutoff has no meaning. The right coordinate is $\tau=1/\log y$: with $g=yW$ the equation becomes $W\bigl(1+\tau/(1+\tau\log W)\bigr)=1$, analytic near $(W,\tau)=(1,0)$, so $W$ is an ordinary convergent power series in $\tau$, $W=1-\tau+\tau^2-2\tau^3+\tfrac72\tau^4-\cdots$, i.e.\ $g(y)=y\bigl(1-1/\log y+1/\log^2y-2/\log^3y+\cdots\bigr)$. The general-core formula produces these terms, but the scale is the reciprocal-logarithm scale, not the power--log scale, and Section~\ref{sec:logarithmic-scales} implements the required reciprocal-logarithm valuation.

\subsection{Flat terms}

$e^{-1/x}$ is smaller than every power of $x$; a germ $x+e^{-1/x}$ has the same asymptotic power--log expansion as $x$. Its inverse obeys $g=y-e^{-1/g}$, so $y-e^{-1/y}\le g<y$ and $(y-g)/y^2\to0$. Consequently $e^{-1/g}/e^{-1/y}\to1$, proving $g-y=-e^{-1/y}(1+\oo(1))$. The marker formula begins $y-e^{-1/y}+y^{-2}e^{-2/y}+(y^{-3}-\tfrac32y^{-4})e^{-3/y}+\cdots$. These exponential sectors are invisible to a pure power--log expansion. Section~\ref{sec:flat-sectors} computes them with a separate sector cutoff and a complete-tail bound; this differs from the exact exponential cores above, whose inverses reduce to Lambert coordinates.

\subsection{Oscillatory coefficients}

$x^2\sin(\log x)$ is $\OO(x^2)$ with all derivatives of the right size, so Theorem~\ref{thm:tame} applies and the marker formula gives a valid expansion $y-y^2\sin L+y^3\sin L(2\sin L+\cos L)+\OO(y^4)$ whose ``coefficients'' are bounded oscillating functions of $L=\log y$. They are not polynomials in $L$, the coefficient algebra is different, and the sharp remainder theory of Section~\ref{sec:remainders} (which uses degrees) does not apply; only the tame theorem does. Section~\ref{sec:fourier-coefficients} supplies a finite Fourier coefficient algebra and a nonoscillatory remainder bound. The lower-level \wl{PerturbativeInverse} remains a formula generator.

\subsection{Symbolic exponents and nonuniformity}

When a gap $\delta$ is a symbol, the order of the weights $k\cdot\bdelta$ depends on the parameter: the order of $2\delta_1$ and $\delta_2$ changes at $\delta_2=2\delta_1$. Exponent truncation needs provable comparisons; depth truncation (Theorem~\ref{thm:depth}) only needs positive gaps after a valid leading term and real branch have been selected. Nothing is uniform as a gap tends to $0$: the $y$-radius $R_\alpha^{1/(\alpha-1)}$ of the $x+x^\alpha$ inverse tends to $0$ as $\alpha\downarrow1$, although $R_\alpha\to1$, and at $\alpha=1$ the inverse is $y/2$.

\subsection{Nested powers of the logarithm}

Non-integer powers $(-\log x)^s$, $\log(-\log x)$, and negative powers of $\log x$ in a correction leave the polynomial coefficient algebra. For example $\varphi(x)=x^{1+\delta}(-\log x)^{1/2}$ is real for small positive $x$ and satisfies Theorem~\ref{thm:tame} with $M=1$, but repeated differentiation produces negative half-integer as well as positive half-integer powers of $-L$. A suitable coefficient algebra therefore contains $(-L)^{1/2}$ and its inverse. The Euler identity and composition recurrence still apply when the coefficient functions are differentiable or analytic as needed; closure and remainder bounds in the enlarged scale must also be established.

The finite hierarchy of Section~\ref{sec:logarithmic-scales} establishes these closure and remainder contracts for generalized Laurent monomials in finitely many positive logarithmic levels. More general coefficient functions still need an independently justified algebra and ordering.

\section{Inversion through explicit coordinate transformations}
\label{sec:coordinates}

The leading expression in the original variable need not belong to the scale used by the inverse engine. A change of source or target coordinate can put the equation into a supported scale while preserving the requested real branch. The transformation is part of the mathematical result: it determines the expansion variable, how the cutoff is interpreted, and how an inner remainder becomes an error in the original variable.

\subsection{Target transformations}

\begin{theorem}[An exact change of target coordinate]
\label{thm:target-coordinate}
Let $F$ be continuous and strictly monotone on a source interval $I$, and let $J$ be continuous and one-to-one on $F(I)$. Put $\Phi=J\circ F$. If $h$ is the inverse of $\Phi$ on $J(F(I))$, then the inverse of $F$ on $F(I)$ is
\[
 g(y)=h(J(y)).
\]
If an observable of that inverse has an expansion $h(v)^r=A(v)+\OO(R(v))$ as $v$ approaches the transformed target endpoint, then
\[
 g(y)^r=A(J(y))+\OO(R(J(y)))
\]
on the corresponding source branch, provided the same real power is defined there.
\end{theorem}
\begin{proof}
$\Phi(h(J(y)))=J(y)$ gives $J(F(h(J(y))))=J(y)$. Injectivity of $J$ gives $F(h(J(y)))=y$, and the source interval fixes the inverse branch. The remainder statement is direct substitution along the transformed endpoint; no output reconstruction is involved.
\end{proof}

The theorem distinguishes two operations that are sometimes conflated. Changing the target of the equation does not amplify the error of its inverse observable. Changing the source variable and then reconstructing it generally does, as discussed below.

\subsection{Exponential phases}

Suppose an expression can be written on the selected real source branch as
\begin{equation}
 F(x)=y_0+\varepsilon e^{\Phi(x)},\qquad \varepsilon\in\{1,-1\}.
 \label{eq:coordinate-phase}
\end{equation}
An eventually signed amplitude is included in the phase: if the original expression is $y_0+A(x)e^{P(x)}$, prove the eventual sign $\varepsilon$ of $A$, then set $\Phi=P+\log(\varepsilon A)$. The inverse equation becomes
\begin{equation}
 \Phi(x)=v,\qquad v=\log\bigl(\varepsilon(y-y_0)\bigr),
 \qquad \varepsilon(y-y_0)>0.
 \label{eq:coordinate-log-target}
\end{equation}
The phase may contain several powers and polynomial--logarithmic terms even when the original function has exponential growth or decay. It is inverted on the same source interval. If $\Phi\to+\infty$, then $F\to\varepsilon\infty$; if $\Phi\to-\infty$, the target approaches $y_0$ from the side with sign $\varepsilon$.

\begin{proposition}[The original residual and the phase residual]
\label{prop:phase-residual}
For a real approximation $X$ on that source interval, put $D=\Phi(X)-v$. Then
\begin{equation}
 \frac{F(X)-y}{y-y_0}=e^D-1.
 \label{eq:phase-original-residual}
\end{equation}
In particular, if $D\to0$, the relative original residual is $D+\OO(D^2)$. If $|\Phi'|\ge\mu>0$ between $X$ and the true inverse $g(y)$, then $|X-g(y)|\le|D|/\mu$.
\end{proposition}
\begin{proof}
Since $y-y_0=\varepsilon e^v$, divide $\varepsilon(e^{\Phi(X)}-e^v)$ by $\varepsilon e^v$. The remaining statements follow from the Taylor expansion of $e^D$ at $0$ and the mean value theorem for $\Phi$.
\end{proof}

Thus a phase residual can be numerically useful without forming enormous exponentials. A numerical solve in the phase remains numerical evidence unless the derivative and residual have rigorous bounds; it nevertheless refers to the original equation through the exact identity \eqref{eq:phase-original-residual}.

\begin{example}[Several powers in an exponential]
For $F(x)=e^{x^2+x}$ at $+\infty$, put $v=\log y$. The inverse is
\[
 x=\frac{-1+\sqrt{1+4v}}2
 =\sqrt v-\frac12+\frac1{8\sqrt v}
 -\frac1{128v^{3/2}}+\OO(v^{-5/2}).
\]
This is an ordinary inverse expansion in the phase variable $v$. Its local target variable is $1/v=1/\log y$, and its remainder is measured in that coordinate. For $e^{-x^2-x}$ at the same source endpoint, replace $v$ by $-\log y$ and take $y\to0^+$.
\end{example}

\begin{example}[A nonconstant amplitude]
For $F(x)=xe^{x^2+x}$ at $+\infty$, the phase is $\Phi=x^2+x+\log x$. In terms of $v=\log y$, inversion begins
\[
 x=\sqrt v-\frac12+
 \frac{\tfrac18-\tfrac14\log v}{\sqrt v}
 +\OO(v^{-1}(1+|\log v|)^2).
\]
The logarithm of the amplitude contributes at the displayed order. Omitting it would preserve the leading term but give the wrong next coefficient. Translating $x$ to $x-c$ translates the reconstructed inverse by $c$ and leaves the phase calculation in that shifted source coordinate.
\end{example}

\begin{example}[A flat forward function with an elementary inverse]
$F(x)=e^{-1/x}$ on $x\to0^+$ has no nonzero power--log forward expansion, but its inverse is exactly
\[
 x=-\frac1{\log y},\qquad 0<y<1.
\]
The logarithmic target transformation turns the equation into $-1/x=\log y$. This example does not require an expansion into flat perturbation sectors: the exponential is part of the exact coordinate transformation.
\end{example}

Variable powers can enter the same route when their real domain is established. For $x^x$ on $x>0$, the phase is $x\log x$, so its inverse at $+\infty$ reduces to the Lambert-coordinate inversion of that phase. The identity $A(x)^{B(x)}=\exp(B(x)\log A(x))$ is used only where $A(x)>0$.

\subsection{Source reconstruction and its error}

\begin{theorem}[A change of source coordinate]
\label{thm:source-coordinate}
Let $x=H(t)$ be a continuously differentiable one-to-one parametrization of the selected source branch, and write $K(t)=F(H(t))$. Let $h$ be its inverse and suppose
\[
 h(y)=A(y)+E(y),\qquad E(y)=\OO(R(y)).
\]
If $H$ is defined on the interval between $A$ and $A+E$, and $|H'|\le B(y)$ there, then
\[
 F^{-1}(y)=H(A(y))+\OO(B(y)R(y)).
\]
The same conclusion applies to an observable after replacing $H$ by that observable composed with $H$.
\end{theorem}
\begin{proof}
The branch identity gives $F^{-1}=H\circ h$. Apply the mean value theorem between $A$ and $h=A+E$.
\end{proof}

For the coordinate transformations used here, more informative exact identities specify when a small error survives reconstruction.

\begin{corollary}[Exponential reconstruction]
\label{cor:exponential-reconstruction}
For a positive source coordinate $u=e^{s t}$, $s\ne0$, if $h=A+E$ and $E\to0$, then
\begin{equation}
 u=e^{sA}\bigl(1+\OO(E)\bigr),\qquad
 u^r=e^{rsA}\bigl(1+\OO(E)\bigr)
 \label{eq:exponential-reconstruction}
\end{equation}
for every fixed real $r$. If $E$ does not tend to zero, an additive expansion for $h$ alone does not justify these relative error assertions.
\end{corollary}
\begin{proof}
$u=e^{sA}e^{sE}$, and $e^{sE}=1+\OO(E)$ when $E\to0$. Replace $s$ by $rs$ for the observable.
\end{proof}

\begin{corollary}[Logarithmic reconstruction]
\label{cor:logarithmic-reconstruction}
For $x=-\log t/\lambda$, $\lambda\ne0$, suppose the positive inner inverse satisfies $h=A+E$ with $A>0$ and $E/A\to0$. Then
\begin{equation}
 x=-\frac{\log A}{\lambda}
 -\frac1\lambda\log\left(1+\frac EA\right)
 =-\frac{\log A}{\lambda}+\OO(E/A).
 \label{eq:logarithmic-reconstruction}
\end{equation}
Thus the inner remainder must be divided by the inner leading scale before it becomes an error in $x$.
\end{corollary}
\begin{proof}
Factor $h=A(1+E/A)$ and expand the logarithm of the positive unit factor. The hypothesis $E/A\to0$ also guarantees $h>0$ eventually.
\end{proof}

\subsection{Pure logarithmic functions}

For $x\to+\infty$, set $t=\log x\to+\infty$. For a positive source distance $u\to0^+$, set $t=-\log u\to+\infty$, then reconstruct $u=e^{-t}$ and the original point $x=x_0\pm u$. A pure logarithmic forward expression becomes a power--log expression in $t$. Its inverse is computed there before exponential reconstruction.

\begin{example}[The constant term in an exponent matters]
For $F(x)=(\log x)^2+\log x$ at $+\infty$, the transformed equation is $t^2+t=y$. Hence
\[
 t=\sqrt y-\frac12+\frac1{8\sqrt y}+\OO(y^{-3/2}),
\]
and one valid reconstruction is
\[
 x=\exp\left(\sqrt y-\frac12+\frac1{8\sqrt y}\right)
 \bigl(1+\OO(y^{-3/2})\bigr).
\]
In contrast, $t=\sqrt y+\OO(1)$ does not justify $x=e^{\sqrt y}(1+\oo(1))$: the omitted constant would change the leading multiplicative factor by $e^{-1/2}$. An implementation must retain enough inner terms to make its omitted additive error tend to zero, even when the initial term request would otherwise stop earlier.
\end{example}

This precision requirement belongs to reconstruction, not to the mere existence of the inner asymptotic expansion. When a user requests a particular relative order in the reconstructed expression, the implementation must translate that request into an inner additive precision and record both meanings.

\subsection{Finite sums of exponentials}

Let
\[
 F(x)=\sum_{j=1}^m a_j e^{\mu_j x},
 \qquad a_j,\mu_j\in\R,
\]
with equal rates merged. At $x\to+\infty$ choose $\lambda>0$, and at $x\to-\infty$ choose $\lambda<0$, so $t=e^{-\lambda x}\to0^+$ in either case. Then
\begin{equation}
 F\left(-\frac{\log t}{\lambda}\right)
 =\sum_j a_j t^{-\mu_j/\lambda}.
 \label{eq:finite-exponential-coordinate}
\end{equation}
This is a finite real-power germ. Rational relations among the rates are unnecessary: local finiteness of the positive correction gaps supplies finite exponent truncations. Its leading coefficient and exponent determine the target side, while the chosen sign of $\lambda$ preserves the source endpoint. After inversion, Corollary~\ref{cor:logarithmic-reconstruction} transports the error back to $x$.

\begin{example}[An exact oracle]
For $F(x)=e^x+e^{2x}$ at $+\infty$, take $t=e^{-x}$. Then $F=t^{-1}+t^{-2}$, and an exact inverse is
\[
 x=\log\left(\frac{\sqrt{1+4y}-1}{2}\right)
 =\frac12\log y-\frac1{2\sqrt y}
 +\frac1{48y^{3/2}}+\OO(y^{-5/2}).
\]
The inverse of $e^x+e^{\sqrt2 x}$ uses the same coordinate with irrational exponents. Constant terms, duplicate rates and cancellations are handled before extracting the leading transformed monomial.
\end{example}

\subsection{Result contracts and verification}

The transformed result retains its inner inverse, coordinate substitution, original function and selected domains. Its scale is identified as \wl{"Transformed"}. The logarithmic target route has coordinate kind \wl{"TargetLog"}; it uses \eqref{eq:coordinate-log-target}. The target expression is recorded explicitly. The stored cutoff belongs to the underlying series coordinate after substitution. An ordinary \wl{SeriesData} in the original target variable is unavailable when that representation would hide the coordinate change. The package README lists the corresponding property names.

The transformed residual describes the model actually composed, and its metadata states that scope. Numerical checks solve the transformed equation at high precision and verify the original source side; they remain numerical evidence. A separate certificate can use Proposition~\ref{prop:phase-residual} with rigorous derivative and residual bounds. An additive input error for $F$ must be transported through the target map before it can be treated as an error in $\Phi$; until that contract is supplied, rejecting an unsupported declaration is mathematically necessary.

These coordinate theorems justify the stated reductions independently of how many recognizers are implemented. The staged implementation plan tracks the accepted families, source reconstruction operations, tests and remaining extensions; the existence of a reduction alone does not establish that every corresponding expression has already passed the package's parser and native test suite.

\section{Weighted enumeration and grouped coefficient formulas}
\label{sec:incremental}

The coefficient theorem admits several finite constructions at a prescribed cutoff. This section organizes them by exact weight, by increasing Newton precision, and by marker degree. All constructions use finitely many positive gaps $\delta_i$ and preserve the meaning of a complete power--log block.

\subsection{Enumeration by complete weight layers}

Let $I$ be the set of retained multi-indices, initially empty, and let its
candidate frontier initially consist of the zero multi-index. At each step,
select the smallest candidate weight $\sigma$, promote \emph{every} candidate
of that weight into $I$, and add all their one-step successors
$\bk+e_i$ to the candidate set. The frontier is a set of integer vectors, so a candidate reached through several predecessors occurs only once. Distinct candidates of equal weight remain distinct until their coefficients are summed.

\begin{proposition}[Complete layers and coefficient reuse]
After the layer at weight $\sigma$ is processed, $I$ is exactly
$\{\bk\in\N^m:\bk\cdot\bdelta\le\sigma\}$, and the candidate set is its
one-step outer boundary. This construction determines every retained multi-index coefficient once. Summing the coefficients of an entire layer before
counting a block preserves all resonances and exact cancellations.
\end{proposition}
\begin{proof}
Every nonzero multi-index has a predecessor of strictly smaller weight because
all gaps are positive. Inductively, every index of the next candidate weight
has already been generated by a processed predecessor. A successor of an index
in the current layer has strictly larger weight, so processing one member of
the layer cannot expose a previously unseen member of the same layer.
The set union identifies repeated occurrences of a candidate through different predecessors. The
coefficient of each weight is the complete finite sum in
Theorem~\ref{thm:coefficients}, whether or not that sum is zero.
\end{proof}

Extending a cutoff requires only subsequent layers: coefficients below the
previous cutoff do not depend on the new cutoff. A prescribed exponent
cutoff is reached in finitely many layers by local finiteness. A prescribed
number of nonzero blocks is a different requirement, because complete
layers can cancel. Without a separate nontermination or exact-termination
argument, enumeration alone does not prove that another nonzero block exists.

\subsection{Newton iteration with growing precision}

Write $E(U)=0$ for \eqref{eq:normalized}, and let
$\delta_* = \min_i\delta_i$. The zero jet is correct below $\delta_*$.
If $U_s$ is correct strictly below weight $s$, one Newton step performed
strictly below
\[
 h=\min(2s,H),\qquad
 U_h=\left[U_s-\frac{E(U_s)}{E'(U_s)}\right]_{<h},
\]
is correct strictly below $h$. Indeed, $E'(U_s)$ has nonzero constant
coefficient $p$, and the Newton error is quadratic in the old error. Products,
unit powers, and logarithmic composition do not decrease the weight of the
error. The argument uses exponent valuation, with polynomial logarithms as
coefficients; logarithmic degree need not stay bounded as $h$ increases.

The working cutoffs are
$\min(2\delta_*,H),\min(4\delta_*,H),\ldots,H$, rather than $H$ at every
step. Their construction uses exact comparisons, with no floating-point
estimate of the required iteration count. The identity $[E(U_h)]_{<h}=0$ characterizes the formal accuracy of each step. Together with triangular uniqueness it proves that the retained coefficients agree with those of the inverse. A numerical bound at a particular nonzero $z$ additionally requires an analytic remainder estimate.

If a unit jet is already correct below $s$, refinement to $H>s$ starts with
that same jet and the same doubling argument. Truncation to a smaller
cutoff preserves correctness. For a pure monomial forward function the
unit correction is exactly zero, so no refinement is needed.

\subsection{Combine resonances before applying Euler operators}

Put
\[
 H_1(z,L)=\sum_i z^{\delta_i}B_i(L),\qquad
 H_1(z,L)^n=\sum_\sigma z^\sigma A_{n,\sigma}(L).
\]
The multinomial theorem gives
\[
 A_{n,\sigma}=
 \sum_{\substack{|\bk|=n\\\bk\cdot\bdelta=\sigma}}
 \frac{n!}{\bk!}\,B^{\bk}.
\]
At fixed depth $n$ and weight $\sigma$, all summands of
\eqref{eq:master} have the same differential operator. By linearity, the unit
bracket of $g^r$ can therefore be computed as
\begin{equation}
 \frac{g(y)^r}{z^r}
 =1+\sum_{n\ge1}\sum_\sigma z^\sigma
 \frac{(-1)^n r}{p^n n!}
 \prod_{j=1}^{n-1}(\D+r+\sigma+pj)A_{n,\sigma}(L).
 \label{eq:grouped-lagrange}
\end{equation}
For cutoff $H$, successive sparse products compute
$[H_1^n]_{<H}$. Only $n$ with $n\delta_*<H$ can contribute, so the computation
is finite. Coefficients of equal weight are combined during each sparse
product, before differentiation. Contributions at the same weight but
different depths are then combined in the result; they too can cancel.

Equation~\eqref{eq:grouped-lagrange} is especially useful when the generators
have many additive relations. It does not require commensurable gaps, and it
does not replace the individual-coefficient formula: that formula
retains the distinct perturbation markers and their provenance. Grouping replaces several applications of the same linear differential
operator by one application to the sum. The boundary argument of
Section~\ref{sec:remainders} is still needed to identify a complete omitted
weight. Neither grouping nor marker-depth truncation changes the ordering
of the final nonzero blocks. The equivalence of the grouped formula,
the individual multi-index formula and Newton iteration follows from
linearity and formal uniqueness.

\section{Refinement by retaining complete coefficient states}
\label{sec:refinement-state}

The public operation \wl{SeriesRefine[s,h]} reuses an ordinary inverse
when its stored forward model remains sufficient at the requested cutoff.
The implementation in \path{Kernel/RefinementState.wl} supports the
individual Lagrange and Newton methods. Its state is stored in the returned
object; there is no global cache, and refinement leaves the source object
unchanged. Coordinate wrappers, grouped coefficient generation, and inputs
whose forward expansion needs additional terms use the existing source
replay path.

\paragraph{Complete Lagrange prefixes.}
Write the normalized inverse observable as
\[
 z^r\sum_{\omega<H}z^\omega Q_\omega(\log z),
 \qquad \omega=\boldsymbol{k}\cdot\boldsymbol{d},\quad
 \boldsymbol{k}\in\mathbb N^m,\quad d_j>0.
\]
The retained state contains a downward closed set of processed indices,
its one-step outside boundary, and the complete coefficient polynomial at
every processed weight. Zero polynomials are omitted from the displayed
blocks but their indices remain processed. Advancing the state takes
\emph{all} boundary indices of the smallest weight simultaneously. Every
predecessor of such an index has smaller weight, because the gaps are
positive. Consequently, no as yet undiscovered index can have the same
weight, and resonant contributions are combined before the new block is
used or displayed.

An inverse produced with a term goal may already retain this state. For
an ordinary inverse constructed at an explicit cutoff, the complete
stored blocks instead seed the state. Enumerating
$\{\boldsymbol{k}:\boldsymbol{k}\cdot\boldsymbol{d}<H\}$ and its boundary
recovers the index bookkeeping without reevaluating the previously
computed Euler coefficients. The aggregate blocks remain authoritative;
the seed does not attempt to split a resonant coefficient polynomial into
its individual contributions.

The first omitted block is also evaluated by complete weight layers.
These frontier coefficients are retained for later refinement, including
zero layers skipped after cancellation. Caching a frontier can require
more indices than merely evaluating that frontier for a remainder. If
the additional state would exceed \wl{"MaxTerms"}, the implementation
performs the ordinary independent frontier calculation without caching
it. Thus this optimization does not reject an otherwise admissible
requested cutoff solely because it tried to retain an extra frontier.
As in ordinary construction, the bounded search through cancelled
frontier layers may return a conservative remainder at the first possible
omitted weight.

\paragraph{A bounded cache of logarithmic polynomial powers.}
The Lagrange state also retains powers of each nonconstant logarithmic
coefficient polynomial. Exponents $0$ and $1$ refer to $1$ and the original
model coefficient; larger cached powers form an initial consecutive
segment. A new power is obtained by multiplying the last retained power
by the original polynomial. Thus the product entering an Euler coefficient
can reuse polynomial powers requested by earlier indices or earlier
refinement calls. The Euler differential operators and normalization
remain unchanged, and the ordinary uncached coefficient routine provides
an independent exact comparison.

Cached powers use only unused \wl{"MaxTerms"} slots after the processed
index region and its boundary. When the region grows, the largest retained
powers are evicted until their entry count fits the remaining capacity.
If no slot is available, the requested coefficient product is evaluated
directly. Eviction never discards an inverse coefficient or changes the
remainder calculation. In particular, optional caching cannot turn an
otherwise admissible index region into a resource failure. This budget
counts retained algebraic objects; it is not a byte bound on a polynomial
whose individual coefficients may themselves be large exact expressions.

\paragraph{Recovering and extending a Newton state.}
Suppose the displayed normalized observable is
$(1+U)^r$ with $r\ne0$. Its constant coefficient is $1$, so the formal
power with exponent $1/r$ uniquely recovers $1+U$ in the filtered unit
algebra. This recovers an initial Newton state from an existing observable,
including inverse powers arising from a source endpoint at infinity.
The recovered unit is checked by substituting it into the normalized
implicit equation: the residual must vanish exactly below the claimed
precision. Only then is the state admitted for further Newton steps.

If a Newton state is already present, refinement retains its unit jet,
precision, and recorded step cutoffs. Each new step increases the
precision to the smaller of twice the old precision and the requested
precision. The existing exact residual check applies at every step.
After a coarser display is requested, a stronger retained state remains
available to subsequent calls. The formal branch is selected by constant
term $1$ throughout; no numerical root selection enters this recovery.

\paragraph{Model identity and input precision.}
Retained states are bound to the normalized model, observable power,
assumptions, original function, and source endpoint and direction. A
stored signature and direct comparisons detect incompatible states, and
the displayed prefix must agree with the corresponding state prefix.
This consistency check prevents accidental reuse after changing a model;
it is not a certificate for arbitrarily edited internal associations.

A declared forward error remains a hard cap on usable inverse precision.
For a leading exponent $p$ and internal observable exponent $r$, a forward
error $\OO(u^\rho |\log u|^k)$ permits target cutoff at most
\[
 \frac{\rho-p+r}{|p|}.
\]
The remainder's logarithmic degree is combined with that transported
error at equal power. An error introduced by automatic expansion of an
analytic source is distinguished from a declared error: if refinement
requires more source terms, the constructor is replayed and the old
model state is not reused. No cached inverse coefficients justify going
beyond an explicit input-error declaration.

\paragraph{What the counters establish.}
Each refined object exposes \wl{"RefinementStatistics"} and appends it to
\wl{"RefinementHistory"}. The Lagrange counters have distinct meanings:
\begin{description}[style=nextline,leftmargin=0pt,labelwidth=0pt,labelsep=0pt,font=\normalfont]
\item[\wl{"ReusedCoefficientEvaluations"}]
Available individual index contributions. For a seed reconstructed from
stored complete blocks this is the size of the old processed region,
not work performed during seeding.
\item[\wl{"NewCoefficientEvaluations"}]
Actual new Euler coefficient evaluations, including any uncached frontier
work. Supplying a cached polynomial product still counts as one coefficient
evaluation.
\item[\wl{"ReusedBlocks"}]
Old nonzero blocks retained below the requested cutoff.
The full old display has its own count,
\wl{"AvailableSourceBlocks"}.
\item[\wl{"ReusedPolynomialPowerRequests"}]
Requests satisfied by already retained polynomial powers during this call.
\item[\wl{"NewPolynomialPowerEvaluations"}]
New expanded polynomial-power computations, including uncached requests
when the optional cache is full. This count is separate from coefficient
evaluations and excludes constant coefficients. It covers incremental
coefficient generation; an independent uncached frontier calculation is
counted separately at the coefficient level.
\item[\wl{"RetainedPolynomialPowers"}]
The number of powers of degree at least $2$ retained in the final state.
Capacity and eviction counters record how the resource budget was applied.
\end{description}
Newton statistics separately report the old and new precision, the new
step cutoffs, and Euler work used for the remainder frontier. These are
exact bookkeeping quantities; they are not a count of all arithmetic
operations and do not prove a runtime bound.

The reproducible driver \path{Tests/BenchmarkRefinement.wl} compares a
sequence of fresh public inverse constructions with one construction
followed by public refinements. It covers dense rational gaps, two and
three irrational gaps, resonant weights, and high logarithmic degrees,
using both Lagrange and Newton examples. A deep one-gap case reaches target
cutoff $64$ and also compares with the independent Catalan-number formula
for the inverse of $x+x^2$. Acceptance requires equality of
every complete block and remainder at every requested cutoff, together
with the retained-state invariants. The driver reports elapsed times and
state counters. It measures retained expression and state sizes with
\wl{ByteCount} and evaluation peaks with \wl{MaxMemoryUsed[expr]}; these
kernel measurements are not total-process memory bounds. No
machine-dependent speed or memory threshold determines acceptance.

\section{Calculus of expansions with transported precision}
\label{sec:series-calculus}

The operations \wl{SeriesAdd}, \wl{SeriesMultiply} and \wl{SeriesPower}
provide arithmetic without discarding remainders. The functions
\wl{SeriesLog} and \wl{SeriesExp} handle logarithms and exponentials;
\wl{SeriesCompose} and \wl{SeriesObservable} handle composition and other
observables. Truncation, refinement and justified differentiation use
\wl{SeriesTruncate}, \wl{SeriesRefine} and \wl{SeriesDifferentiate}.
All ten operations share a common representation
\begin{equation}
 F(y)=b(y)+c(y)\left[J(w(y))+\OO(w(y)^P\M(w(y))^D)\right],
 \qquad w(y)\longrightarrow0^+,
 \label{eq:calculus-carrier}
\end{equation}
where $b,c,w$ are exact expressions and $J$ is a finite power--log jet.
The offset $b$ and prefactor $c$ are recorded separately. Thus a cutoff in
the jet describes relative precision with respect to $c$, and the absolute
remainder contains $|c|$. Ordinary forward and inverse objects, Lambert
objects, and explicitly recorded coordinate substitutions can be converted
to this representation. An operation returns a derived object and its replay
recipe; it does not retain an inverse-model label that would falsely suggest
that the original inverse residual checker applies to an arbitrary observable.

For addition and multiplication, the precision rules of
Section~\ref{sec:scale} are used after expressing both operands in the same
positive coordinate. An exact prefactor may be absorbed into the jet when
it is a finite power--log expression. Distinct prefactors can be aligned
when their ratio lies in that algebra. The package rejects incompatible
coordinates or independently ordered exponential carriers rather than
assigning them an unsupported single cutoff.

\paragraph{Powers and logarithms.}
For a positive monomial leading term $a w^\alpha$, write
$J=a w^\alpha(1+U)$. A constant power and a logarithm reduce to the unit
series of $(1+U)^r$ and $\log(1+U)$. The input uncertainty first becomes
$\OO(w^{P-\alpha}\M^D)$ in $U$. Multiplication by the output prefactor
then transports it to the correct absolute scale. Real noninteger powers
require positivity on the selected branch. These operations retain any
earlier input precision ceiling; requesting a higher output cutoff does
not create unknown coefficients.

\paragraph{Exponentials and exact carriers.}
Suppose
\[
 H=H_{\le0}+H_{>0}+\OO(w^P\M^D),\qquad P>0,
\]
where $H_{\le0}$ contains the nonpositive-weight blocks. Then
\[
 e^H=e^{H_{\le0}}
 \left[e^{H_{>0}}+\OO(w^P\M^D)\right].
\]
Indeed the omitted argument tends to zero and $e^r=1+\OO(r)$ there.
The operation keeps $e^{H_{\le0}}$ as an exact carrier and expands only
the vanishing part. For example $H=w^{-1}+\log w+w$ yields
$w e^{1/w}(1+w+w^2/2+\cdots)$. This also supplies the precision rule
needed when a source logarithm is inverted and the source variable must
be reconstructed by exponentiation. A nonvanishing argument uncertainty
is rejected: it cannot justify a relative error tending to zero after
exponentiation.

\paragraph{Composition and regular observables.}
The outer local coordinate must become $a w^\alpha(1+o(1))$ with
$a,\alpha>0$ after substitution of the inner expansion. This condition
checks both the limiting point and its selected side. An outer remainder
$\OO(v^P\M(v)^D)$ becomes $\OO(w^{\alpha P}\M(w)^D)$, since
$\log v=\alpha\log w+\OO(1)$. The uncertainty of the inner argument is
propagated independently through every jet operation, and the worse
precision is retained. Regular analytic observables at a finite limiting
argument use their Taylor series and the same precision arithmetic.
Poles, branch points and logarithmic leading coefficient functions require
their respective scale normalizations; a general analytic-observable call
does not silently assume them to be regular.

\paragraph{Differentiation requires derivative information.}
An estimate $R(w)=\OO(w^P\M(w)^D)$ alone does not imply any estimate on
$R'$. For example $w^P\sin(e^{1/w})$ has that magnitude bound and a much
larger derivative. Therefore a nonzero remainder can be differentiated
$n$ times only with a recorded or explicitly supplied contract
\[
 R^{(j)}(w)=\OO(w^{P-j}\M(w)^D),\qquad 0\le j\le n.
\]
The option \wl{"RemainderDerivativeOrder" -> n} is a declaration of
these hypotheses, not a numerical test of them. Each differentiation
consumes one order of the contract. The jet derivative uses
$D_w[w^\beta Q(\log w)]=w^{\beta-1}(\beta Q+Q')$; the derivatives
of the exact coordinate, offset and prefactor in
\eqref{eq:calculus-carrier} are included by the chain and product rules.
Exact finite expressions require no remainder hypothesis. A discarded tail
of a known finite power--log expression also has derivative bounds of every
order. In contrast, a derivative bound declared for one unknown remainder
is not automatically promoted to a stronger exponent when the source is
refined; that stronger bound must be established separately.

\paragraph{Truncation and refinement.}
Truncation moves the first discarded block, including its logarithmic
degree, into the remainder. It cannot improve uncertainty inherited from
the input. Refinement instead replays a retained forward expression,
inverse problem, coordinate substitution or operation recipe at a higher
working order. The replay still respects transported input precision
ceilings. Without a retained source or recipe, the package reports that
refinement is unavailable; changing an inert remainder exponent would not
be a mathematical refinement.

An automatically computed forward remainder is distinct from a declared
input error. The object records the latter separately as
\wl{DeclaredInputRemainder}. When refining an analytic input such as
$\sin x$, the original function can be expanded further; its previous
Taylor remainder must not become a permanent declared accuracy ceiling.
A remainder explicitly supplied for an unspecified omitted function
remains a binding ceiling. Tests exercise both cases with independent
inverse coefficients.

\section{Perturbations of a retained exact inverse}
\label{sec:core-perturbations}

An inverse such as a real branch of the Lambert function can be useful as an exact building block even when it also has a convergent inverse--logarithmic expansion. Retaining that building block separates two different refinements: expansion of the core inverse itself, and correction for terms outside the core. The operation \wl{AsymptoticCoreInverse} implements the second refinement. Its truncation parameter counts powers of an auxiliary perturbation marker. The returned coefficients are complete functions of the exact core inverse; they are not individual terms of an exponent-sorted power--logarithmic series.

\subsection{The coefficient identity}

Choose the positive source coordinate $u\to0^+$ and write the original source point as
\[
 x=H(u),\qquad
 H(u)=x_0+\sigma u\quad\hbox{or}\quad H(u)=\frac{\sigma}{u},
 \qquad \sigma\in\{1,-1\}.
\]
Let $F_0(u)$ be the chosen core and $R(u)$ its perturbation in that coordinate. Let $u_0(y)$ be the selected exact inverse of $F_0$, and set $\phi(y)=H(u_0(y))$. Introduce a marker $\lambda$ by
\begin{equation}
 F_0(u_\lambda)+\lambda R(u_\lambda)=y,
 \qquad u_{\lambda=0}=u_0(y).
 \label{eq:core-marker-equation}
\end{equation}
At any fixed point with $F_0'(u_0)\ne0$, the inverse function theorem gives an analytic solution in $\lambda$ near zero whenever the data are analytic locally. The Lagrange--B\"urmann identity gives
\begin{equation}
 H(u_\lambda)=\phi(y)+\sum_{n\ge1}\lambda^n C_n(u_0(y)),
 \label{eq:core-marker-series}
\end{equation}
where
\begin{equation}
 C_n(u)=\frac{(-1)^n}{n!}
 \left(\frac{1}{F_0'(u)}\frac{d}{du}\right)^{n-1}
 \left(\frac{H'(u)R(u)^n}{F_0'(u)}\right).
 \label{eq:core-marker-coefficients}
\end{equation}
Indeed, with $t=F_0(u)$, the equation becomes
$t=y-\lambda R(u_0(t))$, and $d/dt=(F_0'(u))^{-1}d/du$.
Applying the ordinary Lagrange identity to the observable $\phi(t)$ proves
\eqref{eq:core-marker-coefficients}. This form of the identity differentiates
the explicit forward core. It avoids repeated symbolic differentiation of a
large expression containing \wl{ProductLog}.

For arbitrary data, this identity alone does not justify setting $\lambda=1$,
does not order the coefficients by their size as $u_0\to0$, and does not
bound the tail by its first omitted coefficient. The following restricted
class supplies those missing conclusions with a complete tail majorant.

\subsection{A class with an asymptotic tail bound}

Let $\M(u)=1+|\log u|$. Consider finite real power--logarithmic expressions
\begin{align}
 F_0(u)&=y_0+u^p Q_0(\log u)
       +\sum_{j=1}^{J}u^{p+\eta_j}Q_j(\log u),
       &p&\ne0,\quad \eta_j>0,\label{eq:core-admissible-core}\\
 R(u)&=\sum_{i=1}^{I}u^{p+\delta_i}P_i(\log u),
       &\delta_i&>0.\label{eq:core-admissible-perturbation}
\end{align}
All polynomials have real coefficients, $Q_0$ is nonzero, and all powers
are fixed real numbers. Equal powers are merged before the leading term
is selected. Put
\begin{equation}
 q=\deg Q_0,\qquad
 \delta=\min_i\delta_i,\qquad
 D=\max\left(0,\max_i\deg P_i-q\right),\qquad
 \chi(u)=u^\delta\M(u)^D.
 \label{eq:core-small-scale}
\end{equation}
For a zero perturbation the inverse remains exactly $\phi$, so these
parameters and the following tail estimate are unnecessary.

\begin{theorem}[An exact-core expansion with a complete marker tail]
\label{thm:core-marker-majorant}
For fixed data satisfying \eqref{eq:core-admissible-core}--\eqref{eq:core-admissible-perturbation}, the selected real inverse exists sufficiently near the source endpoint. Define $r_*=1$ for a finite source endpoint and $r_*=-1$ for an infinite source endpoint. There are positive constants $C,K$ and a source threshold such that, for every integer $N\ge0$ and all sufficiently small $u_0>0$ with $K\chi(u_0)<1$, the marker series converges at $\lambda=1$ and
\begin{equation}
 \left|F^{-1}(y)-\phi(y)-\sum_{n=1}^{N}C_n(u_0(y))\right|
 \le C u_0^{r_*}
 \frac{(K\chi(u_0))^{N+1}}{1-K\chi(u_0)}.
 \label{eq:core-geometric-tail}
\end{equation}
In particular, for every fixed $N$,
\begin{equation}
 F^{-1}(y)=\phi(y)+\sum_{n=1}^{N}C_n(u_0(y))
 +\OO\!\left(u_0^{r_*+(N+1)\delta}
                  \M(u_0)^{(N+1)D}\right).
 \label{eq:core-asymptotic-tail}
\end{equation}
The estimate concerns the complete omitted marker tail, including possible
cancellations in its first coefficient.
\end{theorem}
\begin{proof}
Write $u=u_0W$ and divide the equation
$F_0(u_0W)+\lambda R(u_0W)=F_0(u_0)$ by the leading scale
$u_0^p(\log u_0)^q$ and its nonzero leading coefficient.
Use $\tau=1/\log u_0$ and finitely many additional parameters for the
higher-power terms. For example,
\[
 \frac{(\log(u_0W))^k}{(\log u_0)^q}
 =(\log u_0)^{k-q}(1+\tau\log W)^k.
\]
If $k>q$, include the factor
$u_0^\eta(\log u_0)^{k-q}$ in its small parameter; if $k\le q$,
retain the nonnegative power $\tau^{q-k}$ as an analytic coefficient.
All such parameters tend to zero. Powers of $W$ and $\log W$ are analytic
in a fixed complex neighborhood of $W=1$ after fixing their branches.

At the limiting parameter origin, the normalized equation is $W^p-1=0$,
whose derivative at $W=1$ is $p\ne0$. The analytic implicit function theorem
therefore gives a holomorphic solution in a common polydisc. The core
parameters remain fixed when $\lambda$ varies, whereas every perturbation
parameter is multiplied by $\lambda$. Each perturbation parameter is
bounded in magnitude by a constant times $\chi(u_0)$ for small $u_0$.
Cauchy estimates in a smaller fixed polydisc consequently bound the total
coefficient of marker degree $n$ by $C K^n\chi(u_0)^n$. Equivalently,
the solution is holomorphic and uniformly bounded on a marker disc whose
radius is a fixed positive multiple of $1/\chi(u_0)$.

For a finite source endpoint, the nonconstant part of the observable is
$\sigma u_0W$; at infinity it is $\sigma u_0^{-1}W^{-1}$.
Both normalized observables are analytic near $W=1$, giving the additional
factor $u_0^{r_*}$. Summing the resulting geometric coefficient bound gives
\eqref{eq:core-geometric-tail}. The local real branch is unique, and
$F_0'(u)\sim p u^{p-1}Q_0(\log u)$ while $R'/F_0'\to0$,
so this analytic solution agrees with the selected real inverse near the
endpoint. Finally $\chi(u_0)\to0$, which proves
\eqref{eq:core-asymptotic-tail}.
\end{proof}

The constants in this theorem are existence bounds for fixed input data.
The implementation does not compute a numerical value of $C$, $K$, or the
source threshold. Its \wl{"MajorantContract"} therefore identifies the
bound as asymptotic existence and sets \wl{"NumericCertificate"} to
\wl{False}. The stated \wl{"SourceRadius"} controls a supplied inverse's
symbolic branch check; it is not asserted to be the unknown uniform
threshold in \eqref{eq:core-geometric-tail}.

The nonnegative value of $D$ in \eqref{eq:core-small-scale} is deliberately
conservative. A quotient by an extra power of $|\log u_0|$ can make an
individual coefficient smaller than the reported polynomial--log bound.
Keeping that improvement would require an additional logarithmic ordering
contract. Also, if the perturbation contains several different powers,
one complete marker coefficient may contain contributions lying beyond
the leading contribution of a later marker coefficient. Marker depth does
not promise a globally sorted list of all resulting monomials.

\subsection{The Lambert core and its first correction}

For the exact core
\[
 F_0(u)=y_0+a u^p(\log u+b)^q,\qquad p\ne0,
 \qquad q\in\Z_{\ge0},
\]
the parser verifies that the leading logarithmic polynomial is exactly the
displayed affine power. For $q=0$, the positive local inverse is
$u_0=((y-y_0)/a)^{1/p}$. For $q>0$, put
\[
 A=a(-1)^q,\qquad Y=\frac{y-y_0}{A}>0,\qquad
 k=-\frac pq,\qquad z=kY^{1/q}e^{pb/q}.
\]
Then
\begin{equation}
 u_0=Y^{1/p}\left(\frac{W_\nu(z)}{k}\right)^{-q/p},
 \qquad
 \nu=\begin{cases}-1,&p>0,\\0,&p<0.\end{cases}
 \label{eq:core-exact-lambert-inverse}
\end{equation}
These are the real branches for which $u_0\to0^+$. On the lower branch,
the real argument condition $-1/e\le z<0$ is recorded explicitly. The
quotient $W_\nu(z)/k$ is positive on the selected eventual branch, so all
displayed real powers have positive bases.

\begin{example}[Correcting $u\log u$ without expanding its exact inverse]
\label{ex:core-lambert-quadratic-perturbation}
For $F(u)=u\log u+u^2$ as $u\to0^+$, let
$u_0=y/W_{-1}(y)$. Formula~\eqref{eq:core-marker-coefficients} gives
\begin{align}
 C_1(u_0)&=-\frac{u_0^2}{1+\log u_0},\label{eq:core-lambert-first-correction}\\
 C_2(u_0)&=\frac{u_0^3(4\log u_0+3)}{2(1+\log u_0)^3}.
 \label{eq:core-lambert-second-correction}
\end{align}
Here $p=q=\delta=1$ and $D=0$. Thus truncation after $C_1$ has the
proved remainder $\OO(u_0^3)$, and truncation after $C_2$ has remainder
$\OO(u_0^4)$. The explicit omitted coefficient can exhibit a sharper
logarithmic scale, but the package reports the complete-tail bound
established above. Each displayed correction retains the exact
$W_{-1}$ expression in $u_0$.
\end{example}

\subsection{A declared forward remainder}

Suppose the actual forward function also contains an unknown error $E(u)$,
with the declared bounds
\begin{equation}
 E(u)=\OO(u^\rho\M(u)^k),\qquad
 E'(u)=\OO(u^{\rho-1}\M(u)^k),\qquad \rho>p.
 \label{eq:core-declared-input-error}
\end{equation}
The derivative bound is a hypothesis of this declaration, not a consequence
of the first bound alone. Since the retained forward derivative has scale
$u^{p-1}\M(u)^q$, the mean value argument of
Section~\ref{sec:transport} gives an additional error in the original
source variable of
\begin{equation}
 \OO\!\left(u_0^{r_*+\rho-p}
       \M(u_0)^{\max(0,k-q)}\right).
 \label{eq:core-input-error-transport}
\end{equation}
The gap $\rho>p$ ensures that this error is relatively small in the
positive source coordinate. Both the perturbed solution and the core
solution are asymptotic to the same $u_0$, so they can be used
interchangeably in these scale estimates. The output combines
\eqref{eq:core-input-error-transport} with the marker-tail bound by taking
the smaller source power, and the larger logarithmic degree if those
powers coincide. Increasing marker depth cannot improve this inherited
accuracy ceiling. The displayed marker coefficients still describe the
specified finite model; coefficients smaller than the unknown input error
are not asserted to be coefficients of every admissible actual function.

\subsection{The implemented operation and its evidence boundary}

\begin{lstlisting}[language=Mathematica]
AsymptoticCoreInverse[
  x Log[x], x^2, {x, 0}, {y, 2}]

AsymptoticCoreInverse[
  x + x^2, x^3, {x, 0}, {y, 1},
  "CoreInverse" -> (Sqrt[1 + 4 y] - 1)/2]
\end{lstlisting}

The first call automatically selects the lower real Lambert branch. The
second retains a supplied algebraic inverse after a bounded symbolic check
of $\phi(F_0(H(u)))=H(u)$ for $0<u<$\wl{"SourceRadius"} on the selected
side. An unverified inverse is rejected. Supplied parameter assumptions
must establish the reality of the coefficients and the eventual leading
sign. Exponents in this operation are exact real numeric constants; they
need not be rational.

The result identifies its scale as \wl{"ExactCorePerturbation"}.
Its kind is \wl{"CoreInverse"}, and \wl{"MarkerTerms"} stores pairs $\{n,C_n\}$,
with degree zero holding $\phi$. The full exact inverse of the core,
the local core inverse, the separately computed first omitted marker
coefficient, the branch certificate, and the tail contract are all retained.
The ordinary \wl{Normal} operation returns the retained expression. The
object does not claim an ordinary \wl{SeriesData} representation or an
ordinary exponent cutoff. Its independent regressions compare marker
coefficients with direct composition and exact algebraic inverses, check
both source orientations and infinite-endpoint error scaling, and compare
the Lambert correction with the original forward equation at high
precision. Such numerical checks supplement the theorem; they do not
turn its unspecified constants into a numerical certificate.

This implemented class requires strictly positive source-power gaps in
the perturbation. Perturbations of equal source power but smaller
logarithmic size, purely logarithmic leading cores, and general
exponential exact cores require additional scale and error-transport
contracts. They are not silently accepted by the marker identity.
The broader roadmap also includes refinement of the retained core,
composition of core objects with other expansion carriers, and computable
majorants. The present operation supplies the finite power--logarithmic
perturbation theorem and its exact-core implementation as one independently
testable stage.

\section{Executable residual certificates}
\label{sec:executable-certificates}

An asymptotic remainder describes a limiting scale; it need not supply a
constant or a numerical neighborhood. A root obtained by numerical iteration
is useful independent evidence, but is not a verified enclosure.
The function \wl{InverseCertificate} implements the residual bracket of
Theorem~\ref{thm:bracket} on an explicit source interval. Its successful output
contains rational endpoints enclosing a unique root, a bound for the error of
the reported center, and the derivative and residual enclosures used to prove
those statements.

\subsection{The certificate theorem with interval data}

Let $I=[A,B]$, let $c\in(A,B)$, and suppose a structural interval evaluator
proves that $F$ is continuous on $I$ and differentiable inside, with
\[
 F'(I)\subseteq[d_-,d_+],\qquad
 F(c)-y\in[\rho_-,\rho_+].
\]
Assume either $d_->0$ or $d_+<0$, and put
\[
 \mu=\min(|d_-|,|d_+|),\quad
 M=\max(|d_-|,|d_+|),\quad
 \epsilon=\max(|\rho_-|,|\rho_+|),\quad h=\epsilon/\mu.
\]
If $[c-h,c+h]\subseteq I$, Theorem~\ref{thm:bracket}, after changing the sign
of $F$ if necessary, proves a unique root $x_*$ in $I$ and
\[
 |x_*-c|\le h.
\]
The mean-value identity further gives
\[
 x_*\in I\cap[c-h,c+h]\cap
 \left(c-\frac{[\rho_-,\rho_+]}{[d_-,d_+]}\right).
\]
Division is interval division; its denominator is separated from zero. This
last interval is the returned \wl{"RootEnclosure"}. If the residual interval
does not contain zero, there is also the lower bound
\[
 |x_*-c|\ge\frac{\min(|\rho_-|,|\rho_+|)}{M}.
\]
It can prove that a requested accuracy is impossible for a fixed center.

If the symmetric residual bracket does not fit inside $I$, the evaluator
also encloses the residual at $A$ and $B$. Proved opposite signs in the
orientation of $F'$ establish existence by continuity, and the separated
derivative still establishes uniqueness. Once existence in $I$ is known,
the mean-value theorem gives the same bound $|x_*-c|\le h$ and the same
intersection enclosure, without requiring the symmetric bracket to be
contained in $I$. The output identifies this alternative as exact endpoint
sign evidence. This case matters for a coarse asymptotic seed close to one
end of a valid interval: increasing arithmetic precision cannot repair
an unsuitable symmetric-bracket geometry.

The theorem asserts uniqueness on the supplied interval, which the evaluator
also checks lies on the requested side of the source endpoint. It does not
certify global injectivity or identify disconnected monotone intervals as the
same continuation of an inverse germ. This scope is recorded in the result.

\subsection{Rational arithmetic and elementary enclosures}

The proof path uses exact rational endpoints. Addition, multiplication,
reciprocals, and integer powers use the usual enclosing interval operations.
After arithmetic, an endpoint $q\ne0$ is rounded outward to a dyadic grid with
a prescribed number of significant bits. Its magnitude exponent is determined
from the binary lengths of the exact numerator and denominator, followed by
one rational comparison. Rounding uses $\lfloor q/\eta\rfloor\eta$ at a lower
endpoint and $\lceil q/\eta\rceil\eta$ at an upper endpoint. Thus even rounding
is justified by integer arithmetic. Approximate decimal values and
floating-point comparisons are absent from every enclosure and every final
sign or containment test.

For $0\le t\le\tfrac12$, the exponential evaluator uses
\[
 S_N(t)=\sum_{j=0}^N\frac{t^j}{j!},\qquad
 0\le e^t-S_N(t)\le
 \frac{t^{N+1}}{(N+1)!}\frac1{1-t/(N+2)}.
\]
The bound follows by majorizing the ratio of successive omitted terms by
$t/(N+2)$. A positive argument is reduced by powers of two, and the enclosure
is squared back using directed rounding. Negative arguments use the reciprocal
of the positive exponential enclosure. An explicit resource bound on the
argument prevents unbounded integer growth.

For $1\le q\le2$, put $v=(q-1)/(q+1)\in[0,\tfrac13]$. Then
\[
 T_N(q)=2\sum_{j=0}^{N-1}\frac{v^{2j+1}}{2j+1},\qquad
 0\le\log q-T_N(q)\le
 \frac{2v^{2N+1}}{(2N+1)(1-v^2)}.
\]
Every positive rational argument is written exactly as $2^k q$ with
$1\le q<2$, using binary numerator and denominator lengths. The identity
$\log(2^kq)=k\log2+\log q$ therefore reduces all logarithms to the displayed
bound. Both tails are positive rational expressions. Increasing $N$ improves
their analytic precision independently of the directed arithmetic precision.

Before enclosing a logarithm, exact structural identities can remove an
exponential whose real exponent has been enclosed, or split a product whose
factors each have a real logarithm enclosure. Positivity is therefore checked
before distributing logarithms. If, for example, a positive product has
negative factors, the evaluator retains the direct interval route. This avoids
materializing a huge exponential only to take its logarithm while preserving
the real-branch conditions.

Noninteger powers on a strictly positive base are evaluated through
$a^b=\exp(b\log a)$, including supported exact irrational exponents. A positive
rational power also admits a zero lower base endpoint. Expressions are handled
by structural recursion through sums, products, powers, exponentials and
logarithms. A reciprocal crossing zero, a logarithm reaching a nonpositive
argument, or an unsupported head produces a failure. In particular, a
numerical value of an unsupported expression is never accepted as a proof.
Successful structural evaluation establishes the needed continuity on the
whole interval; successful evaluation of the symbolic derivative establishes
the derivative enclosure. Symbolic parameters without exact supported values
are currently rejected with the unresolved expression recorded.

\subsection{Interface, phase coordinates, and adaptive accuracy}

The basic call is
\begin{lstlisting}
s = AsymptoticInverse[x + x^2, {x, 0}, {y, 3}];
InverseCertificate[s, 6, "Interval" -> {1, 3},
  "Center" -> 21/10]
\end{lstlisting}
The independently known root is $2$. The derivative enclosure is $[3,7]$;
the certified bracket encloses $2$ even though the asymptotic expansion was
not requested at a target where its truncation is especially accurate. The
residual theorem is independent of the formal truncation theorem.

Targets must be exact supported numeric expressions. Interval endpoints and
an explicitly supplied center must be rational, with the center strictly
inside the interval. With \wl{"Center" -> Automatic}, numerical evaluation of
the expansion chooses a rational center; an unusable seed is replaced by the
interval midpoint. This decimal calculation selects a candidate only. All
proofs are subsequently rebuilt with exact rational arithmetic. The returned
error bound applies to that recorded center, including its rationalization.
It does not silently claim an error bound for a different displayed expression.

The equation being certified is always the object's stored explicit
\wl{"Function"}. A declared \wl{"InputRemainder"} describes further possible
forward functions, but supplies no computable interval bound for their values
or derivatives. The certificate therefore does not cover those unspecified
terms. Results retain that descriptor, set
\wl{"CertifiesInputRemainderFamily" -> False}, and label the scope
\wl{"StoredExpressionOnly"} when it is nontrivial. An automatically computed
forward truncation can also populate the remainder field, so this label does
not guess the descriptor's provenance. The original function and exact target
are recorded separately from any phase equation. Asymptotic remainder data
are never silently promoted into a numerical error certificate.

For an object using a logarithmic target transformation, or a recognized
Lambert exponential core admitting the same exact transformation, the certificate
evaluates the exact phase and transformed target. The source-domain condition
of the coordinate map is proved on the entire interval first. Since the
exponential is strictly monotone there, the resulting root is a root of the
original equation. This avoids constructing enormous values such as
$\exp(x^2+x)$ merely to subtract two nearly equal numbers.

The option \wl{"TargetError" -> epsilon} requests an exact positive rational
tolerance. After an initial certificate, a sharper interval-Newton enclosure
can supply a new center and verification interval. The preceding certificate
already proves a root exists in that interval, so subsequent error estimates
need not re-prove existence by the initial bracket-containment condition.
Derivative bounds and residual bounds are nevertheless recomputed. Unless
disabled by \wl{"RefineExpansion" -> False}, refinement also tries a larger
asymptotic term goal as a source of a better seed; every accepted seed is
verified independently. A bounded failed or unsupported expansion refinement
does not invalidate an already established interval theorem.

An \wl{AsymptoticCoreInverse} object is also accepted as a seed. Its exact
Lambert core and marker corrections may be evaluated numerically to select
the center, while the certificate evaluates the full explicit
\wl{"Core" + "Perturbation"} equation with rational enclosures. If expansion
refinement is requested, the stored exact core inverse is replayed at a larger
marker depth. Its symbolic identity check and computation are time bounded;
a failed replay falls back to the already certified interval information.
The marker majorant's uncomputed constants are not used as numerical bounds.

With an explicit fixed center, refinement increases enclosure precision without
silently changing the center. A certified lower error bound above the target
causes \wl{Failure["AccuracyFloor", ...]}. Exhausting the refinement budget
causes \wl{Failure["AccuracyNotReached", ...]}, retaining the best established
certificate if one exists. Neither outcome is reported as a successful accuracy
request. The \wl{"History"} property records the attempted centers and
enclosure orders. Resource failures and unresolved derivative signs are
reported separately from mathematical impossibility.

\paragraph{Relative accuracy.}

An optional \wl{"RelativeError" -> tau} requests a relative tolerance with
an independently proved root magnitude. If the root interval $I$ is
separated from zero, set $m_I=\min_{z\in I}|z|$ and
$M_I=\max_{z\in I}|z|$. An absolute error bound $r$ proves the requested
combined tolerance when
\[
 r\le\max(\epsilon,\tau m_I),
\]
where $\epsilon$ is an explicitly supplied \wl{TargetError}, or zero when
none is supplied. This implies
$|A-x_*|\le\max(\epsilon,\tau|x_*|)$ for the unique enclosed root.
The reported relative error bound is $r/m_I$. A proved zero root requires
an explicit positive absolute fallback; no relative quotient at zero is
invented. If an interval merely crosses zero, further refinement can
establish a nonzero lower magnitude or report that the goal was not reached.
For a fixed-center impossibility claim, the opposite bound is needed:
the certified lower error must exceed $\max(\epsilon,\tau M_I)$.
Using $m_I$ in that last test would incorrectly turn a conservative
success criterion into an impossibility assertion.

\subsection{Validation scope}

The regression cases exercise outward rounding at both signs and widely
different magnitudes, explicit exponential and logarithmic bounds, interval
powers crossing zero, positive and negative derivatives, and a zero residual.
Known rational roots provide independent checks of the final enclosures.
Further cases cover the lower Lambert branch, a logarithmic phase, successful
adaptive accuracy, a proved fixed-center accuracy floor, and rejection of
rounded targets, turning points, poles, invalid logarithm domains, and
unsupported heads. These tests validate an executable exact-arithmetic
implementation of the theorem; they do not constitute a separate formal
verification of the Wolfram evaluator.

The package's two motivating functions
$x+x^2(1+\log x)$ and $x+x^{\sqrt2}$ are included, together with a nonunit
leading power and coefficient, both real Lambert branches, and exact-core
perturbations. A supplied interval immediately below $x=e^{-1}$ tests the
small derivative near the turning point of $x\log x$. There the reported
derivative lower bound is small but strictly positive in absolute value; the
result remains an interval theorem, without relying on a logarithmic
asymptotic truncation being accurate near that threshold.

\section{Implemented source charts and reconstruction checks}
\label{sec:source-chart-implementation}

The implementation of Section~\ref{sec:coordinates} lives in
\path{Kernel/SourceCoordinates.wl}. The two chart families
are distinguished by the function used to reconstruct the source:
\wl{"SourceExp"} reconstructs an exponential, while
\wl{"SourceLog"} reconstructs a logarithm. The underlying inverse always
has observable power $1$. Powers of the original inverse are formed only
after that chart variable has been computed.

\paragraph{Pure logarithmic source dependence.}
First choose the positive local source coordinate $u$ from the endpoint
and direction. Substitute $u=e^{-h}$, so $h\to+\infty$. A pure
logarithmic expression becomes a function $\Phi(h)$ with no residual
exponential dependence on $h$, and the existing inverse dispatcher solves
$\Phi(h)=y$. For a finite endpoint $x_0$ and side $s\in\{-1,1\}$,
\[
 x=x_0+s e^{-h};
\]
at the source infinity with sign $s$, the reconstruction is $x=s e^h$.
Thus the noninteger real powers $(-\log u)^q$, polynomial and rational
combinations of $\log u$, and both finite source directions can use the
same chart whenever the transformed inverse belongs to a supported scale.

For example, $F(u)=(\log u)^2+\log u$ becomes $h^2-h$, and its small
positive inverse begins
\[
 u=e^{-\sqrt y-1/2}
 \left(1-\frac1{8\sqrt y}+\OO(y^{-1})\right).
\]
The exponential factor multiplies the absolute remainder. A finite-order
approximation to $h$ is usable here only when its absolute error tends to
zero. A finite relative Lambert expansion for a growing $h$ need not have
this property; increasing its logarithmic depth alone does not resolve
that obstruction.

An exact transformed inverse may instead be kept as an exact carrier.
Repeated exact charts therefore give, for instance, the exact inverse
$u=\exp(-e^y)$ of $\log(-\log u)$. Recursion is bounded to avoid an
uncontrolled recognizer loop. Exactness is recorded separately: it can
come from an exact inner series, an exact inner inverse, or a bounded
symbolic check that a small candidate composes exactly to the target.
The asymptotic branch of the inner inverse fixes which exact solution is
being reconstructed.

\paragraph{Finite exponential families.}
At $x\to s\infty$, the implementation takes $u=e^{-s x}$, so
$x=-s\log u$. Finite sums of real exponential rates become finite
generalized-power germs in $u$. Equal rates and cancellations are handled
by the ordinary canonicalization before a leading term is selected.
Rates need not be commensurable. The recognizer requires a nonzero leading
algebraic exponent after removing a constant target offset; additional
logarithmic factors are permitted when the resulting finite power--log
germ is supported by the inner engine.

For example,
\[
 e^{-x}+e^{-2x}=y
 \quad\Longrightarrow\quad
 x=-\log y+y-\frac32y^2+\frac{10}3y^3+\OO(y^4),
 \qquad y\to0^+.
\]
The inverse of $e^x+e^{-x}$ at $+\infty$ begins
$\log y-y^{-2}$; selecting the source endpoint $-\infty$ gives its
negative. The chart preserves this distinction rather than choosing a
numerical root without an endpoint condition.

\paragraph{Original observables and term counts.}
At a finite source endpoint, the power option represents
$(x-x_0)^r$ when $r\ne1$, and the inverse itself when $r=1$.
At infinity it represents $x^r$. A negative source side therefore requires
integer $r$ for these real-valued observable conventions. Exponential
reconstruction multiplies $h$ by $\pm r$ before exponentiating; a finite
source offset is added separately when $r=1$.

For a non-polynomial power of a logarithmic reconstruction, the finite
approximation can be retained inside a specified observable carrier.
If $X\sim c|\log w|$ and $x-X=\OO(w^P\M^D)$ with $P>0$, the mean
value theorem gives
\[
 \frac{x^r-X^r}{X^r}
 =\OO\left(\frac{x-X}{X}\right)
 =\OO(w^P\M^D).
\]
This deliberately conservative relative bound makes the retained carrier
$X^r$ usable without claiming a finite polynomial--log expansion of its
fractional power. An explicit cutoff is required for that representation;
a unit-block term count would not count the corrections retained inside
the carrier. In the ordinary reconstructed unit, the term-goal option
counts the actual nonzero reconstructed blocks.

\paragraph{Residuals and numerical comparisons.}
The object records both directions of the chart and an expression that
recovers the chart variable from the requested observable. Formal
residual checking applies that expression to the \emph{returned finite
observable}. It does not reuse the unmodified inner inverse blocks after
an exponential or logarithmic reconstruction has introduced a new
truncation. The checked residual is
\[
 \frac{\Phi(h_{\mathrm{returned}})-y}{y-y_0},
\]
with denominator $y$ at an infinite target. If reconstruction requires a
coefficient algebra outside the supported jet scale, the exact residual
expression is still returned, accompanied by an explicit indication that
no jet-order assertion was made.

Numerical comparison solves $\Phi(h)=y$ at high precision, reconstructs
the source, checks its original side, and compares the requested observable
with the returned approximation. This avoids unnecessarily evaluating a
very large or very small original exponential during root finding. It is
numerical evidence, not an interval certificate. Refinement replays the
original inverse problem and reconstruction, preserving inverse provenance.
An externally supplied forward remainder must first be transported through
the source chart; such a declaration is rejected until that transport has
been specified.

\section{Finite logarithmic hierarchies and reciprocal-logarithmic order}
\label{sec:logarithmic-scales}

A power gap and a logarithmic gap describe different orders of smallness.
For example, the correction $u/\log u$ has the same source power as $u$,
whereas $u^2\sqrt{-\log u}$ has a positive source-power gap and a
nonpolynomial logarithmic coefficient. These two situations, and a
leading product of logarithmic factors, require different normalizations. Exact
source-coordinate reductions, including repeated logarithmic charts, are
handled by the coordinate machinery of Section~\ref{sec:coordinates}.

\subsection{The finite coefficient hierarchy}

In the positive source coordinate $u\to0^+$, put
\begin{equation}
 L_1(u)=-\log u,\qquad L_{j+1}(u)=\log L_j(u),
 \qquad 1\le j<d.
 \label{eq:logarithmic-hierarchy}
\end{equation}
The depth $d$ is finite. All retained levels are positive and tend to
infinity sufficiently near the endpoint. A generalized logarithmic
monomial is
\begin{equation}
 c\prod_{j=1}^d L_j^{\alpha_j},\qquad
 c\in\R,\quad \alpha_j\in\R.
 \label{eq:generalized-log-monomial}
\end{equation}
Finite sums of these monomials form an algebra under addition and
multiplication. Powers in \eqref{eq:generalized-log-monomial} have positive
bases, including when an exponent is noninteger or negative. An arbitrary
real power of a sum is not silently included in this finite algebra.

\begin{proposition}[Closure under the source Euler derivative]
\label{prop:logarithmic-euler-closure}
For a differentiable coefficient $C(L_1,\ldots,L_d)$,
\begin{equation}
 u\frac{d}{du}C(L_1(u),\ldots,L_d(u))
 =\mathcal E C,
 \qquad
 \mathcal E=-\sum_{j=1}^d
 \frac{1}{L_1\cdots L_{j-1}}\frac{\partial}{\partial L_j},
 \label{eq:hierarchy-euler-operator}
\end{equation}
where the empty denominator is $1$. In particular, the finite generalized
logarithmic algebra is closed under every derivative used in the inverse
coefficient formula.
\end{proposition}
\begin{proof}
The chain rule gives $uL_1'(u)=-1$ and
$uL_j'(u)=-(L_1\cdots L_{j-1})^{-1}$. Each summand in
\eqref{eq:hierarchy-euler-operator} decreases finitely many monomial
exponents and creates no new logarithmic level.
\end{proof}

This is the needed closure for coefficients such as $L_1^{1/2}$:
its derivatives contain $L_1^{-1/2},L_1^{-3/2},\ldots$.
Restricting the coefficient algebra to positive logarithmic powers would
lose closure even in this simplest example. An expression containing
$\sqrt{\log u}$ instead has an eventually negative base on $u\to0^+$;
the real hierarchy does not admit it.

\subsection{Positive source-power gaps with logarithmic coefficients}

Consider
\begin{equation}
 F(u)=y_0+a u^p\left(1+\sum_{i=1}^m
 u^{\delta_i}C_i(L_1(u),\ldots,L_d(u))\right),
 \qquad a\ne0,\quad p\ne0,\quad \delta_i>0,
 \label{eq:generalized-log-forward}
\end{equation}
where every $C_i$ belongs to the finite algebra above. All exponents are
fixed real constants, and the coefficients and the sign of $a$ are fixed
on the parameter domain under consideration.
Set $z=((y-y_0)/a)^{1/p}>0$, so $z\to0^+$. For the observable $u^r$, the
multi-index coefficient of weight $w=\bk\cdot\bdelta$ is
\begin{equation}
 \frac{(-1)^{|\bk|}r}{p^{|\bk|}\prod_i k_i!}
 \prod_{j=1}^{|\bk|-1}(\mathcal E+r+w+pj)
 \prod_i C_i^{k_i},
 \label{eq:generalized-log-coefficient}
\end{equation}
evaluated at the logarithmic levels of $z$ and multiplied by $z^{r+w}$.
For $\bk=0$, the coefficient is $1$. Formula
\eqref{eq:generalized-log-coefficient} is the same Lagrange identity as
in the polynomial case; Proposition~\ref{prop:logarithmic-euler-closure}
supplies its enlarged coefficient algebra.

\begin{theorem}[Convergence and a conservative complete-tail bound]
\label{thm:generalized-log-convergence}
The inverse of \eqref{eq:generalized-log-forward} on its selected eventual
real branch has the convergent multi-index expansion
\eqref{eq:generalized-log-coefficient}. Choose integers $D_i\ge0$ such
that every monomial in $C_i$ satisfies
\[
 \prod_j L_j(u)^{\alpha_j}=\OO(\M(u)^{D_i}).
\]
For example, the ceiling of the largest sum of the positive monomial
exponents is sufficient. Let $\mathcal B$ be the finite minimal excluded
multi-index boundary of a positive source-weight cutoff, and put
\begin{equation}
 \beta=\min_{\bk\in\mathcal B}\bk\cdot\bdelta,
 \qquad D=\max_{\bk\in\mathcal B}\sum_i k_iD_i.
 \label{eq:generalized-log-boundary}
\end{equation}
Then the complete omitted expansion of $u^r$ is
$\OO(z^{r+\beta}\M(z)^D)$. At a finite source endpoint the original
inverse has $r=1$; at an infinite source endpoint it has $r=-1$,
with the selected source sign applied separately.
\end{theorem}
\begin{proof}
Write $u=zW$ near $W=1$. For each level, the quotient
$L_j(zW)/L_j(z)$ is analytic in $W$ and the finite reciprocal logarithmic
parameters near their limiting values. For instance, the first quotient
is $1-\log W/L_1(z)$; subsequent quotients are obtained by taking the
logarithm of the preceding quotient and dividing by the next level.
Each factor raised to a fixed real power is holomorphic near $1$.
Introduce a separate small parameter for every finite coefficient monomial
$z^{\delta_i}\prod_jL_j(z)^{\alpha_j}$. All tend to zero. The normalized
implicit equation has derivative $p\ne0$ in $W$ at its parameter origin,
so the analytic implicit function theorem proves convergence in a fixed
polydisc of these parameters.

Cauchy estimates, with the finite coefficient sums grouped by $i$, give a
geometric multi-index majorant in
$K_i z^{\delta_i}\M(z)^{D_i}$ for suitable constants $K_i$.
The omitted upward-closed index set is covered by the translates of its
finite minimal boundary. Sum the majorant over each translate. Its
denominators remain bounded for sufficiently small $z$, and the boundary
monomial is bounded by $z^\beta\M(z)^D$ using
\eqref{eq:generalized-log-boundary}. Multiplication by $z^r$ completes
the estimate. The estimate does not assume that the first boundary
coefficient is nonzero.
\end{proof}

The resulting logarithmic degree can be larger than the optimal one. The
finite boundary supplies a valid envelope for the entire omitted sum,
while the retained coefficients remain exact generalized expressions.
The cutoff uses the ordinary positive target coordinate: a source weight
$w$ contributes target power $(r+w)/|p|$. This convention preserves the
nonunit leading-power and infinite-endpoint scaling.

\begin{example}[Nonpolynomial logarithmic coefficients]
For $T=-\log y$, inversion at $0^+$ gives
\begin{align}
 (u+u^2\sqrt{-\log u})^{-1}(y)
 &=y-y^2\sqrt T+y^3(2T-\tfrac12)+\OO(y^4\M(y)^3),
 \label{eq:inverse-sqrt-log-coefficient}\\
 (u+u^2\log(-\log u))^{-1}(y)
 &=y-y^2\log T
 +y^3\left(2\log^2T-\frac{\log T}{T}\right)
 +\OO(y^4\M(y)^3).
 \label{eq:inverse-iterated-log-coefficient}
\end{align}
Here the superscript $-1$ denotes function inversion. The displayed
remainders are conservative complete-tail bounds.
The derivative term in each cubic coefficient is essential and follows
directly from \eqref{eq:hierarchy-euler-operator}.
\end{example}

\subsection{Reciprocal-logarithmic units}

Now consider a core of the form
\begin{equation}
 F_0(u)=y_0+a u^p H(-1/\log u),
 \qquad p\ne0,\quad H(0)=1,
 \label{eq:reciprocal-log-core}
\end{equation}
where $H$ is a rational function analytic at $0$ with real coefficients.
Let $z=((y-y_0)/a)^{1/p}>0$, $t=-1/\log z>0$, and write $u=zW$.
The normalized inverse equation is exactly
\begin{equation}
 W^p H\left(\frac{t}{1-t\log W}\right)=1.
 \label{eq:reciprocal-log-unit-equation}
\end{equation}
Its derivative in $W$ is $p$ at $(t,W)=(0,1)$, so its real solution
$W(t)$ is analytic at $0$. Iterating
\begin{equation}
 W\longmapsto
 H\left(\frac{t}{1-t\log W}\right)^{-1/p}
 \label{eq:reciprocal-log-fixed-point}
\end{equation}
from $W=1$ determines successively correct Taylor coefficients: dependence
on the previous $W$ carries a positive power of $t$. Exact finite Taylor
arithmetic determines the finite bracket and its first omitted orders.
An exclusive cutoff $h$ retains exactly the integer degrees $n<h$.
Its remainder has the form $\OO(t^\beta)$, where $\beta$ is the first
unretained nonzero degree, or any smaller degree beyond the retained terms
when no sharper order has been established.

For the original source inverse, the bracket is multiplied by the exact
prefactor $\sigma z$ at a finite endpoint or $\sigma/z$ at infinity.
An observable uses the corresponding power of $W$ and transports the
prefactor and remainder together. This is a logarithmic cutoff inside
that prefactor, not a cutoff in $y$ or $1/y$.

\begin{example}[A zero source-power gap]
At $u\to0^+$, the inverse of $u+u/\log u$ begins
\begin{equation}
 y\left(1+t+t^2+2t^3+\frac72t^4+\OO(t^5)\right),
 \qquad t=-\frac1{\log y}.
 \label{eq:reciprocal-log-inverse-example}
\end{equation}
Equivalently, its bracket is
$1-1/\log y+1/\log^2y-2/\log^3y+7/(2\log^4y)$.
For the same forward expression at $+\infty$, use $t=1/\log y$ and the
bracket $1-t+t^2-2t^3+7t^4/2$.

The expression $u+u/\log u-u/(1+\log u)$ has a cancelled first
reciprocal-logarithmic correction. Its forward unit is
$1+t^2/(1-t)$ in the source logarithmic coordinate, and its inverse bracket
begins $1-t^2+\cdots$. Coefficients are merged before detecting the leading
unit, so a cancelled coefficient does not create a fictitious retained term.
\end{example}

Higher source-power terms can accompany \eqref{eq:reciprocal-log-core}.
If each has a strictly positive power gap and an admitted logarithmic
coefficient envelope, their induced inverse correction is smaller than
the prefactor times every fixed power of $t$. For
$u(1+1/\log u)+u^2$, the logarithmic bracket therefore agrees with
\eqref{eq:reciprocal-log-inverse-example} at every finite logarithmic
order. The omitted source-power expression nevertheless belongs to a
distinct, smaller asymptotic sector. A residual computed solely from the
logarithmic bracket refers to the leading core. Equality of those brackets
does not assert that the full inverses are equal.

\subsection{A leading logarithmic monomial}

The same normalization extends to
\begin{equation}
 F_0(u)=y_0+a u^p\prod_{j=1}^d L_j(u)^{q_j},
 \qquad p\ne0,\quad q_j\in\R.
 \label{eq:leading-hierarchy-core}
\end{equation}
Let $Y=(y-y_0)/a>0$ and define
\begin{equation}
\begin{gathered}
 T=-\frac{\log Y}{p},\quad t=\frac1T,\quad
 B_1=T,\quad B_{j+1}=\log B_j,\\
 d_0=\frac1p\sum_jq_j\log B_j,
 \quad z=Y^{1/p}\prod_jB_j^{-q_j/p}.
\end{gathered}
 \label{eq:leading-hierarchy-normalization}
\end{equation}
All $B_j$ are positive sufficiently near the selected endpoint. Put
\begin{align}
 R_1(W)&=1+t(d_0-\log W),\label{eq:leading-level-first-ratio}\\
 R_{j+1}(W)&=1+B_{j+1}^{-1}\log R_j(W).
 \label{eq:leading-level-next-ratio}
\end{align}
Then $L_j(zW)/B_j=R_j(W)$ exactly, and the normalized inverse equation is
\begin{equation}
 W^p\prod_jR_j(W)^{q_j}=1.
 \label{eq:leading-hierarchy-unit-equation}
\end{equation}

\begin{theorem}[A convergent leading-hierarchy bracket]
\label{thm:leading-hierarchy-bracket}
Equation~\eqref{eq:leading-hierarchy-unit-equation} has a unique real
solution near $W=1$ for the selected asymptotic branch. Its expansion in
powers of $t$, retaining the lower-level coefficients exactly, converges
sufficiently near the endpoint. A truncation beginning its omitted degrees
at $\beta$ has bracket remainder
\begin{equation}
 \OO\!\left(t^\beta(1+|\log t|)^\beta\right).
 \label{eq:leading-hierarchy-tail}
\end{equation}
\end{theorem}
\begin{proof}
Treat $t$, $td_0$, and $B_2^{-1},\ldots,B_d^{-1}$ as independent small
parameters. The equation is analytic in these parameters and $W$ near
their origin and $W=1$, with derivative $p\ne0$ in $W$.
The analytic implicit function theorem gives a common polydisc and a
convergent series there. Since $d_0=\OO(\log T)$, all the parameters
indeed tend to zero on the actual asymptotic curve.

The total degree in the first two parameters equals the formal degree in
$t$ after substituting $td_0$. Cauchy estimates, uniform in the other
parameters on a smaller polydisc, bound the tail of those degrees by a
constant times $[t(1+|d_0|)]^\beta$ for sufficiently small $t$.
This proves \eqref{eq:leading-hierarchy-tail}. The fixed-point expression
$W=\prod_jR_j(W)^{-q_j/p}$ has positive formal valuation in its dependence
on the previous $W$, so finite exact iteration computes the same
coefficients. Real powers use the positive local factors near $1$.
\end{proof}

For $F(x)=x\log\log x$ at $+\infty$, let
$T=\log y$, $B=\log T$, and $t=1/T$. The reconstructed inverse starts
\begin{equation}
 F^{-1}(y)=\frac{y}{B}
 \left(1+\frac{\log B}{B}t
 +\OO\!\left(t^2(1+|\log t|)^2\right)\right).
 \label{eq:leading-iterated-log-example}
\end{equation}
Thus an iterated logarithm in the leading coefficient need not be replaced
by a constant. Products of levels and real powers of those levels are
covered by the same positive local normalization. Positive source-power
perturbations again belong to separate smaller sectors.

\subsection{Precision, exactness, and the boundary of the hierarchy}

The three normalizations above have different precision meanings. Positive
source-power gaps lead to a cutoff in the positive target power coordinate.
Reciprocal-logarithmic units and leading logarithmic monomials instead use
a cutoff in the small logarithmic coordinate inside an exact prefactor.
All these cutoffs are exclusive. Terms at a common source weight must be
combined before counting nonzero blocks; cancellation can remove a complete
block without removing its contribution to the index boundary used for the
tail estimate.

A finite run of zero coefficients does not establish termination. A
sufficient exactness argument reconstructs a positive candidate source from
the finite observable and verifies the original forward equation identically.
If no smaller sector has been omitted, local uniqueness then identifies
this candidate with the inverse germ. In particular, a polynomial source
unit does not by itself make a truncated noninteger power of that unit exact.

The normalized equations also give direct residual identities. Substituting
a finite unit into \eqref{eq:reciprocal-log-unit-equation} or
\eqref{eq:leading-hierarchy-unit-equation} determines its defect in the
leading-core equation. Such a defect is not a residual for a different
forward function containing unretained source-power sectors.

The reciprocal-logarithmic unit has a single ordinary Taylor bracket.
The other normalizations retain an enlarged coefficient algebra, so their
composition requires additional closure arguments. The convergence and
tail theorems give existence of constants and neighborhoods; they do not
supply explicit numerical values for them. An unknown forward remainder
likewise requires a separate transport estimate through the selected
logarithmic hierarchy.

These constructions form a finite hierarchy rather than a general
transseries field. Arbitrary real powers of coefficient sums, sums of
incomparable leading logarithmic monomials, and independent simultaneous
cutoffs at every logarithmic level require further hypotheses and additional
notions of precision. The formulas above establish closure under the Euler
derivatives used in inversion, convergence for the stated families, and
complete-tail estimates with the distinct precision meanings made explicit.

\section{Finite commensurate flat sectors}
\label{sec:flat-sectors}

An exponential that is smaller than every power of the source coordinate
needs its own truncation index. A finite part of this larger scale arises
from the following exact forward expression in a positive source coordinate
$u\to0^+$:
\begin{equation}
 F(u)=y_0+a u^q+\sum_{j=1}^J A_j(u)e^{-k_j c/u^p},
 \quad p,c>0,\quad q\ne0,\quad a\ne0,\quad k_j\in\N_{>0},
 \label{eq:flat-forward}
\end{equation}
where each $A_j$ is a finite real power--log expression. The phase power and
rates are fixed real numbers, with all rates integer multiples of a common
positive rate. Coefficients may depend on fixed real parameters subject to
the stated sign assumptions. Different phase powers and incommensurate
rates require more than one exponential coordinate and fall outside this
single-phase construction.

The zero sector must be exact. An ordinary finite asymptotic approximation
to a nonmonomial core leaves a power-sized error, which dominates every
sector in \eqref{eq:flat-forward}. The exact shifted monomial core
$y_0+a u^q$ avoids this obstruction: its zero-sector inverse is exact. Put
\[
 u_0=\left(\frac{y-y_0}{a}\right)^{1/q}>0,
 \qquad E_0=e^{-c/u_0^p}.
\]
At a finite source endpoint the observable is
$H(u)=x_0+\sigma u$ for power one and $H(u)=\sigma^r u^r$ otherwise.
At an infinite endpoint it is $H(u)=\sigma^r u^{-r}$.
For negative source branches, take an integer observable power.

\subsection{Computing complete sectors}

Introduce a formal symbol $Z$ and the polynomial
$R(u,Z)=\sum_j A_j(u)Z^{k_j}$. Differentiation of the actual exponential
coordinate is represented by the exact operator
\[
 \mathcal D=\partial_u+cp\,u^{-p-1}Z\partial_Z,
 \qquad \frac{d}{du}e^{-c/u^p}=cp\,u^{-p-1}e^{-c/u^p}.
\]
The marker identity of Section~\ref{sec:core-perturbations} now gives
\begin{equation}
 \frac{(-1)^n}{n!}
 \left(\frac{\mathcal D}{a q u^{q-1}}\right)^{n-1}
 \left(\frac{H'(u)R(u,Z)^n}{a q u^{q-1}}\right).
 \label{eq:flat-marker-coefficients}
\end{equation}
Sum this expression over $1\le n\le N$, group equal powers of $Z$, and
evaluate $u=u_0$, $Z=E_0$. Since every $k_j\ge1$, no marker degree larger
than $N$ contributes to an exponential degree at most $N$. Also
$\mathcal D$ preserves exponential degree, so intermediate polynomial
truncation discards no contribution to a retained sector. Each resulting
coefficient is a complete finite power--log function of $u_0$.

For example, the inverse of $F(x)=x+e^{-1/x}$ at $0^+$ has the
first three exponential corrections
\begin{equation}
 y-E+y^{-2}E^2+
 \left(y^{-3}-\frac32 y^{-4}\right)E^3,
 \qquad E=e^{-1/y}.
 \label{eq:flat-primary-inverse}
\end{equation}
For $F(x)=x+x^2e^{-1/x}$, the first two sectors are
$y-y^2E+(y^2+2y^3)E^2$. Adding $2e^{-2/x}$ to the first example changes
the second coefficient to $y^{-2}-2$. Thus the sector index distinguishes
input sectors from powers of the perturbation marker.

\subsection{A bound for the complete omitted tail}

Write every coefficient as a finite sum $u^\alpha P(\log u)$, and choose
\[
\begin{aligned}
 b=\max(0,\max_{\alpha}(q+p-\alpha)),\qquad
 d&=\max(0,\max_P\deg P),\\
 B(u)=u^{-b}\M(u)^d,\qquad \chi(u)=B(u)e^{-c/u^p}.
\end{aligned}
\]
Let $r_*=r$ at a finite source endpoint and $r_*=-r$ at infinity.
For fixed admitted data, there are $C,K>0$ and a source threshold such that
the complete omitted sector tail after degree $N$ satisfies
\begin{equation}
 |\operatorname{Tail}_N|
 \le C u_0^{r_*+p}
 \frac{(K\chi(u_0))^{N+1}}{1-K\chi(u_0)}
 \quad\hbox{when }K\chi(u_0)<1.
 \label{eq:flat-complete-tail}
\end{equation}
Consequently its asymptotic class is
\[
 \OO\!\left(E_0^{N+1}
 u_0^{r_*+p-(N+1)b}\M(u_0)^{(N+1)d}\right).
\]
This bound is allowed to be conservative; it is independent of whether the
first omitted coefficient cancels.

To justify the estimate, set $u=u_0(1+u_0^p V)$ and divide the forward
equation by $u_0^{q+p}$. The normalized monomial difference has the analytic
extension
\[
 a\frac{(1+\epsilon V)^q-1}{\epsilon},
 \qquad \epsilon=u_0^p,
\]
whose derivative at $V=0$, $\epsilon=0$ is $a q\ne0$. Each exponential
factor, after extracting $E_0^{k_j}$, becomes
\[
 \exp\!\left(k_jc\,
 \frac{1-(1+\epsilon V)^{-p}}{\epsilon}\right),
\]
which is analytic and uniformly bounded for small $\epsilon,V$. The
power--log factors have the same local analytic property after their
powers of $u_0$ and $\log u_0$ are extracted. Each normalized coefficient
is bounded by a constant times $B(u_0)$. Since $B\ge1$ for small $u_0$,
replacing $E_0$ by a complex variable $Z$ gives a uniform implicit solution
on $|Z|<1/(K B(u_0))$. One may regard the bounded normalized coefficients
as additional parameters in a compact polydisc; the derivative condition
at $Z=0$ is uniform. Cauchy estimates therefore bound the coefficient of
$Z^k$ by $C(KB)^k$. The nonconstant observable contributes the factor
$u_0^{r_*+p}$. Summing the geometric tail proves
\eqref{eq:flat-complete-tail}, and $\chi(u_0)\to0$ establishes applicability
at $Z=E_0$.

The constants and threshold in \eqref{eq:flat-complete-tail} are
existential. They are not explicit numerical error bounds until suitable
values have been established. Sector depth is inclusive: a depth $N$
retains $k\le N$, whereas an ordinary power cutoff retains powers strictly
below its boundary. The first omitted coefficient and a bound for the
complete omitted tail are distinct mathematical objects.

The primary example also admits a direct formal check. Substitute
$g=y+\sum_{k=1}^N C_k(y)Z^k$ into
\[
 g+Z\exp(1/y-1/g)-y.
\]
Vanishing through degree $N$ in $Z$ identifies the same sector coefficients
by the uniqueness of the formal implicit solution. This identity verifies
the finite coefficients; the uniform analytic argument above supplies the
bound for the complete infinite tail.

\section{Finite Fourier coefficients in a logarithmic variable}
\label{sec:fourier-coefficients}

The polynomial coefficient algebra extends naturally to finite sums
\[
 B(L)=\sum_{\omega\in\Omega}P_\omega(L)e^{i\omega L},
 \qquad L=\log u,\qquad \Omega\subset\R\text{ finite},
\]
where every $P_\omega$ is a polynomial and the frequencies are exact real
numbers. This includes $\sin(\log u)$, several incommensurable frequencies,
and polynomial amplitudes multiplying those oscillations. The source powers
and the Fourier frequencies have different roles: source weights determine
asymptotic truncation, while frequencies label coefficient modes. No Fourier
mode is discarded merely because its frequency is large.

\subsection{Coefficient arithmetic and reality}

Multiplication convolves frequencies, and the Euler derivative is
\[
 \frac{\dd}{\dd L}\bigl(P_\omega(L)e^{i\omega L}\bigr)
 =\bigl(P_\omega'(L)+i\omega P_\omega(L)\bigr)e^{i\omega L}.
\]
Thus this algebra is closed under precisely the operations used by
Theorem~\ref{thm:coefficients}. Distinct syntactic frequencies are combined
when exact comparison proves them equal. Algebraic frequencies are reduced
before comparison. No floating-point tolerance is used to identify modes;
an undecidable comparison produces a failure.

The real-valued condition is the conjugacy relation
$P_{-\omega}(L)=\overline{P_\omega(L)}$ on real $L$. The parser checks it
under the supplied parameter assumptions. Intermediate complex coefficients
are permitted, and real trigonometric expressions are reconstructed for
display. Affine trigonometric phases, such as $\sin(2L+1)$, contribute their
constant phase factor to the polynomial amplitude.

\subsection{Inverse coefficients and independent composition}

The admitted forward family is
\[
 f(u)=y_0+a u^p\left(1+\sum_j u^{\delta_j}B_j(\log u)\right),
 \qquad ap\ne0,\quad \delta_j>0,
\]
with finitely many real Fourier-polynomial coefficients. The leading core is
a monomial with a proved nonzero real coefficient. An oscillatory leading
block, even one that happens to remain positive, requires a different inverse
normalization and is rejected explicitly by this engine.

For a multi-index $\bk$ of depth $n$ and weight
$\sigma=\bk\cdot\bdelta$, the unchanged coefficient formula is
\[
 C^{(r)}_{\bk}(L)=\frac{(-1)^n r}{p^n\bk!}
 \prod_{j=1}^{n-1}(\D+r+\sigma+pj)\prod_i B_i(L)^{k_i}.
\]
Convolution forms the product, and the Fourier Euler operation applies each
factor. The implementation preserves the individual coefficients as well as
the complete sum at a source weight.

Residual checking uses a separate composition recurrence. If
$u=z(1+U)$, then
\[
 (1+U)^\alpha B(L+\log(1+U))=\sum_{k\ge0}Q_k(L)U^k,
 \qquad Q_0=B,\qquad
 Q_{k+1}=\frac{(\alpha-k)Q_k+\D Q_k}{k+1}.
\]
This follows by Taylor expansion in the multiplicative increment; it holds
for each exponential-polynomial mode as well as for their sum. Its use in
the normalized forward equation independently checks the inverse coefficients.
For a nonunit observable power, the residual first reconstructs the source
unit by raising the observable unit to its reciprocal power.

\subsection{Envelopes, convergence, and omitted terms}

For real $L$, Fourier modes have modulus one. If
$P_\omega(L)=\sum_j c_{\omega,j}L^j$, then
\[
 |B(L)|\le\sum_{\omega,j}|c_{\omega,j}|(1+|L|)^j.
\]
This nonoscillatory envelope is stored with each retained block. Euler
differentiation preserves the polynomial degree; frequency factors affect
coefficient constants, rather than the power of the logarithmic envelope.
Consequently a multi-index coefficient has polynomial degree at most
$\sum_i k_i\deg B_i$, just as in the polynomial case.

The convergence argument also extends. After replacing each finite source
mode by an independent marker, logarithmic composition uses
\[
 e^{i\omega(L+\log W)}=e^{i\omega L}W^{i\omega},
\]
which is analytic for $W$ near $1$ with its local logarithm. The implicit
derivative at the origin is still $p\ne0$. The finite-dimensional analytic
implicit function theorem gives an absolutely convergent marker expansion
for sufficiently small markers. Along the real source slice these markers
are bounded by constants times $z^{\delta_i}(1+|\log z|)^{d_i}$ and tend to
zero. The positive source gaps also make the forward derivative asymptotic
to $ap u^{p-1}$, so the selected real inverse branch exists near the endpoint.

The complete one-step outer boundary of the retained multi-index region
therefore bounds the omitted tail. Its least weight supplies the remainder
power, and the maximum summed polynomial degree over that finite boundary
supplies a conservative logarithmic envelope. This bound need not be sharp
after cancellation, but remains valid even at zeros of every displayed
oscillatory coefficient. A declared forward remainder is transported with
the same observable-dependent exponent as Theorem~\ref{thm:transport}; its
matching derivative hypothesis remains required. Constants and source
thresholds are not computed numerical certificates.

\subsection{Interface and acceptance examples}

\begin{lstlisting}
s = AsymptoticFourierInverse[x + x^2 Sin[Log[x]],
  {x, 0}, {y, 4}];
FourierInverseResidual[s]
FourierInverseCoefficient[s, {2}]
\end{lstlisting}
The result is
\[
 y-y^2\sin L+y^3\sin L(2\sin L+\cos L)+\OO(y^4),
 \qquad L=\log y.
\]
With an additional $x^3\cos(2\log x)$, the coefficient of $y^3$ acquires
$-\cos(2L)$, combining a directly supplied frequency with frequencies
generated by convolution. The identity
\[
 \sin(\sqrt2L)\sin(\sqrt3L)
 +\tfrac12\cos\!\bigl(\sqrt{5+2\sqrt6}\,L\bigr)
 =\tfrac12\cos\!\bigl((\sqrt2-\sqrt3)L\bigr)
\]
tests exact algebraic frequency equality and cancellation.

The cutoff is exclusive in the positive target power coordinate, with the
same finite-endpoint, infinity, and observable conventions as the ordinary
engine. This first Fourier engine requires an explicit cutoff and uses the
individual Lagrange formula. Frequencies are exact numeric real constants;
arbitrary symbolic frequency relations and nonlinear logarithmic phases are
outside this interface. \wl{"MaxFrequencies"} bounds distinct modes across
the retained jet and intermediate coefficient operations, and
\wl{"MaxTerms"} bounds source indices and convolution work. Exhaustion
returns an explicit failure; there is no undocumented frequency truncation.

Regression cases compare the displayed formulas with direct trigonometric
differentiation, independently compose normalized residuals, check source and
target translations and an infinite source endpoint, and exercise observable
powers, algebraic resonance, polynomial envelopes, remainder transport, and
frequency budgets. The Fourier envelope test deliberately evaluates at a zero
of the oscillatory factor, where a sign-based or oscillatory error scale would
give a misleading result.

\section{Special functions, asymptotic forward models, and real thresholds}
\label{sec:special-function-adapters}

The operation \wl{AsymptoticSpecialInverse} exposes explicit adapters for
\wl{"Erfc"}, \wl{"LogGamma"}, \wl{"Gamma"},
\wl{"LambertThreshold"}, and \wl{"QuadraticThreshold"}.
These examples need two different contracts. The tail adapters invert a
finite Poincar\'e forward model and transport its omitted error. The
threshold adapters use exact local equations in a ramified coordinate.
A convergent inverse expansion for a finite model does not make the
original Poincar\'e forward series convergent.

Every adapter permits a real affine change of target,
$y=b+aF(x)$, with a provably real offset and a provably nonzero real scale.
It stores the original function, the forward model or exact local equation,
the transformed target, the selected real branch, and the meaning of the
requested cutoff. The following formulas use $a=1,b=0$ unless stated
otherwise.

\subsection{The complementary-error-function tail}

For $x>0$, define
\begin{equation}
 S(x)=\sqrt\pi\,x e^{x^2}\operatorname{erfc}x,
 \qquad
 S_m(x)=\sum_{k=0}^{m-1}\frac{(-1)^k(1/2)_k}{x^{2k}},
 \qquad A_m=(1/2)_m.
 \label{eq:erfc-normalized-tail}
\end{equation}
The real positive-axis expansion has a remainder bounded by its first
neglected term; see
\href{https://dlmf.nist.gov/7.12.i}{DLMF~\S7.12(i)}.
For inverse-error transport we also require a derivative bound.

\begin{lemma}[Value and derivative bounds for the normalized tail]
\label{lem:erfc-normalized-bounds}
For every integer $m\ge1$ and every $x>0$,
\begin{equation}
 |S(x)-S_m(x)|\le A_m x^{-2m},\qquad
 |S'(x)-S_m'(x)|\le 2m A_m x^{-2m-1}.
 \label{eq:erfc-normalized-bounds}
\end{equation}
\end{lemma}
\begin{proof}
Starting with the integral definition of $\operatorname{erfc}$ in
\href{https://dlmf.nist.gov/7.2.E2}{DLMF~(7.2.2)}, substitute
$s^2=x^2+t$ to obtain
\begin{equation}
 S(x)=\int_0^\infty e^{-t}(1+t/x^2)^{-1/2}\,dt.
 \label{eq:erfc-normalized-integral}
\end{equation}
Put $w=x^{-2}$. The derivatives of $(1+tw)^{-1/2}$ alternate in sign, and
their magnitudes for $w\ge0$ are at most their values at $w=0$.
Taylor's integral remainder, followed by integration in $t$, bounds the
remainder of the first $m$ terms by $(1/2)_m w^m$.
Applying the same argument to the derivative with respect to $w$ bounds
its remainder by $m(1/2)_m w^{m-1}$. For $m=1$, the latter statement is
the direct bound on the first derivative. Finally
$|dw/dx|=2x^{-3}$ gives the second inequality in
\eqref{eq:erfc-normalized-bounds}. The exponential factor in the integral
justifies these finite differentiations and integrations.
\end{proof}

Let $Q_m(w)$ be the Taylor polynomial of degree $m-1$ of
$-\log S_m(w^{-1/2})$ at $w=0$. The adapter constructs the finite phase
\begin{equation}
 \Phi_m(x)=x^2+\log x+Q_m(x^{-2}),
 \qquad v=-\log(\sqrt\pi\,y).
 \label{eq:erfc-finite-phase}
\end{equation}
The exact phase is $\Phi(x)=-\log(\sqrt\pi\operatorname{erfc}x)$.
For fixed $m$, Lemma~\ref{lem:erfc-normalized-bounds} implies
\begin{equation}
 \Phi(x)-\Phi_m(x)=\OO(x^{-2m}),\qquad
 \Phi'(x)-\Phi_m'(x)=\OO(x^{-2m-1}).
 \label{eq:erfc-phase-errors}
\end{equation}
These are value and derivative assertions; the second is not inferred
merely by differentiating an unspecified big-$O$ term.

The implementation records explicit conditional bounds as well. Write
$\delta=A_mx^{-2m}$, $\delta_1=2mA_mx^{-2m-1}$, and
$\widetilde\Phi_m=x^2+\log x-\log S_m$. Wherever $S_m>\delta$,
\begin{align}
 |\Phi-\Phi_m|
 &\le\frac{\delta}{S_m-\delta}
       +|\widetilde\Phi_m-\Phi_m|,
       \label{eq:erfc-explicit-phase-bound}\\
 |\Phi'-\Phi_m'|
 &\le\frac{\delta_1}{S_m-\delta}
 +\frac{|S_m'|\delta}{S_m(S_m-\delta)}
 +|\widetilde\Phi_m'-\Phi_m'|.
 \label{eq:erfc-explicit-phase-derivative-bound}
\end{align}
The terms involving $\widetilde\Phi_m$ are exact elementary expressions.
The positivity condition is retained in the result; it is satisfied
eventually because $S_m\to1$ and $\delta\to0$.

Now invert $\Phi_m(x)=v$ on $x\to+\infty$. Its leading source scale is
$x^2$, so the inverse error transported from \eqref{eq:erfc-phase-errors}
is $\OO(x^{-2m-1})=\OO(v^{-m-1/2})$. The ordinary inverse engine receives
the explicit local input remainder $\{2m,0\}$ in $u=1/x$. A requested
cutoff larger than $m+1/2$ is rejected unless the forward model is refined.
This bound is combined with the truncation error of the model inverse.

For example, with $v=-\log(\sqrt\pi y)$,
\begin{equation}
 \operatorname{erfc}^{-1}y
 =\sqrt v-\frac{\log v}{4\sqrt v}
 +\frac{-\log^2v+4\log v-8}{32v^{3/2}}
 +\OO\!\left(v^{-5/2}(1+|\log v|)^3\right).
 \label{eq:erfc-inverse-first-terms}
\end{equation}
The cubic numerator follows by inserting
$x=\sqrt v+A/\sqrt v+B/v^{3/2}$ into
$x^2+\log x+(2x^2)^{-1}+\OO(x^{-4})=v$.
The expansion has a power cutoff in $1/v$, although $y\to0^+$.

The negative-source tail uses the exact identity
$\operatorname{erfc}(-x)=2-\operatorname{erfc}x$ and changes the
reconstructed source sign. For $b+a\operatorname{erfc}x$ at $-\infty$,
the positive small tail is $2-(y-b)/a$ and the target endpoint is $b+2a$.
These identities handle reflected, shifted and sign-scaled tails without
changing the forward-error proof.

\subsection{A Stirling model with a retained Lambert inverse}

For $x>0$, let
\begin{equation}
 G_m(x)=x(\log x-1)-\frac12\log x+\frac12\log(2\pi)
 +\sum_{k=1}^{m-1}\frac{B_{2k}}{2k(2k-1)x^{2k-1}}.
 \label{eq:stirling-finite-model}
\end{equation}
The positive-axis Stirling and digamma error bounds in
\href{https://dlmf.nist.gov/5.11.ii}{DLMF~\S5.11(ii)} give
\begin{align}
 |\log\Gamma(x)-G_m(x)|&\le A_mx^{-\rho},
 \label{eq:stirling-value-bound}\\
 |\psi(x)-G_m'(x)|&\le\rho A_mx^{-\rho-1},
 \label{eq:stirling-derivative-bound}\\
 \rho&=2m-1,\qquad
 A_m=\frac{|B_{2m}|}{2m(2m-1)}.
 \notag
\end{align}
The finite value and derivative bounds apply for all positive $x$.
The adapter selects the increasing inverse branch $x>2$; on this interval
$\psi(x)\ge\log x-(2x)^{-1}-(12x^2)^{-1}>0$.
This branch choice excludes the separate behavior near the positive
minimum of the gamma function.

Use the exact core
\begin{equation}
 F_0(x)=x(\log x-1),\qquad
 \phi(v)=\frac{v}{W_0(v/e)}.
 \label{eq:stirling-exact-core}
\end{equation}
The difference $G_m-F_0$ is a finite higher-power perturbation in
$u=1/x$. The operation of Section~\ref{sec:core-perturbations} therefore
retains $\phi$ exactly and computes its complete marker corrections.
With $u_0=1/\phi(v)$, the omitted Stirling error transports to
\begin{equation}
 \OO(u_0^{2m-1})
 \label{eq:stirling-inverse-input-floor}
\end{equation}
in the implementation's conservative polynomial--log bound. The actual
derivative of the core provides an additional logarithmic denominator,
but discarding that improvement preserves a valid bound in the stated
remainder class. The final error combines
\eqref{eq:stirling-inverse-input-floor} with the complete marker tail.

For $m=1$, the first marker correction gives
\begin{equation}
 (\log\Gamma)^{-1}(v)
 =\phi(v)+\frac12-\frac{\log(2\pi)}{2\log\phi(v)}
 +\OO(\phi(v)^{-1}).
 \label{eq:stirling-first-corrected-inverse}
\end{equation}
Increasing marker depth while holding $m$ fixed cannot improve the inherited
Stirling accuracy limit. Marker coefficients beyond that limit still
describe the chosen finite model, with that interpretation explicit in
the result. Neither the finite marker expansion nor its analytic
complete-tail theorem establishes convergence of the infinite Stirling
series.

For inversion of $\Gamma(x)$, use the exact target transformation
$v=\log y$. The implementation solves and checks the equation
$\log\Gamma(x)=\log y$; it can therefore accept a target such as
$e^{10000}$ without forming that enormous value in the forward residual.
For a real affine target $b+a\Gamma(x)$, replace $y$ by $(y-b)/a>1$.
The analogous log-gamma adapter uses $v=(y-b)/a>0$ directly.

\subsection{The two real Lambert branches at their common threshold}

The two real inverses of $F(x)=xe^x$ meet at $x=-1$, $y=-1/e$.
The branch description is given in
\href{https://dlmf.nist.gov/4.13}{DLMF~\S4.13}.
Writing $x=-1+h$ gives the exact local equation
\begin{equation}
 eF(-1+h)+1
 =\frac{h^2}{2}+\frac{h^3}{3}+\frac{h^4}{8}
 +\frac{h^5}{30}+\cdots.
 \label{eq:Lambert-threshold-forward}
\end{equation}
The leading power is two. Ordinary analytic reversion in the positive
ramified coordinate
\begin{equation}
 t=\sqrt{2(1+ey)}
 \label{eq:Lambert-threshold-uniformizer}
\end{equation}
therefore yields a convergent local expansion. Its first coefficients are
\begin{equation}
 W_\nu(y)=-1+\sigma t-\frac{t^2}{3}
 +\sigma\frac{11t^3}{72}-\frac{43t^4}{540}+\OO(t^5),
 \label{eq:Lambert-threshold-inverse}
\end{equation}
where $(\nu,\sigma)=(0,1)$ or $(-1,-1)$.
The implementation computes these coefficients from the exact forward
equation, not from a large-argument Lambert asymptotic formula.

The option \wl{"LambertBranch"} selects $0$ or $-1$.
Its source direction must agree: the principal branch approaches $-1$
from above and the lower branch from below. Contradictory choices are
rejected. The recorded real target neighborhood is $-1/e<y<0$,
transformed by the requested affine target map. A target below the real
threshold is rejected by the numerical checker. The cutoff is an
exclusive power cutoff in the positive distance $y+1/e$, so a cutoff
$5/2$ retains the terms through $t^4$.

\subsection{An exact coalescing quadratic branch}

For an exact quadratic model
\begin{equation}
 F(x)=b+a c(x-x_0)^2,\qquad ac\ne0,
 \label{eq:quadratic-threshold-forward}
\end{equation}
the two branches are exactly
\begin{equation}
 x=x_0+\sigma\sqrt{\frac{y-b}{ac}},\qquad
 \sigma\in\{1,-1\},\qquad\frac{y-b}{ac}>0.
 \label{eq:quadratic-threshold-inverse}
\end{equation}
The source direction selects $\sigma$. The remainder is zero because this
is an exact inverse identity. For $F(x)=x^2-\mu$, the branches at target
$y=0$ are $\pm\sqrt\mu$ and coalesce as $\mu\to0^+$. The same coordinate
retains the selected sign throughout this limiting comparison. It does
not choose a different branch because the two roots become close.

\subsection{API, numerical evidence, and switching policy}

\begin{lstlisting}[language=Mathematica]
AsymptoticSpecialInverse["Erfc", {x, Infinity}, {y, 3}]
AsymptoticSpecialInverse["LogGamma", {x, Infinity}, {y, 3}]
AsymptoticSpecialInverse["Gamma", {x, Infinity}, {y, 2}]
AsymptoticSpecialInverse[
  "LambertThreshold", {x, -1}, {y, 5/2},
  "LambertBranch" -> -1]
AsymptoticSpecialInverse[
  "QuadraticThreshold", {x, 0}, {y, 1},
  "TargetOffset" -> -mu, Assumptions -> mu > 0,
  Direction -> "FromBelow"]
\end{lstlisting}

The gamma calls use an integer marker depth. The error-function call uses
a power cutoff in its logarithmic target coordinate. The threshold calls
use local target-power cutoffs. Each result identifies this distinction.
\wl{"ModelTerms"} controls the number of retained forward asymptotic
terms, with a default sufficient for the requested error-function cutoff
or gamma marker depth. A model that cannot support an error-function
cutoff is rejected; a deliberately coarse Stirling model retains its
explicit accuracy floor.

\wl{SpecialInverseNumericalCheck} uses the original special-function
equation or the exact threshold inverse as its independent reference.
Exact target substitutions precede numerical evaluation, preserving
$\log(e^v)=v$ and avoiding unnecessary cancellation in reflected tiny
error-function tails. If a numerical input has already lost the digits
needed by its transformed target, the checker reports insufficient
precision. Its output sets \wl{"Certified"} to \wl{False}: a
high-precision solve is not an interval proof.

The regression families compare successive model orders at common target
values, test the first-neglected value and derivative bounds against the
original special functions, compare local Lambert expansions with exact
\wl{ProductLog} values on both real branches, and use the exact quadratic
formula as an oracle. These are explicit overlap comparisons. They do not
establish a universal automatic crossover threshold between a local
expansion, a tail expansion and another numerical representation.

The current adapters are selected explicitly. Their metadata states this
switching policy. Complex branches, automatic certified crossover
selection, the gamma minimum and its second positive inverse branch,
arbitrary coalescing higher-order critical points, and additional special
functions remain outside these admitted families. The four roadmap
acceptance types are represented here with their own branch and error
contracts; their existence does not make these remaining cases automatic.

\section{Retaining an exact growing exponential core}
\label{sec:exponential-core-perturbation}

An inverse--logarithmic expansion of the Lambert function does not resolve
the much smaller correction caused by adding $x^2$ to $xe^x$. The operation
\wl{AsymptoticExponentialCoreInverse} retains the Lambert inverse exactly
and expands those corrections in a separate small exponential scale.
Its zero sector is an exact inverse, and each retained coefficient is a
complete function of that inverse.

\subsection{The admitted core and source charts}

Use a positive coordinate $v\to+\infty$ and consider the exact model
\begin{equation}
 F(v)=F_0(v)+R(v),\qquad
 F_0(v)=y_0+a v^b e^{c v^p},
 \qquad c,p>0,\quad a\ne0,
 \label{eq:exponential-core-forward}
\end{equation}
where $b$ is a fixed real number and $R$ is a finite real power--log
expression $\sum_\alpha v^\alpha P_\alpha(\log v)$.
Every such perturbation, including its finitely many derivatives, is
smaller than the growing core sufficiently near the endpoint.

At an infinite original source endpoint, the reconstruction is
$x=H(v)=s+\sigma v$ with $\sigma\in\{1,-1\}$. The source translation
$s$ can be supplied explicitly or recognized from affine factors in the
core. At a finite endpoint, use $x=H(v)=x_0+\sigma/v$.
Thus the same operation covers both $xe^x+x^2$ at infinity and
$e^{1/x}+x$ at $0^+$. Define $r_*=1$ for the first reconstruction and
$r_*=-1$ for the second.

Let $Y=(y-y_0)/a>0$. The selected exact core inverse is
\begin{equation}
 v_0=
 \begin{cases}
 (\log Y/c)^{1/p},&b=0,\\[0.3em]
 \displaystyle
 \left(\frac{b}{cp}
 W_\nu\!\left(\frac{cp}{b}Y^{p/b}\right)\right)^{1/p},&b\ne0,
 \end{cases}
 \quad
 \nu=\begin{cases}0,&b>0,\\-1,&b<0.\end{cases}
 \label{eq:exponential-core-exact-inverse}
\end{equation}
These are the real branches with $v_0\to+\infty$.
For $b<0$, the argument is negative and the lower branch selects the
increasing large-$v$ part of the core. The target domain records the real
Lambert argument condition, positivity of the real-power base, and
$b+cpv_0^p>0$. In all cases the exact inverse in the original source is
$\phi(y)=H(v_0(y))$.

An optional supplied \wl{"CoreInverse"} is retained only after a bounded
symbolic comparison establishes equality with this selected exact
expression. Additional terms inside the stated core are rejected if they
would make \eqref{eq:exponential-core-exact-inverse} merely an approximation.
They can instead be placed in the separately declared finite perturbation.

\subsection{Complete coefficients without expanding the core inverse}

Put
\begin{equation}
 D(v)=a v^{b-1}(b+cpv^p),\qquad
 F_0'(v)=e^{cv^p}D(v),\qquad E(v)=e^{-cv^p}.
 \label{eq:exponential-core-derivative}
\end{equation}
Introduce the marker $\lambda$ in $F_0(v)+\lambda R(v)=y$.
The exact-core Lagrange formula of Section~\ref{sec:core-perturbations}
shows that its degree-$n$ correction has the form
$C_n(v_0)E(v_0)^n$. The coefficient $C_n$ can be computed using only
power--logarithmic and rational operations on $v$:
\begin{align}
 Q_{n,1}(v)&=\frac{H'(v)R(v)^n}{D(v)},
 \label{eq:exponential-core-coefficient-initial}\\
 Q_{n,j+1}(v)&=
 \frac{Q_{n,j}'(v)-jcpv^{p-1}Q_{n,j}(v)}{D(v)},
 \qquad 1\le j<n,\label{eq:exponential-core-coefficient-recursion}\\
 C_n(v)&=\frac{(-1)^n}{n!}Q_{n,n}(v).
 \label{eq:exponential-core-coefficient-final}
\end{align}
Indeed, differentiating $Q(v)E(v)^j$ gives
$(Q'-jcpv^{p-1}Q)E^j$, and division by $F_0'$ adds one factor of $E$.
The implementation therefore avoids intermediate derivatives of
\wl{ProductLog} and avoids mixing the exact exponential factor into each
coefficient calculation.

The exact core identity gives a second, often more convenient expression
for the small scale:
\begin{equation}
 E_0=e^{-cv_0^p}=\frac{v_0^b}{Y}.
 \label{eq:exponential-core-target-sector}
\end{equation}
The result records both expressions and uses the right-hand side in the
retained formula. It is positive on the selected branch. Through inclusive
sector degree $N$, the approximation is
\begin{equation}
 \phi(y)+\sum_{n=1}^N C_n(v_0(y))E_0(y)^n.
 \label{eq:exponential-core-retained-sectors}
\end{equation}
Here one complete marker degree is also one exponential degree, because
the forward perturbation contains no additional exponentials.

\begin{example}[The inverse of $xe^x+x^2$]
Let $v_0=W_0(y)$ and $E_0=v_0/y$. The first coefficients are
\begin{align}
 C_1(v)&=-\frac{v^2}{v+1},
 \label{eq:xexpx-core-first-coefficient}\\
 C_2(v)&=\frac{v^3(-v^2+2v+4)}{2(v+1)^3}.
 \label{eq:xexpx-core-second-coefficient}
\end{align}
Thus the first correction to the exact core inverse is
\[
 (xe^x+x^2)^{-1}(y)
 =W_0(y)-\frac{W_0(y)^3}{(W_0(y)+1)y}+\cdots.
\]
This is much smaller than the error of any fixed inverse--logarithmic
truncation of $W_0(y)$. Retaining $W_0$ exactly makes that correction
meaningful without requiring a second, deeper core approximation.
\end{example}

\subsection{A uniform complete-tail theorem}

For the finite perturbation in \eqref{eq:exponential-core-forward}, choose
\begin{equation}
\begin{gathered}
 d=\max(0,\max_\alpha(\alpha-b)),\qquad
 k=\max_\alpha\deg P_\alpha,\\
 B(v)=v^d(1+|\log v|)^k,\qquad \chi(v)=B(v)e^{-cv^p}.
\end{gathered}
 \label{eq:exponential-core-envelope}
\end{equation}
The zero perturbation is treated separately and leaves the exact inverse
unchanged.

\begin{theorem}[Complete tail around a growing exact core]
\label{thm:exponential-core-complete-tail}
For fixed data in \eqref{eq:exponential-core-forward}, there exist positive
constants $C,K$ and a source threshold such that the selected real inverse
has a convergent expansion \eqref{eq:exponential-core-retained-sectors}
with its complete omitted tail bounded by
\begin{equation}
 |\operatorname{Tail}_N|
 \le C v_0^{r_*-p}
 \frac{(K\chi(v_0))^{N+1}}{1-K\chi(v_0)},
 \qquad K\chi(v_0)<1.
 \label{eq:exponential-core-geometric-tail}
\end{equation}
In particular, for fixed $N$ the remainder class is
\begin{equation}
 \OO\!\left(E_0^{N+1}
 v_0^{r_*-p+(N+1)d}(1+|\log v_0|)^{(N+1)k}\right).
 \label{eq:exponential-core-asymptotic-tail}
\end{equation}
The bound covers the whole omitted series independently of cancellation
in the first omitted coefficient.
\end{theorem}
\begin{proof}
Set
\[
 v=v_0(1+\epsilon V),\qquad\epsilon=v_0^{-p}.
\]
The displacement scale $v_0^{1-p}$ keeps the change in the exponential
phase bounded. Dividing the forward equation by $a v_0^b e^{cv_0^p}$
gives a core difference
\begin{equation}
 (1+\epsilon V)^b
 \exp\left(c\frac{(1+\epsilon V)^p-1}{\epsilon}\right)-1.
 \label{eq:exponential-core-normalized-equation}
\end{equation}
It extends analytically to $\epsilon=0$, and its derivative in $V$ at
$(\epsilon,V)=(0,0)$ is $cp\ne0$.

The normalized perturbation is
$E_0R(v_0(1+\epsilon V))/(a v_0^b)$.
After extraction of its powers of $v_0$ and $\log v_0$, each of its finite
coefficient factors is analytic and uniformly bounded near
$\epsilon=0,V=0$, with $1/\log v_0$ as an additional small parameter.
Its absolute envelope is a constant times $B(v_0)E_0$.
Replace $E_0$ temporarily by a complex variable $Z$ and rescale
$\zeta=B(v_0)Z$. The remaining normalized coefficient parameters range
over a bounded finite polydisc. Since the derivative $cp$ at $Z=0$ is
independent of them, the analytic implicit function theorem gives a
uniform solution near $V=0$ for $|\zeta|$ smaller than a fixed radius.
Equivalently, it is analytic on $|Z|<1/(K B(v_0))$.

For the observable, subtract its exact core value. The nonconstant part is
\[
 H(v_0(1+\epsilon V))-H(v_0)
 =\sigma v_0^{r_*-p}
 \frac{(1+\epsilon V)^{r_*}-1}{\epsilon}.
\]
The quotient is analytic near $(\epsilon,V)=(0,0)$ and uniformly bounded
on a smaller polydisc. Cauchy estimates bound its $Z^n$ coefficient by
$C(KB)^n$, with the displayed factor $v_0^{r_*-p}$.
The coefficient identity above identifies these coefficients with the
exact-core Lagrange coefficients. Summing the geometric tail and setting
$Z=E_0$ proves \eqref{eq:exponential-core-geometric-tail}.
Finally $\chi(v_0)\to0$, and the eventual derivative of the real forward
function has the core's nonzero sign, so this solution is the selected
real inverse near the endpoint.
\end{proof}

For $xe^x+x^2$, one can take $d=1,k=0,p=1,r_*=1$.
The complete tail after sector $N$ is consequently
$\OO((v_0E_0)^{N+1})$.
For $e^{1/x}+x$ at $0^+$, the reconstruction has $r_*=-1$,
and the first correction is
\begin{equation}
 (e^{1/x}+x)^{-1}(y)
 =\frac1{\log y}+\frac1{y\log^3y}
 +\OO\left(\frac1{y^2\log^2y}\right).
 \label{eq:finite-endpoint-exponential-core-example}
\end{equation}
The last bound is conservative but follows directly from
\eqref{eq:exponential-core-asymptotic-tail} with $d=0$.

\subsection{An unknown input error remains a separate first sector}

Suppose the actual perturbation includes a real continuously differentiable
unknown term $A(v)$ satisfying
\begin{equation}
 A(v)=\OO(v^{-\rho}(1+|\log v|)^\ell),\qquad
 A'(v)=\OO(v^{-\rho-1}(1+|\log v|)^\ell).
 \label{eq:exponential-core-input-error}
\end{equation}
Here $\rho$ may be any fixed real number: the growing exponential core
dominates every such finite power scale. Its derivative is asymptotic to
$acp\,v^{b+p-1}e^{cv^p}$. The mean value transport argument, with the
derivative of $H$, gives an additional inverse error
\begin{equation}
 \OO\!\left(E_0v_0^{r_*-b-p-\rho}
                  (1+|\log v_0|)^\ell\right).
 \label{eq:exponential-core-input-transport}
\end{equation}
The matching derivative assumption ensures eventual monotonicity and
permits comparison on the same local inverse branch. The model inverse
and the actual inverse are both asymptotic to the exact core coordinate,
so the scale can be expressed in $v_0$.

Equation~\eqref{eq:exponential-core-input-transport} has exponential degree
one, even when the finite model has been expanded through a much larger
sector depth. The implementation therefore stores it separately and adds
it to the model-tail remainder. It does not relabel it as a higher sector
or claim that additional exact model coefficients overcome unknown input
data. The retained terms beyond this error still describe the explicitly
specified finite model.

\subsection{Implemented behavior and verification}

\begin{lstlisting}[language=Mathematica]
AsymptoticExponentialCoreInverse[
  x Exp[x], x^2, {x, Infinity}, {y, 2}]

AsymptoticExponentialCoreInverse[
  7 + 3 (x - 2) Exp[x - 2], 5 (x - 2)^2,
  {x, Infinity}, {y, 1}]

AsymptoticExponentialCoreInverse[
  Exp[1/(2 - x)], 2 - x, {x, 2}, {y, 1},
  Direction -> "FromBelow"]
\end{lstlisting}

The result identifies scale \wl{"ExactExponentialCoreSectors"}.
It records the exact inverse and branch certificate, the divergent source
coordinate and reconstruction, the complete sector coefficients, the first
omitted coefficient, both representations of $E_0$, and the separate model
and input-error remainders. Its \wl{"MajorantContract"} is an asymptotic
existence theorem, with \wl{"NumericCertificate"} set to \wl{False}; the
constants and uniform threshold are not computed numerical bounds.

Independent regression checks insert
$v+\sum_{n=1}^N C_n(v)Z^n$ into
$X e^{X-v}-v+ZX^2=0$ for the primary example and perform ordinary Taylor
expansion in $Z$. Further checks use original forward targets generated
from known source points, compare successive sector depths at high
precision, and exercise affine shifts and scales, negative source infinity,
the lower real Lambert branch, nonunit phase powers and finite endpoint
reconstruction. Rejection cases cover a decaying core, an additional term
incorrectly placed in the exact core, an extra exponential rank in the
perturbation, an unverified inverse, and unproved coefficient reality.

The admitted phase is a single growing monomial $cv^p$. General phase
polynomials would need an appropriate exact inverse supplied with its own
branch contract, while perturbations such as $e^{v/2}$ need more than the
single exponential sector used here. The operation does not replace those
missing contracts with a finite inverse--logarithmic core approximation.

\section{Separate inner precision and calculus of flat sectors}
\label{sec:flat-sector-operations}

A finite calculus for the scale of Section~\ref{sec:flat-sectors}
uses a common positive monomial target coordinate
\[
 u=\left(\frac{y-y_0}{a}\right)^{1/q},\qquad
 E=e^{-c/u^p},\qquad p,c>0,
\]
for every operand. The phase and chart are part of the scale; combining
different phases requires an additional change-of-scale argument.

The asymptotic representation distinguishes two errors:
\begin{equation}
 S(y)=C_0(u)+\sum_{k=1}^{N}E^k
 \left(T_k(u)+\OO(u^{\rho_k}\M(u)^{d_k})\right)
 +\OO(E^{N+1}u^{\rho_*}\M(u)^{d_*}).
 \label{eq:flat-two-precision}
\end{equation}
Here $C_0$ is exact, each $T_k$ is a finite power--log expression, and
$\rho_k=+\infty$ means that the coefficient is exact. The remainder
is a sum with one term for each nonzero inner error and a separate term
for the complete omitted exponential tail. These are asymptotic bounds for
fixed data, whose constants and threshold are existential.

\subsection{Inner truncation and multiplication}

An inner truncation at $h$ retains powers strictly below $h$ in every
positive sector. The zero sector remains exact, including any of its powers
at or above $h$. If a coefficient loses its first term at power $\alpha$,
its inner error has power $\alpha$ and the degree of the corresponding
logarithmic polynomial. An existing, less precise bound is preserved.
Increasing the formal cutoff after truncation does not recover discarded
coefficient information.

For example, truncating the inner coefficients of
\[
 S=y-y^2E+(y^2+2y^3)E^2+\OO(y^2E^3),\qquad E=e^{-1/y},
\]
at $h=3$ gives
\[
 S=y-y^2E+y^2E^2+E^2\OO(y^3)+\OO(y^2E^3).
\]
The error $E^2\OO(y^3)$ cannot be put into the third sector: its ratio to
any $E^3$ times a fixed power--log expression is unbounded as $y\to0^+$.

Multiplication convolves exponential degrees and uses the ordinary
precision-tracked power--log multiplication inside each degree. If
$A=\sum_{i\le N_A}E^i A_i+R_A$ and
$B=\sum_{j\le N_B}E^j B_j+R_B$, the result retains degrees through
$N=\min(N_A,N_B)$. Its tail includes all three kinds of contributions:
\[
 \sum_{i+j>N}E^{i+j}A_iB_j,
 \qquad R_A\sum_j E^j B_j+R_B\sum_i E^i A_i,
 \qquad R_A R_B.
\]
Each coefficient bound includes its own inner uncertainty. Since $0<E<1$,
a term in degree $k>N$ can conservatively be bounded in degree $N+1$
while retaining its power--log envelope. This can give a looser algebraic
power in the final sector bound; it is always a valid bound. In particular,
unknown tails are added by bounds and never cancelled because some displayed
coefficients cancel.

Polynomial composition $H(S)$ follows from addition and multiplication.
The coefficients of $H$ may be exact finite real power--log expressions
in the target chart. Horner's formula computes the composition through
successive products and gives the exact zero sector $H(C_0)$. A constant
observable has no dependence on the unknown input tail; the zero observable
therefore has zero remainder. General nonpolynomial composition and changes
of phase require additional analytic and scale-compatibility arguments.

\subsection{Qualified differentiation}

A magnitude bound alone does not justify its differentiated bound.
For the exact flat inverse, the uniform holomorphic implicit-function
construction supplies bounds for every fixed derivative order. The exact
finite forward model is essential to this argument. Addition,
multiplication, polynomial composition, and inner truncation preserve the
required local holomorphic estimates. A general real-axis remainder has
no such derivative bound without a separate regularity hypothesis.

For a retained monomial the target derivative is explicit:
\begin{equation}
\begin{aligned}
 \frac{d}{dy}\left(E^k u^\alpha P(\log u)\right)
 =\frac{E^k}{a q}\Bigl[&u^{\alpha-q}
       \bigl(\alpha P(\log u)+P'(\log u)\bigr)\\
       &+kcp\,u^{\alpha-p-q}P(\log u)\Bigr].
\end{aligned}
 \label{eq:flat-target-derivative}
\end{equation}
Thus the exponential degree stays $k$. For $k>0$, a qualified inner
bound at $(\rho_k,d_k)$ is transported to
$(\rho_k-p-q,d_k)$. When known differentiated terms are then discarded at
that same boundary power, their logarithmic degree can increase the final
inner degree bound. The complete tail has the same power shift
$\rho_*\mapsto\rho_*-p-q$. This formula applies also when $q<0$;
the shift is a signed exact number, not its absolute value.

For completeness, the analytic tail estimate supports this derivative
claim. Fix a small positive real $u$ and use a complex disk
\[
 |v-u|<\eta u^{p+1},
\]
with fixed sufficiently small $\eta>0$. On this disk the chosen local
power and logarithm branches are holomorphic, their power--log envelopes
are comparable to those at $u$, and
\[
 v^{-p}-u^{-p}=\OO(1).
\]
Consequently $e^{-c/v^p}$ differs from $e^{-c/u^p}$ by a bounded factor.
The normalized implicit-function construction in
Section~\ref{sec:flat-sectors} extends uniformly to these disks. Its
geometric tail estimate therefore holds there with changed constants.
Cauchy's estimate costs a factor $u^{-p-1}$ for one source derivative.
The chain-rule factor $du/dy=(a q)^{-1}u^{1-q}$ gives the asserted loss
$u^{-p-q}$ for one target derivative. Every fixed higher derivative is
handled on nested disks, or by the corresponding target disk of radius
comparable to $u^{p+q}$; constants may depend on the derivative order.
No differentiation of an arbitrary real-axis Big-O bound is used.

The finitely many discarded inner terms are themselves analytic
power--log functions, so the same bound, and the exact phase derivative,
apply to their errors. After differentiation, terms are retained only within the transported precision. For the previous example, the derivative
of the truncated representation has the valid form
\[
 S'=1-(1+2y)E+2E^2+E^2\OO(y)+\OO(E^3),
\]
where the last bound is the transported $\OO(y^2E^3)$.
The full second-sector derivative contains further powers of $y$ that the
truncated data do not determine.

The two levels of precision remain distinct under all these operations.
A power-sized uncertainty in sector $k$ cannot be replaced by a tail in
sector $k+1$, even if many known coefficient terms cancel. The analytic
regularity required for differentiation is likewise separate from the
magnitude of the remainder. These distinctions prevent formal coefficient
arithmetic from asserting precision not present in the original data.

\section{Composition and differentiation in a reciprocal logarithm}
\label{sec:reciprocal-log-operations}

The reciprocal-logarithmic unit of Section~\ref{sec:logarithmic-scales}
has more regularity than an arbitrary function satisfying a magnitude
Big-$O$ estimate. For an exact rational logarithmic forward unit, its
implicit inverse is holomorphic in the small reciprocal logarithm. This
regularity supplies composition and differentiation rules without
discarding the exact power carrier or asking for an unsupported derivative
of an unknown remainder.

Consider exact rational logarithmic forward units with positive monomial
inverse carriers and zero finite source and target offsets. Known positive
source-power perturbations are not included in this particular analytic
construction: although smaller than every reciprocal-logarithmic order,
they require a separate differentiable-remainder argument. The composition
and differentiation below preserve holomorphy of the coefficient function.

\subsection{A common positive logarithmic coordinate}

Write
\begin{equation}
 t=\frac{\epsilon}{\log y}>0,\qquad
 \epsilon=\begin{cases}-1,&y\to0^+,\\1,&y\to+\infty.\end{cases}
 \qquad g(y)=C y^\alpha A(t),
 \label{eq:reciprocal-calculus-carrier}
\end{equation}
Here $C>0$, $\alpha$ is a fixed real number, and $A$ is holomorphic
near $t=0$. For an original inverse unit, $A(0)=1$. Differentiation can
produce a coefficient function whose constant vanishes or has either sign;
these outcomes remain valid carriers for further differentiation and as
outer functions in composition.

Suppose the original inverse has positive leading amplitude $a$ and
internal leading source power $p\ne0$. Its small coordinate is
$\tau=-1/\log((y/a)^{1/p})$, and therefore
\begin{equation}
 \tau=\frac{|p|t}{1-\epsilon(\log a)t},\qquad
 \epsilon=-\operatorname{sign}p.
 \label{eq:reciprocal-calculus-rechart}
\end{equation}
This is an analytic change of variable at zero with positive derivative.
It transports an analytic $\OO(\tau^\beta)$ tail to
$\OO(t^\beta)$. For an internal observable power $r_*$, the monomial
carrier has $\alpha=r_*/p$ and $C=\sigma^r a^{-\alpha}$.
Assume this $C$ is positive. In addition to the
original target domain, the canonical chart requires $y>0$ and
$\epsilon\log y>0$. These are eventual branch conditions; they are not
a computed radius of convergence.

The Taylor expansions have nonnegative integer powers of $t$ and constant real
coefficients. A finite exclusive cutoff $h$ retains powers $n<h$; the
error exponent is limited by the inherited Taylor remainder and any
discarded coefficient. Cancellation can make the first omitted nonzero
degree larger than the cutoff.

\subsection{Composition of two monomial carriers}

Let the outer function and inner function be
\[
 g(v)=C_o v^\alpha A\!\left(\frac{\epsilon_o}{\log v}\right),
 \qquad h(y)=C_i y^b V(t),\qquad
 t=\frac{\epsilon_i}{\log y},\quad V(0)=1.
\]
A positive nonunit constant in the inner coefficient function is first
absorbed into $C_i$. The composition requires $b\ne0$ and
\begin{equation}
 \epsilon_o b\epsilon_i>0.
 \label{eq:reciprocal-composition-endpoints}
\end{equation}
This condition says that the positive inner function approaches the outer
function's selected endpoint. It permits a negative power to interchange
zero and infinity when the two specified endpoints agree with that change.

Taking the real logarithm of the inner carrier gives the exact identities
\begin{align}
 T(t)&=\frac{\epsilon_o t}
 {b\epsilon_i+t(\log C_i+\log V(t))},
 \label{eq:reciprocal-composed-coordinate}\\
 g(h(y))&=C_o C_i^\alpha y^{\alpha b}
       V(t)^\alpha A(T(t)).
 \label{eq:reciprocal-composition-identity}
\end{align}
The denominator at zero is nonzero, $T(0)=0$, and
$T'(0)=\epsilon_o/(b\epsilon_i)>0$. The functions $\log V$ and
$V^\alpha$ use their analytic branches with $V(0)=1$. Consequently the
new coefficient function is holomorphic near zero and the original
monomial carrier remains exact.

If the outer and inner Taylor remainders are
$\OO(T^{\beta_o})$ and $\OO(t^{\beta_i})$, respectively, a conservative
relative remainder for \eqref{eq:reciprocal-composition-identity} is
\begin{equation}
 \OO\!\left(t^{\min(\beta_o,\beta_i)}\right).
 \label{eq:reciprocal-composition-error}
\end{equation}
To see this, an inner perturbation of size $\OO(t^{\beta_i})$ changes
$\log V$ and $V^\alpha$ by that order, while it changes $T$ by
$\OO(t^{\beta_i+2})$. The outer function has bounded derivatives near
zero, and its original remainder is evaluated at a coordinate comparable
to $t$. The Taylor algebra preserves these bounds and can improve them when
a constant factor or derivative vanishes.
Truncating the displayed polynomial adds its own first discarded order.
A larger formal cutoff never removes either operand's unknown tail.

For example, if both functions have the form $yW(-1/\log y)$ near zero,
the coefficient function of their composition is
\[
 W(t)\,W\!\left(\frac{t}{1-t\log W(t)}\right).
\]
At positive infinity the denominator is instead $1+t\log W(t)$.
These formulas determine the coefficients by ordinary Taylor expansion.
For general powers and scales, \eqref{eq:reciprocal-composed-coordinate}
retains both the exponent factor $b$ and the additive logarithm
$\log C_i$; dropping either gives incorrect coefficients already at low
reciprocal-logarithmic orders.

An inner coefficient function with a zero constant term is excluded from
this composition rule. Its logarithm introduces an additional $\log t$
scale. A negative constant also fails the positive inner branch condition.
Those failures describe a change in the required scale or branch, rather
than a reason to take a logarithm of an unproved positive approximation.

\subsection{Differentiating the analytic remainder}

The coordinate derivative is
\[
 \frac{dt}{dy}=-\frac{\epsilon t^2}{y}.
\]
Hence
\begin{equation}
 \frac{d}{dy}\bigl(Cy^\alpha A(t)\bigr)
 =Cy^{\alpha-1}\left(\alpha A(t)-\epsilon t^2 A'(t)\right).
 \label{eq:reciprocal-carrier-derivative}
\end{equation}
For the $n$th derivative, the coefficient function is obtained by applying
the successive operators
\begin{equation}
 \prod_{k=0}^{n-1}
 \left(\alpha-k-\epsilon t^2\frac{d}{dt}\right)
 \quad\hbox{to }A(t),
 \label{eq:reciprocal-carrier-higher-derivatives}
\end{equation}
and the carrier is $Cy^{\alpha-n}$. The factors differ only by constants
and commute. Applying them successively displays the precision change
and possible cancellation at each derivative order.

If a holomorphic remainder has a zero of order at least $\beta$ at zero,
then its $j$th derivative is $\OO(t^{\beta-j})$ for every fixed $j$.
This follows directly from its convergent Taylor series, or from Cauchy
estimates on a smaller disk. Thus the $t^2 A'$ contribution in
\eqref{eq:reciprocal-carrier-derivative} has remainder
$\OO(t^{\beta+1})$, while a nonzero $\alpha A$ contribution has remainder
$\OO(t^\beta)$. If $\alpha=0$, the latter contribution disappears and the
relative logarithmic error improves by one power. This is a proved
derivative contract for the exact rational-logarithmic implicit model;
it is not a rule for differentiating an arbitrary magnitude Big-$O$.

For the inverse of $x+x/\log x$ near zero, the retained coefficient
function begins
\[
 A(t)=1+t+t^2+2t^3+\tfrac72t^4+\OO(t^5).
\]
The first derivative therefore has the expansion
\[
 1+t+2t^2+4t^3+\tfrac{19}{2}t^4+\OO(t^5).
\]
Its carrier is constant. Differentiating once more gives a carrier $y^{-1}$
and a relative remainder $\OO(t^6)$, provided the known terms generated
through degree five are retained. This illustrates the precision improvement
when a carrier derivative vanishes. At infinity the signs of the odd
powers and the coordinate derivative change consistently.

\subsection{Refinement and the analytic hypothesis}

A stronger approximation to a composed or differentiated function requires
additional coefficients of its underlying analytic functions. Merely
truncating the same finite polynomial at a larger formal degree supplies
no new information about its unknown tail. Once those extra coefficients
are known, the exact substitution formulas and derivative operators above
transport their precision to the result.

The remainder theorem depends on holomorphy in the reciprocal logarithm,
rather than on finite coefficient identities alone. A source-power sector
which is smaller than every power of $t$ can be invisible to the Taylor
series yet obstruct holomorphy at $t=0$. Its differentiation therefore
requires a separate sector estimate. The hypotheses identifying an exact
rational logarithmic model are part of the mathematical statement and
cannot be discarded after its finite Taylor coefficients have been found.

As with the underlying implicit-function theorem, the convergence radius
and bounding constants exist for fixed data. The formulas in this section
do not assign explicit numerical values to those constants.

\section{Implicit inverse composition and real branch selection}
\label{sec:inverse-function-expressions}

An expression $F^{-1}(h(x))$ poses two distinct mathematical problems:
selecting the inverse branch of $F$, and expanding its composition with
$h$. Branch selection must precede the coefficient calculation. A real
inverse germ consists of a defining equation, a source domain, a source
endpoint, and an approach direction. These data distinguish solutions
which may share the same limiting target.

\subsection{Scalar implicit equations and parameters}

For a function $F$ of $n$ real variables, inversion in its $k$th argument
with the other arguments fixed means solving the scalar equation
\begin{equation}
 F(a_1,\ldots,a_{k-1},t,a_{k+1},\ldots,a_n)=a_k.
 \label{eq:inverse-expression-selected-argument}
\end{equation}
This is a one-dimensional inverse problem with parameters. For example,
$a t+t^2=y$, with fixed $a>0$ and $t\ge0$, has the positive local solution
\[
 t=\frac{\sqrt{a^2+4y}-a}{2}
   =\frac ya-\frac{y^2}{a^3}+\frac{2y^3}{a^5}
     +\OO(y^4),\qquad y\to0^+.
\]
The parameter $a$ is fixed in this statement. If it varies with $y$, the
error bound and even the relevant leading balance can change.

An important family with varying coefficients admits the exact reduction
\[
 A(x)F(t)+B(x)=h(x)
 \quad\Longleftrightarrow\quad
 F(t)=\frac{h(x)-B(x)}{A(x)}.
\]
Assume $A(x)\ne0$ on the selected deleted neighborhood and that the
source domain for $t$ is independent of $x$. Then the variation of $A$
and $B$ belongs entirely to the target argument. Nonvanishing at the
endpoint itself is unnecessary: $A(x)=x$ is permitted near $0^+$.
The quotient must nevertheless be expanded with its full uncertainty;
replacing a varying coefficient by its limiting value is generally invalid.

The coupled family $a(x)t+t^2=h(x)$ cannot in general be reduced to a
fixed forward germ by this separation. Joint implicit asymptotics would
require hypotheses on the relative rates and on uniform regularity in the
varying parameters. Such an additional analysis is independent of the
fixed-parameter inversion formulas.

\subsection{Source, target, and parameter domains}

Let $C(t)$ specify a source restriction. The admissible real source domain
is its intersection with the real domain of the chosen forward expression
$F(t)$, under the fixed parameter assumptions. Merely requiring $t$ real
does not imply $t>0$, whereas $t+t^2(1+\log t)$, with the real logarithm,
is defined only for $t>0$. Domain restrictions remain part of the inverse
problem throughout the calculation.

A restriction on the expansion variable $x$ instead specifies the approach
in $F^{-1}(h(x))$. A condition on a fixed parameter specifies the family
for which the conclusion holds. These three kinds of conditions have
different roles: an approach condition must hold eventually along $x$,
a source condition must hold on the selected inverse germ, and a parameter
condition is fixed throughout the limiting argument.

\subsection{Local existence and unambiguous selection}

Let $D\subset\R$ denote the admitted source domain. A source endpoint
$t_*$ is approached through a positive coordinate $u\to0^+$:
\[
 t=t_*+s u\quad(t_*\in\R),\qquad
 t=\frac{s}{u}\quad(t_*=s\infty),\qquad s\in\{-1,1\}.
\]
The condition $t(u)\in D$ concerns a \emph{deleted} neighborhood
$0<u<\delta$. It need not hold at $u=0$, where the source expression
may be undefined. Thus a restriction $0<t<1$ is compatible with
$t\to0^+$.

\begin{proposition}[A selected real inverse germ]
\label{prop:inverse-expression-local-branch}
Suppose there is a deleted source neighborhood on which $F$ is real,
continuous, and differentiable, and $F'$ has a fixed strict sign.
Suppose also that $F(t(u))\to\eta$, with the prescribed one-sided target
approach when $\eta$ is finite. Then a sufficiently small interval on
that target side, or a sufficiently distant target tail when
$\eta=\pm\infty$, has a unique inverse with values in this source
neighborhood. It approaches $t_*$ on the selected source side.
\end{proposition}

\begin{proof}
The source chart parametrizes an interval. The derivative condition
makes $F$ strictly monotone there, and continuity makes its image an
interval. Its endpoint limit and target side show that the image
contains a sufficiently small target interval or tail of the stated
kind. The inverse on that image is unique and continuous, and its
endpoint behavior follows from strict monotonicity.
\end{proof}

A nonzero leading block of the derivative's real asymptotic expansion
can prove its eventual sign. The leading coefficient and the parity of
the leading logarithmic degree determine that sign because
$\log u\to-\infty$. This is a qualitative eventual-sign proof; it does
not compute a numerical radius on which the derivative bound holds.
A neighborhood known to satisfy the source condition is therefore not
automatically a quantitative neighborhood for the full asymptotic error bound.

Proposition~\ref{prop:inverse-expression-local-branch} verifies a proposed
branch. \emph{Selection} among different source endpoints
requires an additional argument. There are two useful sufficient routes.

\begin{enumerate}[leftmargin=*]
\item \textbf{Complete real fiber and boundary information.}
For a finite target $\eta$, all interior points satisfying $F(t)=\eta$
must be accounted for, together with all finite boundaries of $D$ and
both source infinities. Continuity in the interior forces any finite
interior endpoint of an inverse germ to belong to this fiber. Boundary
limits must be checked separately because $F$ need not be defined there.
Every admissible endpoint and source direction is then tested by
Proposition~\ref{prop:inverse-expression-local-branch}. Exactly one
successful choice gives an unambiguous branch. For an infinite target,
the finite interior fiber is replaced by the observation that a
continuous finite-valued function cannot diverge at an interior point.
\item \textbf{Strict monotonicity on a connected real domain.}
If $D$ is an interval and $F$ is strictly monotone throughout $D$, at
most one real inverse can have values in $D$. A separately verified
local candidate therefore identifies that inverse even without an explicit enumeration of the limiting fiber. This route requires both connectedness
and the global monotonicity proof; local monotonicity near one candidate
alone does not exclude another branch elsewhere.
\end{enumerate}

An arbitrary finite list of candidate endpoints does not establish that
either route has succeeded. For example, the two source sides of $t^2$ give
opposite inverses at $\eta=0^+$ when only realness is required. Restricting
$t>0$ or $t<0$ selects one of them. An explicit source endpoint and
direction specifies a local problem even when the global inverse is ambiguous.
The chosen candidate must still satisfy the domain, limit, target-side,
and monotonicity hypotheses.

\subsection{Transporting the target argument's uncertainty}

After selecting the branch, composition must account for uncertainty in
both the inverse and its argument. The following ordinary power--log
lemma makes the two contributions explicit. Exact exponential carriers
or other scales require their corresponding chart rules from
Section~\ref{sec:series-calculus}.

\begin{lemma}[Two independent composition errors]
\label{lem:inverse-expression-composition-error}
Let $w\to0^+$ and $\M(w)=1+|\log w|$. Suppose the positive local target
coordinate satisfies
\[
 q(w)\sim c w^\lambda,\quad c,\lambda>0,\qquad
 q(w)-\widehat q(w)=\OO(w^\rho\M(w)^d),\quad \rho>\lambda.
\]
Let $G$ be the selected inverse, or a differentiable observable of it,
on this positive target coordinate, and suppose
\[
 G(v)-T(v)=\OO(v^\beta\M(v)^k),\qquad
 |G'(v)|\le C v^{\nu-1}\M(v)^\ell
\]
for sufficiently small $v>0$. Here $\beta,\nu\in\R$ and
$d,k,\ell\ge0$. Then
\begin{equation}
 G(q(w))-T(\widehat q(w))
 =\OO(w^{\lambda\beta}\M(w)^k)
  +\OO(w^{\rho+\lambda(\nu-1)}\M(w)^{d+\ell}).
 \label{eq:inverse-expression-two-errors}
\end{equation}
\end{lemma}

\begin{proof}
The target uncertainty is $o(w^\lambda)$, so $q$, $\widehat q$, and
every point between them are positive and comparable to $w^\lambda$.
Their logarithmic envelopes are comparable to $\M(w)$. Split the
difference as
\[
 [G(q)-G(\widehat q)]+[G(\widehat q)-T(\widehat q)].
\]
Substitution into the stated outer remainder gives the first term in
\eqref{eq:inverse-expression-two-errors}. The mean value theorem and the
derivative bound give the second. The derivative estimate is an
independent hypothesis: no differentiation of a magnitude Big-$O$
statement is used.
\end{proof}

Thus the two candidate power cutoffs are $\lambda\beta$ and
$\rho+\lambda(\nu-1)$. The smaller governs the final power order;
at equal powers the logarithmic envelopes must also be combined.
A vanishing derivative can improve transported input precision, while
a singular derivative can lower it. For an inverse itself, a suitable
derivative estimate can be obtained from the original equation and
$G'=1/(F'\circ G)$ on the selected branch. A larger formal cutoff cannot recover coefficients hidden by the input
argument's uncertainty.
Resonant contributions at a common weight are combined before complete
blocks are retained or discarded.

\subsection{Logarithmic and irrational-power examples}

Consider $F(t)=t+t^2(1+\log t)$ on $t>0$. Its derivative satisfies
\[
 F'(t)=1+t(3+2\log t)\ge1-2e^{-5/2}>0,
\]
because $t(3+2\log t)$ has its global minimum at $t=e^{-5/2}$.
The strictly increasing positive real branch tends to $0$ as $y\to0^+$.
Its inverse begins
\[
 G(y)=y-y^2(1+\log y)
 +y^3(2\log^2y+5\log y+3)
 +\OO(y^4\M(y)^3).
\]
This is reversion of a real source germ, although the inverse has no
ordinary Taylor series at $y=0$.

For $F(t)=t+t^{\sqrt2}$ on $t>0$,
$F'(t)=1+\sqrt2\,t^{\sqrt2-1}>0$. At $0^+$ its inverse has the expansion
\[
 G(y)=y-y^{\sqrt2}+\sqrt2\,y^{2\sqrt2-1}
 -\frac{6-\sqrt2}{2}\,y^{3\sqrt2-2}
 +\OO(y^{4\sqrt2-3}).
\]
At $+\infty$, a different normalization is required because $t^{\sqrt2}$
is dominant. Put $p=\sqrt2$, $z=y^{1/p}$, and $\varepsilon=z^{1-p}$.
Writing $G(y)=zW(\varepsilon)$ gives
\[
 W^p+\varepsilon W=1,\qquad W(0)=1.
\]
The implicit derivative in $W$ at the origin is $p\ne0$, so $W$ has a
convergent ordinary power series near $\varepsilon=0$. Its first terms are
\begin{align*}
 G(y)=y^{1/p}\Bigl(&
 1-\frac{\varepsilon}{p}
 +\frac{3-p}{2p^2}\varepsilon^2\\
 &{}
 -\frac{(p-2)(p-4)}{3p^3}\varepsilon^3
 +\frac{(5-p)(5-2p)(5-3p)}{24p^4}\varepsilon^4
 +\OO(\varepsilon^5)\Bigr),
\end{align*}
where $\varepsilon=y^{(1-p)/p}$.
The same globally increasing real function thus has different natural
asymptotic coordinates at its two endpoints.

An arithmetic expression such as $1+G(x)$ or $G(x)^r$ is an observable
of the inverse. Its defining identity is obtained by reconstructing the
source value $G(x)$ before substituting into $F$. It is not itself an
inverse of $F$. On compatible branches, inversion of $G$ recovers $F$;
composition with a general observable has no such interpretation.

\subsection{Exact identities and the role of the selected root}

A closed form does not remove the need to specify its real branch.
For example, $W_0(y)$ inverts $t e^t$ on $t\ge-1$, while $W_{-1}(y)$
inverts it on $t\le-1$. Their common threshold and their distinct
large-argument or near-zero behavior require the corresponding source
endpoint data. An exact algebraic radical has the same obligation to use
the branch matching the source domain.

A residual evaluated at a reconstructed source root must verify both the
forward equation and the source restriction. Checking a powered observable
alone can lose the source sign. A quantitative interval argument must
establish the defining conditions on its verification interval; an open
source condition excludes an endpoint at equality. A numerical root
satisfying the condition supplies evidence at that point, while a uniform
inverse statement requires the neighborhood hypotheses proved above.

\section{Exponential normalization for Gamma and Barnes functions}
\label{sec:exponential-normalization}

Power--logarithmic scales describe algebraic growth and its corrections.
A function such as $\Gamma(x)^2$ grows faster than every power of $x$.
Its logarithm, however, has a power--logarithmic expansion. The appropriate
construction retains the nonvanishing part of that logarithm exactly and
expands only the vanishing part. This yields a relative asymptotic expansion
whose absolute error includes the extracted growth factor.

\subsection{From absolute logarithmic error to relative error}

Let $u\to0^+$ and write $\M(u)=1+|\log u|$. Suppose a real function has fixed
sign $\sigma\in\{1,-1\}$ sufficiently near the endpoint and
\begin{equation}
 \log|f(u)|=A(u)+B(u)+R(u),\qquad
 R(u)=\OO\bigl(u^\rho \M(u)^k\bigr),\quad \rho>0,
 \label{eq:logarithmic-normalization}
\end{equation}
where $B$ is a finite sum of positive-power logarithmic blocks, so that
$B(u)\to0$. The finite expression $A$ includes all retained nonvanishing
blocks, including any constant. Put $C(u)=\sigma e^{A(u)}$.

\begin{proposition}[Relative normalization]
Under these hypotheses,
\begin{equation}
 f(u)=C(u)\left(e^{B(u)}+
           \OO\bigl(u^\rho \M(u)^k\bigr)\right).
 \label{eq:relative-normalization}
\end{equation}
If $e^B=T+\OO(u^\eta \M(u)^\ell)$, the finite expression $C T$ has
absolute error
\begin{equation}
 \OO\left(|C(u)|
    \left(u^\rho \M(u)^k+u^\eta \M(u)^\ell\right)\right).
 \label{eq:normalized-absolute-error}
\end{equation}
\end{proposition}
\begin{proof}
The identity $f=C e^B e^R$ and the mean-value estimate
$|e^R-1|\le |R|e^{|R|}$ apply for real $R$.
Both $B$ and $R$ tend to zero, so $e^B$ and $e^{|R|}$ are bounded.
The first assertion follows, and the triangle inequality gives the second.
\end{proof}

The logarithmic remainder must be known to tend to zero; the hypothesis
$\rho>0$ ensures this. An error bound of order one in the logarithm
does not imply a relative error tending to zero in the original function. Likewise, a relative expansion of the logarithm
alone need not determine a relative expansion of its exponential.

\subsection{Products, varying powers, and coefficient extraction}

For positive real functions $f_j$ and real functions $r_j$, multiplication
and real exponentiation become addition before any truncation:
\begin{equation}
 \log\left(\prod_{j=1}^m f_j(u)^{r_j(u)}\right)
   =\sum_{j=1}^m r_j(u)\log f_j(u).
 \label{eq:product-logarithms}
\end{equation}
Cancellation can remove leading growth terms. It is therefore useful to
form the complete sum at the required precision before extracting $A$.
If $r_j(u)=\OO(u^{-s_j}\M(u)^{d_j})$, a logarithmic remainder
$\OO(u^h \M(u)^k)$ becomes
$\OO(u^{h-s_j}\M(u)^{k+d_j})$. To reach a desired vanishing absolute
error in the logarithm, the expansion of $\log f_j$ must be correspondingly
deeper. A real varying exponent cannot in general be treated as a constant
when propagating precision.

For an ordinary correction series
$B(t)=\sum_{n\ge1}a_n t^n$, write
$e^{B(t)}=\sum_{n\ge0}c_n t^n$. Formal differentiation gives
\begin{equation}
 c_0=1,\qquad
 c_n=\frac1n\sum_{j=1}^n j a_j c_{n-j}\quad(n\ge1).
 \label{eq:exponential-coefficient-recurrence}
\end{equation}
For real positive weights, the finite convolution algebra of
Section~\ref{sec:scale} supplies the same construction. If the smallest
positive weight is $\delta$, powers $B^n$ with $n\delta\ge h$ cannot
contribute below the strict weight cutoff $h$. Polynomial logarithmic
coefficients stay attached to their complete weight blocks.

\subsection{Stirling's logarithmic expansion}

On the positive real axis, Stirling's expansion is
\begin{equation}
 \log\Gamma(x)=
 \left(x-\tfrac12\right)\log x-x+\tfrac12\log(2\pi)
 +\sum_{n=1}^{N}
   \frac{B_{2n}}{2n(2n-1)x^{2n-1}}
 +\OO(x^{-2N-1}),\qquad x\to+\infty,
 \label{eq:stirling-logarithm}
\end{equation}
where $B_{2n}$ are Bernoulli numbers; see \cite[\S5.11]{dlmf}.
This is a Poincar\'e expansion: the assertion holds for every fixed
truncation order. It is not a convergence assertion as $N\to\infty$.
Multiplying, adding, and substituting into finitely many such expansions
is justified at each fixed order with the corresponding error estimates.

For products $\prod_j\Gamma(a_j(x))^{r_j(x)}$, assume that every argument
is eventually positive, tends to $+\infty$, and admits the substitutions
and error bounds needed in \eqref{eq:stirling-logarithm}. Formula
\eqref{eq:product-logarithms} then determines the logarithmic model.
Factors with arguments tending to positive finite points instead use
the local analytic expansion of $\log\Gamma$. Exponential normalization
combines the resulting growth and correction terms in either case.

\begin{example}[A square and a ratio]
Squaring the Gamma function doubles its logarithm. Exponentiating the
vanishing part gives
\begin{equation}
 \Gamma(x)^2=
 2\pi e^{-2x}x^{2x-1}
 \left(1+\frac1{6x}+\frac1{72x^2}
       -\frac{31}{6480x^3}-\frac{139}{155520x^4}
       +\OO(x^{-5})\right).
 \label{eq:gamma-square-expansion}
\end{equation}
For a ratio, subtract the logarithms first:
\begin{equation}
 \log\frac{\Gamma(3x)}{\Gamma(x)}
 = (3x-\tfrac12)\log3+2x\log x-2x
       -\frac1{18x}+\frac{13}{4860x^3}+\OO(x^{-5}).
\end{equation}
\begin{equation}
\begin{aligned}
 \frac{\Gamma(3x)}{\Gamma(x)}
 &=3^{3x-1/2}x^{2x}e^{-2x}
 \Bigl(1-\frac1{18x}+\frac1{648x^2}\\
 &\hspace{7em}+\frac{463}{174960x^3}-\frac{1867}{12597120x^4}
       +\OO(x^{-5})\Bigr).
\end{aligned}
 \label{eq:gamma-ratio-expansion}
\end{equation}
The remainder inside each parenthesis is relative to the prefactor.
The corresponding absolute error is that prefactor times $\OO(x^{-5})$.
\end{example}

\begin{example}[A varying exponent]
Multiplying \eqref{eq:stirling-logarithm} by $x$ turns the first
Stirling correction into the constant $1/12$. It therefore belongs to
the exact prefactor:
\begin{equation}
 \Gamma(x)^x=e^{1/12}
 \left(\sqrt{2\pi}\,e^{-x}x^{x-1/2}\right)^x
 \left(1-\frac1{360x^2}+\frac{1447}{1814400x^4}
       +\OO(x^{-6})\right).
 \label{eq:gamma-varying-power}
\end{equation}
The step from $\OO(x^{-7})$ in $\log\Gamma(x)$ to $\OO(x^{-6})$
in $x\log\Gamma(x)$ illustrates the required increase in working order.
\end{example}

\subsection{Barnes' \texorpdfstring{$G$}{G}-function and the argument shift}

Barnes' $G$-function satisfies
\begin{equation}
 G(z+1)=\Gamma(z)G(z),\qquad G(1)=1.
 \label{eq:barnes-recurrence}
\end{equation}
It is positive on the positive real axis. Write
$c_G=\zeta'(-1)=1/12-\log A$, where $A$ is Glaisher's constant.
For every fixed integer $N\ge0$, its logarithmic expansion is
\begin{equation}
\begin{aligned}
 \log G(z+1)
  ={}&\left(\frac{z^2}{2}-\frac1{12}\right)\log z
      -\frac{3z^2}{4}+\frac z2\log(2\pi)+c_G\\
    &+\sum_{k=1}^{N}
       \frac{B_{2k+2}}{2k(2k+2)z^{2k}}
       +\OO(z^{-2N-2}),\qquad z\to+\infty.
\end{aligned}
 \label{eq:barnes-logarithm}
\end{equation}
This form follows by inserting Stirling's expansion for
$\log\Gamma(z+1)$ into \cite[(5.17.5)]{dlmf}; the relation between
$A$ and $\zeta'(-1)$ is \cite[(5.17.7)]{dlmf}. The factor $z$ multiplying
that logarithm requires one additional order of its absolute error.
Expanding it through the Bernoulli term with index $2N+2$ gives the
remainder displayed in \eqref{eq:barnes-logarithm}. As with Stirling's
formula, each fixed truncation has its stated error, without asserting
convergence of the infinite series.

The first vanishing logarithmic terms are
\begin{equation}
 -\frac1{240z^2}+\frac1{1008z^4}
 -\frac1{1440z^6}+\frac1{1056z^8}+\OO(z^{-10}).
 \label{eq:barnes-logarithmic-corrections}
\end{equation}
Define the positive prefactor
\begin{equation}
 P_G^+(z)=e^{c_G}(2\pi)^{z/2}
          e^{-3z^2/4}z^{z^2/2-1/12}.
 \label{eq:barnes-shifted-prefactor}
\end{equation}
Relative normalization and
\eqref{eq:exponential-coefficient-recurrence} give, for example,
\begin{equation}
 G(z+1)=P_G^+(z)
 \left(1-\frac1{240z^2}+\frac{269}{268800z^4}
       +\OO(z^{-6})\right).
 \label{eq:barnes-shifted-expansion}
\end{equation}
Every correction power in this coordinate is even. A logarithmic
remainder $\OO(z^{-2N-2})$ becomes a relative remainder of that
same order after exponentiation; the absolute remainder includes
the positive factor $P_G^+(z)$.

For $G(x)$ the argument in \eqref{eq:barnes-logarithm} is
$z=x-1$. Retaining powers of $1/(x-1)$ gives the shifted form above.
Expressing the result in powers of $1/x$ changes both the exact
nonvanishing logarithmic part and the correction coefficients. The
recurrence \eqref{eq:barnes-recurrence} provides an equivalent direct
derivation: subtract \eqref{eq:stirling-logarithm} from the expansion
of $\log G(x+1)$. Put
\begin{equation}
 P_G(x)=e^{c_G}(2\pi)^{(x-1)/2}
        e^{-3x^2/4+x}x^{x^2/2-x+5/12}.
 \label{eq:barnes-prefactor}
\end{equation}
Then
\begin{equation}
\begin{aligned}
 \log G(x)={}&\log P_G(x)-\frac1{12x}-\frac1{240x^2}\\
            &+\frac1{360x^3}+\frac1{1008x^4}+\OO(x^{-5}),
\end{aligned}
 \label{eq:barnes-unshifted-logarithm}
\end{equation}
and its first five relative blocks are
\begin{equation}
\begin{aligned}
 G(x)=P_G(x)\Bigl(&1-\frac1{12x}-\frac1{1440x^2}
              +\frac{157}{51840x^3}\\
              &+\frac{65911}{87091200x^4}+\OO(x^{-5})\Bigr).
\end{aligned}
 \label{eq:barnes-five-blocks}
\end{equation}
The term $z=x-1$ is therefore essential when translating the
logarithmic formula to an expansion of $G(x)$. A shift of its argument
does not preserve the correction bracket in a fixed coordinate.

Products of positive Gamma and Barnes factors use
\eqref{eq:product-logarithms}, so their logarithmic expansions can
be combined before choosing the prefactor. Real varying powers
transport the logarithmic error with the precision loss described
above. Positive finite limiting arguments instead use local analytic
logarithms. Integer shifts can also produce exact simplifications;
for example,
\[
 \frac{G(x+1)}{G(x)}=\Gamma(x),\qquad
 \frac{G(x+1)}{\Gamma(x)G(x)}=1\qquad(x>0).
\]
These identities justify exact cancellation independently of a finite
asymptotic coefficient calculation.

\subsection{Exact identities and elementary growth}

The identities $\Gamma(x+1)=x\Gamma(x)$ and
$\mathrm B(a,b)=\Gamma(a)\Gamma(b)/\Gamma(a+b)$ connect the same
analysis to factorials, binomial coefficients, beta functions, and
rising factorials. For example,
\begin{equation}
 \binom{2x}{x}=\frac{4^x}{\sqrt{\pi x}}
 \left(1-\frac1{8x}+\frac1{128x^2}
       +\frac5{1024x^3}+\OO(x^{-4})\right)
 \qquad(x\to+\infty),
\end{equation}
where the binomial coefficient is defined by its Gamma quotient.
The recurrence also proves $\Gamma(x+1)/\Gamma(x)=x$ exactly.
Vanishing coefficients in a finite asymptotic computation alone would
not prove this identity: a function may have all algebraic coefficients
zero and still have a nonzero flat remainder.

The construction applies equally to elementary growth. Two examples are
\begin{align}
 e^{x+1/x}
   &=e^x\left(1+\frac1x+\frac1{2x^2}+\OO(x^{-3})\right),\\
 x^{x+1/x}
   &=x^x\left(1+\frac{\log x}{x}
             +\frac{(\log x)^2}{2x^2}
             +\OO\left(\frac{(\log x)^3}{x^3}\right)\right).
\end{align}
In each case it is the vanishing logarithmic correction, rather than
the entire growing function, that is expanded in a small scale.


\begin{thebibliography}{30}
\bibitem{q1} V.~Reshetnikov, \emph{How to get an asymptotic of the real-valued branch of the inverse function?}, Mathematica Stack Exchange, question 236367, December 2020. \url{https://mathematica.stackexchange.com/questions/236367/}
\bibitem{q2} V.~Reshetnikov, \emph{Asymptotic expansion for a function containing irrational exponents}, Mathematica Stack Exchange, question 274761. \url{https://mathematica.stackexchange.com/questions/274761/}
\bibitem{dlmf} NIST Digital Library of Mathematical Functions, \S1.10(vii) Lagrange inversion, \S4.13 Lambert $W$-function, \S5.11 Gamma-function asymptotics. \url{https://dlmf.nist.gov/}
\bibitem{gessel} I.~M.~Gessel, \emph{Lagrange inversion}, J.\ Combin.\ Theory Ser.\ A \textbf{144} (2016), 212--249. arXiv:1609.05988.
\bibitem{surya} E.~Surya, L.~Warnke, \emph{Lagrange inversion formula by induction}, arXiv:2305.17576.
\bibitem{edgar1} G.~A.~Edgar, \emph{Transseries for beginners}, Real Anal.\ Exchange \textbf{35} (2009/10), 253--309. arXiv:0801.4877.
\bibitem{edgar2} G.~A.~Edgar, \emph{Transseries: composition, recursion, and convergence}, arXiv:0909.1259.
\bibitem{vdh} J.~van der Hoeven, \emph{Transseries and Real Differential Algebra}, Lecture Notes in Mathematics 1888, Springer, 2006.
\bibitem{salvy} B.~Salvy, J.~Shackell, \emph{Symbolic asymptotics: multiseries of inverse functions}, J.\ Symbolic Comput.\ \textbf{27} (1999), 543--563.
\bibitem{cghjk} R.~M.~Corless, G.~H.~Gonnet, D.~E.~G.~Hare, D.~J.~Jeffrey, D.~E.~Knuth, \emph{On the Lambert $W$ function}, Adv.\ Comput.\ Math.\ \textbf{5} (1996), 329--359.
\bibitem{proveit} V.~Reshetnikov, \emph{Transseries: the polynomial--logarithmic calculus, series reversal at infinity, and the inversion of rapidly growing functions} and \emph{Combinatorial transseries and their inverses}, ProveIt repository, directory \path{Analysis/FabiusFunction/docs/semi-formalized-research-frontiers/drafts/series-and-transseries}, September 2026. \url{https://github.com/VladimirReshetnikov/ProveIt}
\end{thebibliography}

\end{document}
